\documentclass{article}
\usepackage{hyperref}
\usepackage{geometry}
\geometry{margin=1in}
\title{Manual for {\tt alg-graph-thry2.sage}\\
  {\small{a graph theory and harmonic morphisms Python module}}}
\author{David Joyner}
\date{2025-07-31}
\begin{document}
\maketitle


\vskip .5in

This is a manual for a collection of 150+ related Python functions
to deal with oriented graphs, edge-weighted graphs,
graphs with partitions, and harmonic morphisms, as in articles
such as Norman Biggs's and Matt Baker's works
(see for example [BN], [B]), and in my books with Caroline Melles,
[JM], [MJ].


\begin{itemize}
\item[][BN] M. Baker and S. Norine,
  {\it Riemann-Roch and Abel-Jacobi theory on a finite graph,} available:
\newline  
\url{https://arxiv.org/abs/math/0608360}

\item[][B] N. Biggs, {\it Chip firing and the critical group of a graph,}
  J. Algeb. Comb. 9 (1999)25-45.

\item[][JM]
D. Joyner and C. Melles, {\bf Adventures in graph theory,}
Birkhauser.

\item[][MJ]
  C. Melles and D. Joyner, {\bf Excursions in graph theory,}
  to appear.
  
\end{itemize}

Happy programming!

\vskip 1in

\tableofcontents
\newpage

The module is especially useful in the context of spectral graph theory,
algebraic graph theory, and tropical geometry. It provides computational
tools that complement theoretical frameworks in these areas.

The included functions support a variety of tasks such as:

\begin{itemize}
\item
Graph transformations and conversions, e.g., converting graphs to
standard form.

\item
Matrix and linear algebra utilities relevant to graphs, such
as computing orthogonal eigenbases, permutation matrices,
and weighted incidence matrices.

\item
Walks and cycles, including plotting walks on graphs,
analyzing edge sets in walks, and computing the image of walks
under graph morphisms.

\item
Graph morphism tools, including checks for harmonicity
(both weighted and unweighted), and computing local/global
horizontal and vertical multiplicities.

\item
Graph Laplacians and spectral tools, including construction
of codifferential matrices and eigenbases aligned with graph topology.

\item
Zeta function-related computations, such as constructing
Terras's oriented edge matrix or computing the Ihara zeta polynomial
of a graph.

\end{itemize}

Each function in the module is presented with its name, syntax,
a brief description of its purpose, and an example of usage
(usually). The manual is structured to facilitate quick
reference for users working with {\tt SageMath}
(www.sagemath.org) or custom
Python codebases in discrete mathematics and graph theory.

\section{\tt \ \ compactify\_list}

\textbf{Syntax:} \verb}compactify_list(L)}\\

\textbf{Description:}

L is a list of integers of length n, possibly with repitition,
    taken from integers in an ordered set 
    S of distinct integers of length k.
    
    This function returns a corresponding list (also of length n)
    of integers taken from the range {0, 1, ..., k-1}, obtained by
    "deleting" gaps in and 0-based reindexing.

\textbf{Example:}
\begin{verbatim}
EXAMPLES:
    sage: compactify_list([10, 30, 10, 20])  # Mapping: 10 -> 0, 20 -> 1, 30 -> 2
    [0, 2, 0, 1]
    sage: compactify_list( [5, 7, 5, 9] ) 
    [0, 1, 0, 2]
\end{verbatim}

\section{\tt \ \ graph\_convert\_to\_standard}

\textbf{Syntax:} \verb}graph_convert_to_standard(Gamma)}\\

\textbf{Description:}

This utility function simply converts the (unlabeled,
    ie, no edge weights) graph Gamma = (V, E), assumed 
    to be in some other indexing (eg, as a grid graph or cube graph) 
    into the standard numerical ordering (so the vertices are
[0, ..., n-1], for example.

    The ordering is preserved (I hope).

\textbf{Example:}
\begin{verbatim}
    EXAMPLES:
        sage: Gamma = graphs.CubeGraph(3)    
        sage: Gamma.vertices(sort=True)
         ['000', '001', '010', '011', '100', '101', '110', '111']
        sage: Gamma.edges(labels=None,sort=True)
         [('000', '001'), ('000', '010'), ('000', '100'), ('001', '011'),
          ('001', '101'), ('010', '011'), ('010', '110'), ('011', '111'),
          ('100', '101'), ('100', '110'), ('101', '111'), ('110', '111')]
        sage: Gamma0 = graph_convert_to_standard(Gamma)
        sage: Gamma0.vertices(sort=True)
         [0, 1, 2, 3, 4, 5, 6, 7]
        sage: Gamma0.edges(labels=None,sort=True)
         [(0, 1), (0, 2), (0, 4), (1, 3), (1, 5), (2, 3), (2, 6),
          (3, 7), (4, 5), (4, 6), (5, 7), (6, 7)]
        sage: Gamma.is_isomorphic(Gamma0)
         True
\end{verbatim}

\section{\tt \ \ find\_matrix\_permutation}

\textbf{Syntax:} \verb}find_matrix_permutation(L1, L2)}\\

\textbf{Description:}

\noindent
INPUT: Two lists (possibly with repeated elements) of length n, L1, L2

\noindent
OUTPUT: A permutation matrix sending L1 to L2

\textbf{Example:}
\begin{verbatim}
EXAMPLE:
        sage: find_matrix_permutation([2,4,4,0,0], [0,0,2,4,4])
        [0 0 0 1 0]
        [0 0 0 0 1]
        [1 0 0 0 0]
        [0 1 0 0 0]
        [0 0 1 0 0]

\end{verbatim}
%REFERENCE:
%     Uses material from
%     \url{https://ask.sagemath.org/question/60779/how-to-find-a-permutation-which-sends-a-given-list-to-another-given-list-in-sagemath/}
     
\section{\tt \ \ orthogonal\_eigenbasis\_sorted}

\textbf{Syntax:} \verb}orthogonal_eigenbasis_sorted(M, F=QQ)}\\

\textbf{Description:}
\begin{verbatim}
Input:
      M -- a nxn symmetric real matrix (with real eigenvalues)

Outputs:
      D -- a matrix with eigenvalues D_i sorted in monotonically
           increasing order
      P -- a matrix whose n columns P_i are orthogonal eigenvectors of M s.t.
	   M*P_i = D_i*P_i
      N -- a diagonal matrix of norms of the P_i, so
           that P*N^(-1) is orthonormal.

    This works because the diagonalization method uses a matrix of eigenvectors for the transformation matrix.
\end{verbatim}

\textbf{Example:}
\begin{verbatim}
EXAMPLES:
        sage: Gamma1 = graphs.PathGraph(3)
        sage: Q1 = Gamma1.laplacian_matrix(); Q1
        [ 1 -1  0]
        [-1  2 -1]
        [ 0 -1  1]
        sage:  orthogonal_eigenbasis_sorted(Q1)
        (
        [0 0 0]  [ 1  1  1]  [sqrt(3)       0       0]
        [0 1 0]  [ 1  0 -2]  [      0 sqrt(2)       0]
        [0 0 3], [ 1 -1  1], [      0       0 sqrt(6)]
        )
        sage: D,P,N = orthogonal_eigenbasis_sorted(Q1)
        sage: i=0; Q1*P.columns()[i] == D.diagonal()[i]*P.columns()[i]
        True
        sage: i=1; Q1*P.columns()[i] == D.diagonal()[i]*P.columns()[i]
        True
        sage: i=2; Q1*P.columns()[i] == D.diagonal()[i]*P.columns()[i]
        True
        sage: U = P*N^(-1)
        sage: U.transpose()*U
        [1 0 0]
        [0 1 0]
        [0 0 1]
\end{verbatim}

\section{\tt \ \ plot\_walk\_in\_graph}

\textbf{Syntax:} \verb}plot_walk_in_graph(Gamma, W)}\\

\textbf{Description:}

This function plots a walk W in a graph Gamma, returned as a graphics
    object (append \verb+.show()+ to view). The graph is plotted as usual 
    and the edges of the walk (more precisely the edges of the path 
    determined by it) are plotted in thick red.
    
    Returns False if W is not a walk in Gamma.

\textbf{Example:}
\begin{verbatim}
EXAMPLES:
        sage: Gamma = graphs.ButterflyGraph()
        sage: W = [(0,4),(4,1)]
        sage: P = plot_walk_in_graph(Gamma, W)

        Type P.show() to view.
\end{verbatim}

\section{\tt \ \ edges\_in\_cycle\_basis}

\textbf{Syntax:} \verb+edges_in_cycle_basis(Gamma)+\\

\textbf{Description:}
A cycle of a graph Gamma = (V, E) is a connected sequence of edges that returns
to where you started without revisiting any edge.

Computes {\tt L = Gamma.cycle\_basis()}, each element of which is a 
list of vertices of a ``basis'' cycle in Gamma, then for each r in L,
returns the list of edges determined by r.

\textbf{Example:}
\begin{verbatim}
EXAMPLES:
        sage: V2 = [0, 1, 2, 3, 4, 5, 6, 7, 8]
        sage: E2 = [(0,1),(0,2),(0,7),(0,8),(1,2),(1,3),(1,7),(2,3),(2,8),
        ....: (3, 4), (3, 5), (4, 5), (4, 6), (4, 8), (5, 6), (5, 7), (6, 7), (6, 8)]
        sage: Gamma2 = Graph([V2,E2])
        sage: Gamma2.cycle_basis()
        [[2, 8, 0],
         [4, 6, 8],
         [7, 6, 8, 0],
         [4, 5, 6],
         [7, 5, 6],
         [1, 3, 5, 6, 8, 0],
         [2, 3, 5, 6, 8],
         [4, 3, 5],
         [1, 7, 0],
         [1, 2, 0]]
        sage: edges_in_cycle_basis(Gamma2)
        [[(2, 8), (8, 0), (0, 2)],
         [(4, 6), (6, 8), (8, 4)],
         [(7, 6), (6, 8), (8, 0), (0, 7)],
         [(4, 5), (5, 6), (6, 4)],
         [(7, 5), (5, 6), (6, 7)],
         [(1, 3), (3, 5), (5, 6), (6, 8), (8, 0), (0, 1)],
         [(2, 3), (3, 5), (5, 6), (6, 8), (8, 2)],
         [(4, 3), (3, 5), (5, 4)],
         [(1, 7), (7, 0), (0, 1)],
         [(1, 2), (2, 0), (0, 1)]]
\end{verbatim}

\section{\tt \ \ tree\_component\_plus}

\textbf{Syntax:} \verb}tree_component_plus(Gamma, e, T, eo)}\\

\textbf{Description:}
\begin{verbatim}
This computes the positive component T_e^+ of a spanning tree T in a graph Gamma, 
with e in T.

    INPUT:
        Gamma - graph
        e  - edge of Gamma
        T  - spanning tree of Gamma  
        eo - a vector of 1's and -1's whose length is the number of edges in Gamma
             (ie, the size of Gamma, M)
\end{verbatim}
        
\textbf{Example:}
\begin{verbatim}
    EXAMPLES:
        sage: Gamma = graphs.PaleyGraph(9)
        sage: T = Gamma.min_spanning_tree()
        sage: E = Gamma.edges(sort=True)
        sage: e = E[5]
        sage: e in T
        True
        sage: eo = [1, -1, 1, 1, -1, -1, 1, -1, 1, 1, -1, -1, 1, -1, 1, 1, -1, -1]
        sage: tree_component_plus(Gamma, e, T, eo)
        Subgraph of (): Graph on 1 vertex
        sage: tree_component_plus(Gamma, e, T, eo).vertices(sort=True)
        [a + 2]
        sage: tree_component_minus(Gamma, e, T, eo)
        Subgraph of (): Graph on 8 vertices
\end{verbatim}

\section{\tt \ \ tree\_component\_minus}

\textbf{Syntax:} \verb}tree_component_minus(Gamma, e, T, eo)}\\

\textbf{Description:}
\begin{verbatim}
This computes the negative component T_e^- of a spanning tree T in a graph Gamma, with e in T.

    INPUT:
        Gamma - graph
        e  - edge of Gamma
        T  - spanning tree of Gamma  
        eo - a vector of 1's and -1's whose length is the number of edges in Gamma
             (ie, the size of Gamma, M)
\end{verbatim}
        
\textbf{Example:}
\begin{verbatim}
    EXAMPLES:
        sage: Gamma = graphs.PaleyGraph(9)
        sage: T = Gamma.min_spanning_tree()
        sage: E = Gamma.edges(sort=True)
        sage: e = E[5]
        sage: e in T
        True
        sage: eo = [1, -1, 1, 1, -1, -1, 1, -1, 1, 1, -1, -1, 1, -1, 1, 1, -1, -1]
        sage: tree_component_plus(Gamma, e, T, eo)
        Subgraph of (): Graph on 1 vertex
        sage: tree_component_plus(Gamma, e, T, eo).vertices(sort=True)
        [a + 2]
        sage: tree_component_minus(Gamma, e, T, eo).vertices(sort=True)
        [0, 1, 2, a, a + 1, 2*a, 2*a + 1, 2*a + 2]
\end{verbatim}

\section{\tt \ \ is\_incident\_to}

\textbf{Syntax:} \verb}is_incident_to(Gamma, e, v)}\\

\textbf{Description:}
\begin{verbatim}
    Returns True if 
       (a) v is a vertex of a graph Gamma, and
       (b) e is an edge of the graph Gamma,
       (c) e is incident to v
    Returns False otherwise.
\end{verbatim}

\textbf{Example:}
\begin{verbatim}
    EXAMPLES:
         sage: V1 = [0,1,2,3]
         sage: E1 = [(0, 1), (1, 2), (1, 3)]
         sage: Gamma1 = Graph([V1,E1])
         sage: is_incident_to(Gamma1, E1[0], V1[2])
         False
         sage: is_incident_to(Gamma1, E1[0], V1[1])
         True
\end{verbatim}

\section{\tt \ \ incidence\_value}

\textbf{Syntax:} \verb}incidence_value(Gamma, v, e, eo)}\\

\textbf{Description:}
\begin{verbatim}
This computes the incidence value of a vertex and edge of a graph Gamma
with edge-orientation vector eo. 

    INPUT:
        Gamma - graph
        v  - vertex of Gamma
        e  - edge of Gamma  
        eo - a pair [E, e0] where e0 is a vector of 1's and -1's of 
             length |E|
\end{verbatim}

\textbf{Example:}
\begin{verbatim}
    EXAMPLES:
        sage: Gamma = graphs.PaleyGraph(9)
        sage: E = list(Gamma.edges(labels=None,sort=True))
        sage: V = Gamma.vertices(sort=True)
        sage: eo = [E, [1]*len(E)]
        sage: incidence_value(Gamma, V[2], E[3], eo)
        0
        sage: incidence_value(Gamma, V[8], E[3], eo)
        -1
\end{verbatim}

\section{\tt \ \ incidence\_matrix}

\textbf{Syntax:} \verb}incidence_matrix(Gamma, eo)}\\

\textbf{Description:}
\begin{verbatim}
This computes the incidence matrix (whose rows are indexed by edges and whose
columns are indexed by vertices) of a graph Gamma with edge-orientation vector eo.
The ordering of the edges and of the vertices is the same as Sage's vertices and edges methods.

    INPUT:
        Gamma - graph
        eo - a pair [E, e0] where e0 is a vector of 1's and -1's whose 
             length is |E|     (ie, the size of Gamma, M)
\end{verbatim}

\textbf{Example:}
\begin{verbatim}
    EXAMPLES:
        sage: Gamma = graphs.PaleyGraph(9)
        sage: E = list(Gamma.edges(labels=None,sort=True))
        sage: V = Gamma.vertices(sort=True)
        sage: eo = [E, [1, -1, 1, 1, -1, -1, 1, -1, 1, 1, -1, -1, 1, -1, 1, 1, -1, -1]]
        sage: incidence_matrix(Gamma, eo)
        [ 1 -1  1  1  0  0  0  0  0  0  0  0  0  0  0  0  0  0]
        [-1  0  0  0 -1 -1  1  0  0  0  0  0  0  0  0  0  0  0]
        [ 0  1  0  0  1  0  0 -1  1  0  0  0  0  0  0  0  0  0]
        [ 0  0  0  0  0  0  0  1  0  1 -1 -1  0  0  0  0  0  0]
        [ 0  0 -1  0  0  0  0  0  0 -1  0  0  1 -1  0  0  0  0]
        [ 0  0  0  0  0  1  0  0  0  0  1  0 -1  0  1  0  0  0]
        [ 0  0  0  0  0  0 -1  0  0  0  0  0  0  0 -1  1 -1  0]
        [ 0  0  0  0  0  0  0  0 -1  0  0  1  0  0  0 -1  0 -1]
        [ 0  0  0 -1  0  0  0  0  0  0  0  0  0  1  0  0  1  1]

        sage: IM = incidence_matrix(Gamma, eo)
        sage: IM1 = IM.delete_rows([0])
        sage: Proj1 = IM1.transpose()*(IM1*IM1.transpose())^(-1)*IM1
        sage: Proj1.image() # row space
        Vector space of degree 18 and dimension 8 over Rational Field
        Basis matrix:
        [ 1  0  0  0  1  0  0  0  0  0 -1  0  1  0  0 -1  1  0]
        [ 0  1  0  0  1  0  0  0  0  1 -1  0  0  0  0 -1  0 -1]
        [ 0  0  1  0  0  0  0  0  0  1  0  0 -1  1  0  0  0  0]
        [ 0  0  0  1  0  0  0  0  0  0  0  0  0 -1  0  0 -1 -1]
        [ 0  0  0  0  0  1  0  0  0  0  1  0 -1  0  1  0  0  0]
        [ 0  0  0  0  0  0  1  0  0  0  0  0  0  0  1 -1  1  0]
        [ 0  0  0  0  0  0  0  1  0  1 -1 -1  0  0  0  0  0  0]
        [ 0  0  0  0  0  0  0  0  1  0  0 -1  0  0  0  1  0  1]
        sage: cut_space(Gamma, eo)
        Vector space of degree 18 and dimension 8 over Rational Field
        Basis matrix:
        [ 1  0  0  0  1  0  0  0  0  0 -1  0  1  0  0 -1  1  0]
        [ 0  1  0  0  1  0  0  0  0  1 -1  0  0  0  0 -1  0 -1]
        [ 0  0  1  0  0  0  0  0  0  1  0  0 -1  1  0  0  0  0]
        [ 0  0  0  1  0  0  0  0  0  0  0  0  0 -1  0  0 -1 -1]
        [ 0  0  0  0  0  1  0  0  0  0  1  0 -1  0  1  0  0  0]
        [ 0  0  0  0  0  0  1  0  0  0  0  0  0  0  1 -1  1  0]
        [ 0  0  0  0  0  0  0  1  0  1 -1 -1  0  0  0  0  0  0]
        [ 0  0  0  0  0  0  0  0  1  0  0 -1  0  0  0  1  0  1]
        sage: B = incidence_matrix(Gamma, eo) 
        sage: B*B.transpose() == Gamma.laplacian_matrix()
        True

        sage: V = range(10)
        sage: E = [(0,1),(0,2),(0,5),(1,2),(1,3),(1,5),(2,4),(2,5),(3,4),(3,5),(3,6),\
        ....: (4,5),(4,7),(5,6),(5,7),(5,8),(6,7),(6,8),(7,8),(8,9)]
        sage: Gamma = Graph([V,E]) # Seifrot graph of Kaballah mysticism
        sage: eo = [E,20*[1]]
        sage: incidence_matrix(Gamma, eo)
        [ 1  1  1  0  0  0  0  0  0  0  0  0  0  0  0  0  0  0  0  0]
        [-1  0  0  1  1  1  0  0  0  0  0  0  0  0  0  0  0  0  0  0]
        [ 0 -1  0 -1  0  0  1  1  0  0  0  0  0  0  0  0  0  0  0  0]
        [ 0  0  0  0 -1  0  0  0  1  1  1  0  0  0  0  0  0  0  0  0]
        [ 0  0  0  0  0  0 -1  0 -1  0  0  1  1  0  0  0  0  0  0  0]
        [ 0  0 -1  0  0 -1  0 -1  0 -1  0 -1  0  1  1  1  0  0  0  0]
        [ 0  0  0  0  0  0  0  0  0  0 -1  0  0 -1  0  0  1  1  0  0]
        [ 0  0  0  0  0  0  0  0  0  0  0  0 -1  0 -1  0 -1  0  1  0]
        [ 0  0  0  0  0  0  0  0  0  0  0  0  0  0  0 -1  0 -1 -1  1]
        [ 0  0  0  0  0  0  0  0  0  0  0  0  0  0  0  0  0  0  0 -1]
\end{verbatim}

\section{\tt \ \ weight\_at\_edge}

\textbf{Syntax:} \verb}weight_at_edge(Gamma, w0, e)}\\

\textbf{Description:}
\begin{verbatim}
Returns the edge weight of e on an edge-weighted graph Gamma
    (the value is independent of edge orientation),
but returns -1 if w0 is not a legal set of edge-weights.

    eo - a pair [E, e0] where e0 is a vector of 1's and -1's of length |E|
    w0 - a pair [E, wt0] where wt0 is a vector of length |E| whose entries are > 0
\end{verbatim}

\textbf{Example:}
\begin{verbatim}
EXAMPLES:
        sage: V = [0,1,2,3]
        sage: E = [(0, 1), (0, 2), (0, 3), (1, 2)]
        sage: Gamma = Graph([V, E])
        sage: w0 = [E, [1/2, 3/5, 4/3, 2/3]]
        sage: weight_at_edge(Gamma, w0, (2, 1))
         2/3
        sage: weight_at_edge(Gamma, w0, (0, 3))
         4/3
        sage: Ew = [(0, 1, 1/2), (0, 2, 3/5), (0, 3, 4/3), (1, 2, 2/3)]
        sage: Gamma_w = Graph([V,Ew], weighted = True)
        sage: Gamma_w.weighted_adjacency_matrix()
         [  0 1/2 3/5 4/3]
         [1/2   0 2/3   0]
         [3/5 2/3   0   0]
         [4/3   0   0   0]
\end{verbatim}

\section{\tt \ \ weighted\_degree}

\textbf{Syntax:} \verb}weighted_degree(Gamma, w0, v)}\\

\textbf{Description:}
Returns the weighted degree of v on an edge-weighted graph Gamma (the value is
independent of edge orientation).

INPUT:
    w0 - a pair [E, wt0] where wt0 is a vector of length |E| whose entries are > 0
    v - a vertex of Gamma

\textbf{Example:}
\begin{verbatim}
EXAMPLES:
        sage: V = [0,1,2,3]
        sage: E = [(0, 1), (0, 2), (0, 3), (1, 2)]
        sage: Gamma = Graph([V, E])
        sage: w0 = [E, [1/2, 3/5, 4/3, 2/3]]
        sage: weighted_degree(Gamma, w0, 0)
         73/30
\end{verbatim}

\section{\tt \ \ weighted\_degree\_matrix}

\textbf{Syntax:} \verb}weighted_degree_matrix(Gamma, w0)}\\

\textbf{Description:}
\begin{verbatim}

Returns the diagonal weighted degree matrix of Gamma.
    (the value is independent of edge orientation)

    eo - a pair [E, e0] where e0 is a vector of 1's and -1's of length |E|
    w0 - a pair [E, wt0] where wt0 is a vector of length |E| whose entries are > 0
\end{verbatim}

\textbf{Example:}
\begin{verbatim}
EXAMPLES:
        sage: V = [0,1,2,3]
        sage: E = [(0, 1), (0, 2), (0, 3), (1, 2)]
        sage: Gamma = Graph([V, E])
        sage: w0 = [E, [1/2, 3/5, 4/3, 2/3]]
        sage: weighted_degree_matrix(Gamma, w0)
         [73/30     0     0     0]
         [    0   7/6     0     0]
         [    0     0 19/15     0]
         [    0     0     0   4/3]
\end{verbatim}


\section{\tt \ \ weighted\_incidence\_matrix}

\textbf{Syntax:} \verb}weighted_incidence_matrix(Gamma, eo, w0)}\\

\textbf{Description:}
\begin{verbatim}
    This computes the incidence matrix (whose rows are indexed by edges
    and whose columns are indexed by vertices) of an edge-weighted
    graph Gamma = (V, E) with 

    eo - a pair [E,e0], where e0 an edge-orientation vector of length |E|
               with entries in {-1, 1}
    w0 - a pair [E,wt0], where wt0 an edge-weight vector of length |E|
               with entries > 0. 
    (These are represented as pairs to ensure the ordering of the edges is the same 
    as that of the orientation list and weight list.)

    INPUT:
        Gamma - graph
        eo - a pair [E, e0] where e0 is a vector of 1's and -1's of length |E|
        w0 - a pair [E, wt0] where wt0 is a vector of length |E| whose entries are > 0
\end{verbatim}

\textbf{Example:}
\begin{verbatim}
    EXAMPLES:
        sage: V = [0,1,2,3]
        sage: E = [(0, 1), (0, 2), (0, 3), (1, 2)]
        sage: Gamma = Graph([V, E])
        sage: eo = [E,[1, 1, 1, 1]]
        sage: w0 = [E, [1/2, 3/5, 4/3, 2/3]]
        sage: B_w = weighted_incidence_matrix(Gamma, eo, w0); B_w
         [   sqrt(1/2)    sqrt(3/5)  2*sqrt(1/3)            0]
         [  -sqrt(1/2)            0            0    sqrt(2/3)]
         [           0   -sqrt(3/5)            0   -sqrt(2/3)]
         [           0            0 -2*sqrt(1/3)            0]
        sage: B_w*B_w.transpose()
         [73/30  -1/2  -3/5  -4/3]
         [ -1/2   7/6  -2/3     0]
         [ -3/5  -2/3 19/15     0]
         [ -4/3     0     0   4/3]
        sage: Ew = [(0, 1, 1/2), (0, 2, 3/5), (0, 3, 4/3), (1, 2, 2/3)]
        sage: Gamma_w = Graph([V,Ew], weighted = True)
        sage: Gamma_w.laplacian_matrix()
         [73/30  -1/2  -3/5  -4/3]
         [ -1/2   7/6  -2/3     0]
         [ -3/5  -2/3 19/15     0]
         [ -4/3     0     0   4/3]
        sage: weighted_degree_matrix(Gamma, e0, w0)
         [73/30     0     0     0]
         [    0   7/6     0     0]
         [    0     0 19/15     0]
         [    0     0     0   4/3]
\end{verbatim}


\section{\tt \ \ oriented\_edge\_matrix}

\textbf{Syntax:} \verb}oriented_edge_matrix(Gamma, eo)}\\

\textbf{Description:}
\begin{verbatim}
Returns the 2mx2m (0,1)-edge matrix as defined by Terras (who denotes it by 
W or W_1), where m is the number of edges of Gamma. Satisfies 

    det(I_2m-x*W) = (Ihara zeta polynomial of Gamma).

    INPUT:
        Gamma - graph, (V, E)
           eo - a pair [E,e0], where e0 an edge-orientation vector of length |E|
                with entries in {-1, 1}

\end{verbatim}

\textbf{Example:}
\begin{verbatim}
    EXAMPLES:
        sage: Gamma = graphs.CycleGraph(4)
        sage: E = list(Gamma.edges(labels=None,sort=True))
        sage: eo = [E, len(E)*[1]]
        sage: I2m = diagonal_matrix(2*m*[1])
        sage: W = oriented_edge_matrix(Gamma, eo); W
        [0 0 1 0 0 0 0 0]
        [0 0 0 0 0 0 0 1]
        [0 0 0 1 0 0 0 0]
        [0 0 0 0 0 1 0 0]
        [0 1 0 0 0 0 0 0]
        [1 0 0 0 0 0 0 0]
        [0 0 0 0 1 0 0 0]
        [0 0 0 0 0 0 1 0]
        sage: det(I2m-x*W)
        x^8 - 2*x^4 + 1
        sage: expand(zeta_polynomial_graph(Gamma))
        x^8 - 2*x^4 + 1
\end{verbatim}

\section{\tt \ \ cycle\_space}

\textbf{Syntax:} \verb}cycle_space(Gamma, eo, F = QQ)}\\

\textbf{Description:}
\begin{verbatim}
This computes the cycle space of Gamma, as a Sage vector space over the field F.

    INPUT:
        Gamma - graph, (V, E)
           eo - a pair [E,e0], where e0 an edge-orientation vector of length |E|
                with entries in {-1, 1}
\end{verbatim}

This vector space is independent of the orientation vector.

\textbf{Example:}
\begin{verbatim}
    EXAMPLES:
        sage: Gamma = graphs.PaleyGraph(9)
        sage: E = Gamma.edges(sort=True)
        sage: V = Gamma.vertices(sort=True)
        sage: eo = [1, -1, 1, 1, -1, -1, 1, -1, 1, 1, -1, -1, 1, -1, 1, 1, -1, -1]
        sage: cycle_space(Gamma, eo)
        Vector space of degree 18 and dimension 10 over Rational Field
        Basis matrix:
        [ 1  0  0 -1  0  0  1  0  0  0  0  0  0  0  0  0 -1  0]
        [ 0  1  0  1  0  0  0  0 -1  0  0  0  0  0  0  0  0  1]
        [ 0  0  1 -1  0  0  0  0  0  0  0  0  0 -1  0  0  0  0]
        [ 0  0  0  0  1  0  1  0 -1  0  0  0  0  0  0  0 -1  1]
        [ 0  0  0  0  0  1  1  0  0  0  0  0  0  0 -1  0  0  0]
        [ 0  0  0  0  0  0  0  1  1  0  0  1  0  0  0  0  0  0]
        [ 0  0  0  0  0  0  0  0  0  1  0  1  0 -1  0  0  0  1]
        [ 0  0  0  0  0  0  0  0  0  0  1 -1  0  0 -1  0  1 -1]
        [ 0  0  0  0  0  0  0  0  0  0  0  0  1  1  1  0 -1  0]
        [ 0  0  0  0  0  0  0  0  0  0  0  0  0  0  0  1  1 -1]
\end{verbatim}

\section{\tt \ \ cocycle\_space}

\textbf{Syntax:} \verb}cocycle_space(Gamma, eo)}\\

\textbf{Description:}
\begin{verbatim}
This computes the cocycle space (= cut space) B of Gamma.

    INPUT:
        Gamma - graph, (V, E)
           eo - a pair [E1,e0], where e0 an edge-orientation vector of length |E1|
                with entries in {-1, 1}

\end{verbatim}

This vector space is independent of the orientation vector.

\textbf{Example:}
\begin{verbatim}
See cut_space(Gamma, eo).
\end{verbatim}

\section{\tt \ \ cut\_space}

\textbf{Syntax:} \verb}cut_space(Gamma, eo)}\\

\textbf{Description:}
\begin{verbatim}
This computes the cut space B of Gamma. 

A cut is a way to partition vertices V of the graph Gamma into two subsets.
If you split up V into A and B, a cut is the set of edges that have one 
end in A and one end in B.

INPUT:
        Gamma - graph, (V, E)
           eo - a pair [E1,e0], where e0 an edge-orientation vector of length |E1|
           with entries in {-1, 1}

\end{verbatim}

\textbf{Example:}
\begin{verbatim}
    EXAMPLES:
        sage: Gamma = graphs.PaleyGraph(9)
        sage: E = list(Gamma.edges(labels=None,sort=True))
        sage: V = Gamma.vertices(sort=True)
        sage: eo = [E, [1, -1, 1, 1, -1, -1, 1, -1, 1, 1, -1, -1, 1, -1, 1, 1, -1, -1]]
        sage: cut_space(Gamma, eo)
        Vector space of degree 18 and dimension 8 over Rational Field
        Basis matrix:
        [ 1  0  0  0  1  0  0  0  0  0 -1  0  1  0  0 -1  1  0]
        [ 0  1  0  0  1  0  0  0  0  1 -1  0  0  0  0 -1  0 -1]
        [ 0  0  1  0  0  0  0  0  0  1  0  0 -1  1  0  0  0  0]
        [ 0  0  0  1  0  0  0  0  0  0  0  0  0 -1  0  0 -1 -1]
        [ 0  0  0  0  0  1  0  0  0  0  1  0 -1  0  1  0  0  0]
        [ 0  0  0  0  0  0  1  0  0  0  0  0  0  0  1 -1  1  0]
        [ 0  0  0  0  0  0  0  1  0  1 -1 -1  0  0  0  0  0  0]
        [ 0  0  0  0  0  0  0  0  1  0  0 -1  0  0  0  1  0  1]
        sage: v = cut_space(Gamma, eo).random_element()
        sage: w = cycle_space(Gamma, eo).random_element()
        sage: v.inner_product(w)
        0
\end{verbatim}

\section{\tt \ \ cuts\_list}

\textbf{Syntax:} \verb}cuts_list(Gamma, verbose=False)}\\

\textbf{Description:}
\begin{verbatim}
Returns list of all cuts of Gamma.
\end{verbatim}

\textbf{Example:}
\begin{verbatim}
    EXAMPLES:
        sage: Gamma =  graphs.DiamondGraph()
        sage: cuts_list(Gamma)

        [[(0, 1), (0, 2)],
         [(1, 0), (1, 2), (1, 3)],
         [(0, 2), (1, 2), (1, 3)],
         [(2, 0), (2, 1), (2, 3)],
         [(0, 1), (2, 1), (2, 3)],
         [(1, 0), (1, 3), (2, 0), (2, 3)],
         [(1, 3), (2, 3)]]
        sage: Gamma = graphs.CompleteGraph(4)
        sage: cuts_list(Gamma)

        [[(0, 1), (0, 2), (0, 3)],
         [(1, 0), (1, 2), (1, 3)],
         [(0, 2), (0, 3), (1, 2), (1, 3)],
         [(2, 0), (2, 1), (2, 3)],
         [(0, 1), (0, 3), (2, 1), (2, 3)],
         [(1, 0), (1, 3), (2, 0), (2, 3)],
         [(0, 3), (1, 3), (2, 3)]]
        sage: Gamma = graphs.HouseGraph()
        sage: cuts_list(Gamma, verbose=True)

        [[[[], [0, 1, 2, 3, 4]]],
         [[[0], [1, 2, 3, 4]], (0, 1), (0, 2)],
         [[[1], [0, 2, 3, 4]], (1, 0), (1, 3)],
         [[[0, 1], [2, 3, 4]], (0, 2), (1, 3)],
         [[[2], [0, 1, 3, 4]], (2, 0), (2, 3), (2, 4)],
         [[[0, 2], [1, 3, 4]], (0, 1), (2, 3), (2, 4)],
         [[[1, 2], [0, 3, 4]], (1, 0), (1, 3), (2, 0), (2, 3), (2, 4)],
         [[[0, 1, 2], [3, 4]], (1, 3), (2, 3), (2, 4)],
         [[[3], [0, 1, 2, 4]], (3, 1), (3, 2), (3, 4)],
         [[[0, 3], [1, 2, 4]], (0, 1), (0, 2), (3, 1), (3, 2), (3, 4)],
         [[[1, 3], [0, 2, 4]], (1, 0), (3, 2), (3, 4)],
         [[[0, 1, 3], [2, 4]], (0, 2), (3, 2), (3, 4)],
         [[[2, 3], [0, 1, 4]], (2, 0), (2, 4), (3, 1), (3, 4)],
         [[[0, 2, 3], [1, 4]], (0, 1), (2, 4), (3, 1), (3, 4)],
         [[[1, 2, 3], [0, 4]], (1, 0), (2, 0), (2, 4), (3, 4)],
         [[[0, 1, 2, 3], [4]], (2, 4), (3, 4)]]
\end{verbatim}

\section{\tt \ \ cut\_matrix}

\textbf{Syntax:} \verb}cut_matrix(Gamma)}\\

\textbf{Description:}
\begin{verbatim}
Returns cut matrix of Gamma.

The cut matrix lists all the possible nontrivial cuts. Each row represents one cut, 
and each column represents an edge in the graph.
1 in a position means that edge is part of the cut.
0 means the edge is not part of the cut.

\end{verbatim}

\textbf{Example:}
\begin{verbatim}
    EXAMPLES:
        sage: Gamma =  graphs.DiamondGraph()
        sage: cut_matrix(Gamma)

        [1 1 0 0 0]
        [1 0 1 1 0]
        [0 1 1 1 0]
        [0 1 1 0 1]
        [1 0 1 0 1]
        [1 1 0 1 1]
        [0 0 0 1 1]
\end{verbatim}


\section{\tt \ \ cut\_oriented\_matrix}

\textbf{Syntax:} \verb}cut_oriented_matrix(Gamma)}\\

\textbf{Description:}
\begin{verbatim}
An oriented cut matrix is like the regular cut matrix, but it takes edge orientation into account.
It's not obvious, but the result is independent of the choice of an orientation vector.

Each row of the matrix corresponds to one oriented cut, and each column to an edge.

Each cut partitions the vertices of the graph into two sets, say V = A \cup B, disjoint and non-empty.

An edge from a vertex in A to a vertex in B contributes +1 to the row.
An edge from B to A contributes -1.
Edges that stay within A or within B contribute 0.

All the matrix entries are in {-1,0,1}, depending on how each edge crosses 
the cut and in what direction.
\end{verbatim}

\textbf{Example:}
\begin{verbatim}
    EXAMPLES:
        sage: Gamma = graphs.CompleteGraph(4)
        sage: cut_matrix_oriented(Gamma,6*[1])

        [ 1  1  1  0  0  0]
        [-1  0  0  1  1  0]
        [ 0  1  1  1  1  0]
        [ 0 -1  0 -1  0  1]
        [ 1  0  1 -1  0  1]
        [-1 -1  0  0  1  1]
        [ 0  0  1  0  1  1]
\end{verbatim}

\section{\tt \ \ cut\_matrix\_oriented}

\textbf{Syntax:} \verb}cut_matrix_oriented(Gamma)}\\

\textbf{Description:}
\begin{verbatim}
Same as \verb+cut_oriented_matrix+.
\end{verbatim}


\section{\tt \ \ incidence\_matrix\_tree\_inverse\_value}

\textbf{Syntax:} \verb}incidence_matrix_tree_inverse_value(Gamma, eo, x, T, e, v)}\\

\textbf{Description:}
\begin{verbatim}
This computes the matrix entries of the inverse of the truncated incidence matrix, 
D(x,T)^{-1}, with edge-orientation vector eo, associated to a vertex x in a spanning tree T.

    INPUT:
        Gamma - graph
        x,v  - distinct vertices of Gamma
        T  - spanning tree of Gamma  
        e  - an edge of Gamma in T
        eo - a vector of 1's and -1's whose length is the number of edges in Gamma
             (ie, the size of Gamma, M)

This is consistent with the inverse matrix example for incidence_matrix_tree(Gamma, eo, x, T).
\end{verbatim}

\textbf{Example:}
\begin{verbatim}
    EXAMPLES:
        sage: Gamma = graphs.PaleyGraph(9)
        sage: E = Gamma.edges(sort=True)
        sage: V = Gamma.vertices(sort=True)
        sage: eo = [1, -1, 1, 1, -1, -1, 1, -1, 1, 1, -1, -1, 1, -1, 1, 1, -1, -1]
        sage: T = Gamma.min_spanning_tree()
        sage: x = V[2]
        sage: [incidence_matrix_tree_inverse_value(Gamma, eo, x, T, e, v) for v in V if v != x]
        [0, 0, 0, 0, 1, 0, 0, 0]
        sage: e = E[0]
        sage: e in T
        True
        sage: [incidence_matrix_tree_inverse_value(Gamma, eo, x, T, e, v) for v in V if v != x]
        [0, -1, 0, 0, -1, -1, 0, 0]

\end{verbatim}

\section{\tt \ \ incidence\_tree\_matrix}

\textbf{Syntax:} \verb}incidence_tree_matrix(Gamma, eo, x, T)}\\

\textbf{Description:}
\begin{verbatim}
This computes the truncated incidence matrix D(x,T) with edge-orientation vector eo, 
associated to a vertex x in a spanning tree T.

    INPUT:
        Gamma - graph
        x  - vertex of Gamma
        T  - spanning tree of Gamma  
        eo - a vector of 1's and -1's whose length is the number of edges in Gamma
             (ie, the size of Gamma)

\end{verbatim}

\textbf{Example:}
\begin{verbatim}
    EXAMPLES:
        sage: Gamma = graphs.PaleyGraph(9)
        sage: E = Gamma.edges(sort=True)
        sage: V = Gamma.vertices(sort=True)
        sage: eo = [1, -1, 1, 1, -1, -1, 1, -1, 1, 1, -1, -1, 1, -1, 1, 1, -1, -1]
        sage: T = Gamma.min_spanning_tree()
        sage: v = V[2]
        sage: incidence_matrix_tree(Gamma, eo, v, T)
        [ 1 -1  0  0  0  0  0  0]
        [-1  0  0  0  0  0  0  0]
        [ 1  0  0 -1  0  0  0  0]
        [ 1  0  0  0  0  0  0 -1]
        [ 0 -1  0  0  1  0  0  0]
        [ 0  1  0  0  0 -1  0  0]
        [ 0  0  1  0  0  0  0  0]
        [ 0  0  0  0  0  0 -1  0]
        sage: DvT^(-1)
        [ 0 -1  0  0  0  0  0  0]
        [-1 -1  0  0  0  0  0  0]
        [ 0  0  0  0  0  0  1  0]
        [ 0 -1 -1  0  0  0  0  0]
        [-1 -1  0  0  1  0  0  0]
        [-1 -1  0  0  0 -1  0  0]
        [ 0  0  0  0  0  0  0 -1]
        [ 0 -1  0 -1  0  0  0  0]
\end{verbatim}

\section{\tt \ \ incidence\_matrix\_tree}

\textbf{Syntax:} \verb}incidence_matrix_tree(Gamma, eo, x, T)}\\

\textbf{Description:}
\begin{verbatim}
Same as \verb+incidence_tree_matrix+.
\end{verbatim}


\section{\tt \ \ fund\_vertex\_cut\_char\_fcn}

\textbf{Syntax:} \verb}fund_vertex_cut_char_fcn(Gamma, v, e, eo)}\\

\textbf{Description:}
\begin{verbatim}
This computes the characteristic function of the fundamental cut defined by 
a vertex v, evaluated at an edge e of a graph Gamma with edge-orientation vector eo. 

    INPUT:
        Gamma - graph
        v  - vertex of Gamma
        e  - edge of Gamma  
        eo - a vector of 1's and -1's whose length is the number of edges in Gamma
             (ie, the size of Gamma, M)

    OUTPUT:
        1, if v is the head of e (defined by eo)
        -1, if v is the tail of e (defined by eo)
        0, otherwise.
\end{verbatim}

\textbf{Example:}
\begin{verbatim}
    EXAMPLES:
        sage: Gamma = graphs.PaleyGraph(9)
        sage: E = Gamma.edges(sort=True)
        sage: V = Gamma.vertices(sort=True)
        sage: eo = [1, -1, 1, 1, -1, -1, 1, -1, 1, 1, -1, -1, 1, -1, 1, 1, -1, -1]
        sage: v = V[2]
        sage: e = E[1]
        sage: fund_vertex_cut_char_fcn(Gamma, v, e, eo)
        1
        sage: [fund_vertex_cut_char_fcn(Gamma, V[2], e, eo) for e in E]
        [0, 1, 0, 0, 1, 0, 0, -1, 1, 0, 0, 0, 0, 0, 0, 0, 0, 0]
        # 3rd column of incidence_matrix(Gamma, eo)
        sage: [fund_vertex_cut_char_fcn(Gamma, v, E[1], eo) for v in V]
        [-1, 0, 1, 0, 0, 0, 0, 0, 0]
        # 2nd row of incidence_matrix(Gamma, eo)
\end{verbatim}

\section{\tt \ \ fund\_cut\_char\_fcn}

\textbf{Syntax:} \verb}fund_cut_char_fcn(Gamma, v, e, eo)}\\

\textbf{Description:}
\begin{verbatim}
This computes the characteristic function of the fundamental cut defined by 
a subset U of vertices, evaluated at an edge e of a graph Gamma with 
edge-orientation vector eo. The cut defined by U is the set of edges having exactly 
one vertex in U.

    INPUT:
        Gamma - graph
        U  - a list of vertices of Gamma
        e  - edge of Gamma  
        eo - a vector of 1's and -1's whose length is the number of edges in Gamma
             (ie, the size of Gamma, M)

    OUTPUT:
        1, if v is the head of e (defined by eo)
        -1, if v is the tail of e (defined by eo)
        0, otherwise.
\end{verbatim}

\textbf{Example:}
\begin{verbatim}
    EXAMPLES:
        sage: Gamma = graphs.PaleyGraph(9)
        sage: E = Gamma.edges(sort=True)
        sage: V = Gamma.vertices(sort=True)
        sage: eo = [1, -1, 1, 1, -1, -1, 1, -1, 1, 1, -1, -1, 1, -1, 1, 1, -1, -1]
        sage: U = [V[2],V[4]]
        sage: e = E[1]
        sage: fund_cut_char_fcn(Gamma, U, e, eo)
        1
        sage: [fund_cut_char_fcn(Gamma, U, e, eo) for e in E]
        [0, 1, -1, 0, 1, 0, 0, -1, 1, -1, 0, 0, 1, -1, 0, 0, 0, 0]
\end{verbatim}

\section{\tt \ \ fund\_cycle\_char\_fcn}

\textbf{Syntax:} \verb}fund_cycle_char_fcn(Gamma, q, e, eo)}\\

\textbf{Description:}
\begin{verbatim}
This computes the characteristic function of the fundamental cycle q, 
evaluated at an edge e of a graph Gamma with edge-orientation vector eo. 

    INPUT:
        Gamma - graph
        q  - fundamental cycle of Gamma
        e  - edge of Gamma  
        eo - a vector of 1's and -1's whose length is the number of edges in Gamma
             (ie, the size of Gamma, M)

    OUTPUT:
        1, if v is the head of e (defined by eo)
        -1, if v is the tail of e (defined by eo)
        0, otherwise.
\end{verbatim}

\textbf{Example:}
\begin{verbatim}
   EXAMPLES:
        sage: Gamma = graphs.PaleyGraph(9)
        sage: Q = Gamma.cycle_basis()
        sage: V = Gamma.vertices(sort=True)
        sage: eo = [1, -1, 1, 1, -1, -1, 1, -1, 1, 1, -1, -1, 1, -1, 1, 1, -1, -1]
        sage: e = E[1]
        sage: [fund_cycle_char_fcn(Gamma, q, e, eo) for q in Q]
        [0, 0, 0, 0, 0, 0, -1, 0, -1, -1]
\end{verbatim}

\section{\tt \ \ fund\_cycle}
\textbf{Syntax:} \verb}fund_cycle(Gamma, T, e, eo)}\\

When you take a graph Gamma and build a spanning tree T from it (a connected graph
with no cycles that includes all the vertices). Now, if you take any edge that
isn't in T and add it back in, it must create a unique cycle in Gamma. This is
the fundamental cycle of that edge with respect to the tree.


\textbf{Description:}
\begin{verbatim}
This computes the fundamental cycle C(T,e), say, associated to a tree T and an edge e of a graph 
Gamma, not in T, with edge-orientation vector eo. 

    INPUT:
        Gamma - graph
        T  - a tree of Gamma
        e  - edge of Gamma not in T
        eo - a vector of 1's and -1's whose length is the number of edges in Gamma
             (ie, the size of Gamma, M)

    OUTPUT:
        C(T,e)
\end{verbatim}

\textbf{Example:}
\begin{verbatim}
    EXAMPLES:
        sage: Gamma = graphs.PaleyGraph(9)
        sage: Q = Gamma.cycle_basis()
        sage: V = Gamma.vertices(sort=True)
        sage: E = Gamma.edges(sort=True)
        sage: eo = [1, -1, 1, 1, -1, -1, 1, -1, 1, 1, -1, -1, 1, -1, 1, 1, -1, -1]
        sage: e = E[9]
        sage: T = Gamma.min_spanning_tree()
        sage: e in T
        False
        sage: fund_cycle(Gamma, T, e, eo)
        [(a, 2, None), (2, 0, None), (0, a + 1, None), (a + 1, a, None)]
\end{verbatim}

\section{\tt \ \ laplacian\_reduced\_matrix}

\textbf{Syntax:} \verb}laplacian_reduced_matrix(Gamma, eo, s)}\\

\textbf{Description:}
\begin{verbatim}
This computes the reduced (vertex) Laplacian of a graph Gamma with edge-orientation vector eo, 
where the row column associated to s has been removed.

    INPUT:
        Gamma - graph having N vertices
        s - a (source) vertex of Gamma
        eo - a vector of 1's and -1's whose length is the number of edges in Gamma
             (ie, the size of Gamma, M)

    OUTPUT:
        (N-1)x(N-1) matrix of integers which agrees with the usual adjacency 
        matrix Laplacian except for the entries associated to the vertex s.
\end{verbatim}

\textbf{Example:}
\begin{verbatim}
    EXAMPLES:
        sage: Gamma = graphs.PaleyGraph(9)
        sage: eo = [1, 1, -1, 1, 1, -1, -1, -1, 1, 1, -1, -1, 1, -1, -1, -1, 1, -1]
        sage: s = V[2]
        sage: LR = laplacian_reduced_matrix(Gamma, eo, s)
        sage: LR
        [ 4 -1  0 -1  0  0  0 -1]
        [-1  4  0  0 -1 -1  0  0]
        [ 0  0  4 -1 -1  0 -1  0]
        [-1  0 -1  4 -1  0  0 -1]
        [ 0 -1 -1 -1  4 -1  0  0]
        [ 0 -1  0  0 -1  4 -1 -1]
        [ 0  0 -1  0  0 -1  4 -1]
        [-1  0  0 -1  0 -1 -1  4]
        sage: LR.dimensions()
        (8, 8)
        sage: A = ZZ^8/LR.image()
        sage: A.cardinality()
        11664
        sage: Gamma.spanning_trees_count()
        11664
        sage: det(LR)
        11664
        sage: factor(11664)
        2^4 * 3^6
\end{verbatim}

\section{\tt \ \ laplacian\_reduced}

\textbf{Syntax:} \verb}laplacian_reduced(Gamma, eo, s)}\\

\textbf{Description:}
\begin{verbatim}
Same as laplacian_reduced_matrix.
\end{verbatim}

\section{\t laplacian\_reduced\_punctured\_matrix}
  
\textbf{Syntax:} \verb}laplacian_reduced_punctured_matrix(Gamma, eo, s, p, q)}\\

\textbf{Description:}
\begin{verbatim}
This computes the reduced, punctured (vertex) Laplacian of a graph Gamma with 
edge-orientation vector eo, where row for p, s are removed and columns for q,s
are removed. It's determinant is equal to (pq||s), in Biggs notation [B], section 13-14.

    INPUT:
        Gamma - graph having N vertices
        s - a (source) vertex of Gamma
        p.q - two vertices, each distinct from s
        eo - a vector of 1's and -1's whose length is the number of edges in Gamma
             (ie, the size of Gamma, M)

    OUTPUT:
        (N-2)x(N-2) matrix of integers which agrees with the usual
        adjacency matrix Laplacian except for the entries associated to the
        vertex s.


    REFERENCES:
        [B] Biggs, Norman, {|it Algebraic potential theory on graphs,} BLMS, 1997.

\end{verbatim}

\textbf{Example:}
\begin{verbatim}
 EXAMPLES:
        sage: Gamma = graphs.PaleyGraph(9)
        sage: eo = [1, 1, -1, 1, 1, -1, -1, -1, 1, 1, -1, -1, 1, -1, -1, -1, 1, -1]
        sage: V = Gamma.vertices(sort=True)
        sage: s = V[2]; p = V[4]; q = V[6]
        sage: LRP = laplacian_reduced_punctured(Gamma, eo, s, p, q)
        sage: LRP
        [ 4 -1  0 -1  0  0 -1]
        [-1  4  0  0 -1 -1  0]
        [ 0  0  4 -1 -1  0  0]
        [-1  0 -1  4 -1  0 -1]
        [ 0 -1  0  0 -1  4 -1]
        [ 0  0 -1  0  0 -1 -1]
        [-1  0  0 -1  0 -1  4]
        sage: LRP.dimensions()
        (7, 7)
        sage: det(LRP)
        2592
\end{verbatim}

\section{\tt \ \ laplacian\_reduced\_punctured}

\textbf{Syntax:} \verb}laplacian_reduced_punctured(Gamma, eo, s, p, q)}\\

\textbf{Description:}
\begin{verbatim}
Same as \verb+laplacian_reduced_punctured_matrix+.
\end{verbatim}


\section{\tt \ \ potential\_induced}

\textbf{Syntax:} \verb}potential_induced(Gamma, e0, j, s, p, q)}\\

\textbf{Description:}
\begin{verbatim}
Returns the potential function f:V -> QQ which is induced by a source
s having magnitude j, with input p and output q and satisfies f(q)=0.
Such an f satisfies the "differential equation" 

                            D^t*f = Ps,

    where 
      * D is the incidence matrix, and
      * P is the orthogonal projection from the space, C^1(Gamma, RR), 
        of functions on E (the edges of Gamma) to the cut space B.
\end{verbatim}

\textbf{Example:}
\begin{verbatim}
    EXAMPLES:
        sage: Gamma = graphs.PaleyGraph(9)
        sage: e0 = [1, 1, -1, 1, 1, -1, -1, -1, 1, 1, -1, -1, 1, -1, -1, -1, 1, -1]
        sage: V = Gamma.vertices(sort=True)
        sage: s = V[2]; p = V[4]; q = V[6]
        sage: j = 5
        sage: potential_induced(Gamma, e0, j, s, p, q)
        2/45
\end{verbatim}

\section{\tt \ \ number\_of\_2\_trees}

\textbf{Syntax:} \verb}number_of_2_trees(Gamma, e0, v, w)}\\

A $2$-tree is a spanning forest of Gamma with two connected components, i.e., a subgraph that
includes all the vertices,
has exactly two connected components (trees),
has no cycles (i.e., it's a forest),
and spans the entire graph.

\textbf{Description:}
\begin{verbatim}
Returns the number of 2-trees for which one component contains v and the other contains w.
\end{verbatim}

\textbf{Example:}
\begin{verbatim}
EXAMPLES:
        sage: Gamma = graphs.PaleyGraph(9)
        sage: eo = [1, 1, -1, 1, 1, -1, -1, -1, 1, 1, -1, -1, 1, -1, -1, -1, 1, -1]
        sage: V = Gamma.vertices(sort=True)
        sage: v = V[4]; w = V[6]
        sage: number_of_2_trees(Gamma, e0, v, w)
        6480
\end{verbatim}

\section{\tt \ \ critical\_group}
\textbf{Syntax:} \verb}critical_group(Gamma)}\\

\textbf{Description:}
\begin{verbatim}
Returns the co-kernel of a reduced Laplacian of the graph Gamma.
\end{verbatim}

\textbf{Example:}
\begin{verbatim}
    EXAMPLES:
        sage: WG6 = graphs.WheelGraph(6)
        sage: critical_group(WG6)
        Finitely generated module V/W over Integer Ring with invariants (11, 11)
        sage: G = graphs.PaleyGraph(9); G
        Paley graph with parameter 9: Graph on 9 vertices
        sage: critical_group(G)
        Finitely generated module V/W over Integer Ring with invariants (6, 6, 18, 18)
        sage: 18*18*6*6
        11664
        sage: G = graphs.CompleteGraph(4); G
        Complete graph: Graph on 4 vertices
        sage: critical_group(G)
        Finitely generated module V/W over Integer Ring with invariants (4, 4)
        sage: WG4 = graphs.WheelGraph(4)
        sage: WG4.delete_edge((2,3))
        sage: WG4.add_vertex(5)
        sage: WG4.add_edges([(2,5),(3,5)])
        sage: critical_group(WG4)
        Finitely generated module V/W over Integer Ring with invariants (24)
        sage: WG4 = graphs.WheelGraph(4)
        sage: critical_group(WG4)
        Finitely generated module V/W over Integer Ring with invariants (4, 4)

        sage: G = DihedralGroup(4)
        sage: G.gens()
        [(1,2,3,4), (1,4)(2,3)]
        sage: S = G.gens()+[G.gens()[0]^(-1)]; S
        [(1,2,3,4), (1,4)(2,3), (1,4,3,2)]
        sage: A = G.cayley_graph(generators=S,simple=True).adjacency_matrix()
        sage: Gamma = Graph(A)
        sage: critical_group(Gamma)
        Finitely generated module V/W over Integer Ring with invariants (2, 8, 24)
\end{verbatim}

\section{\tt \ \ laplacian\_vertex\_matrix}

\textbf{Syntax:} \verb}laplacian_vertex_matrix(Gamma, eo)}\\

\textbf{Description:}
\begin{verbatim}
This computes the vertex Laplacian of a graph Gamma with edge-orientation vector eo. 

    INPUT:
        Gamma - graph having N vertices
        eo - a vector of 1's and -1's whose length is the number of edges in Gamma
             (ie, the size of Gamma, M)

    OUTPUT:
        NxN matrix of 1, -1, 0's (except on the diagonal) which agrees 
        with the usual adjacency matrix Laplacian 
\end{verbatim}

\textbf{Example:}
\begin{verbatim}
    EXAMPLES:
        sage: Gamma = graphs.PaleyGraph(9)
        sage: eo = [1, 1, -1, 1, 1, -1, -1, -1, 1, 1, -1, -1, 1, -1, -1, -1, 1, -1]
        sage: laplacian_vertex(Gamma, eo)
        [ 4 -1 -1  0 -1  0  0  0 -1]
        [-1  4 -1  0  0 -1 -1  0  0]
        [-1 -1  4 -1  0  0  0 -1  0]
        [ 0  0 -1  4 -1 -1  0 -1  0]
        [-1  0  0 -1  4 -1  0  0 -1]
        [ 0 -1  0 -1 -1  4 -1  0  0]
        [ 0 -1  0  0  0 -1  4 -1 -1]
        [ 0  0 -1 -1  0  0 -1  4 -1]
        [-1  0  0  0 -1  0 -1 -1  4]
        sage: Gamma.laplacian_matrix()
        [ 4 -1 -1  0 -1  0  0  0 -1]
        [-1  4 -1  0  0 -1 -1  0  0]
        [-1 -1  4 -1  0  0  0 -1  0]
        [ 0  0 -1  4 -1 -1  0 -1  0]
        [-1  0  0 -1  4 -1  0  0 -1]
        [ 0 -1  0 -1 -1  4 -1  0  0]
        [ 0 -1  0  0  0 -1  4 -1 -1]
        [ 0  0 -1 -1  0  0 -1  4 -1]
        [-1  0  0  0 -1  0 -1 -1  4]
        sage: laplacian_vertex(Gamma, eo).kernel()
        Vector space of degree 9 and dimension 1 over Rational Field
        Basis matrix:
        [1 1 1 1 1 1 1 1 1]
\end{verbatim}


\section{\tt \ \ laplacian\_vertex}
  
\textbf{Syntax:} \verb}laplacian_vertex(Gamma, eo)}\\

\textbf{Description:}
\begin{verbatim}
Same as laplacian_vertex_matrix.
\end{verbatim}


\section{\tt \ \ laplacian\_edge\_matrix}

\textbf{Syntax:} \verb}laplacian_edge_matrix(Gamma, eo)}\\

\textbf{Description:}
\begin{verbatim}
This computes the edge Laplacian of a graph Gamma with edge-orientation vector eo. 

    INPUT:
        Gamma - graph having N vertices
        eo - a vector of 1's and -1's whose length is the number of edges in Gamma
             (ie, the size of Gamma, M)

    OUTPUT:
        MxM matrix of 1, -1, 0's which agrees with the edge
        adjacency matrix Laplacian 
\end{verbatim}

\textbf{Example:}
\begin{verbatim}
    EXAMPLES:
        sage: Gamma = graphs.PaleyGraph(9)
        sage: eo = [1, 1, -1, 1, 1, -1, -1, -1, 1, 1, -1, -1, 1, -1, -1, -1, 1, -1]
        sage: laplacian_edge(Gamma, eo)
        [ 2  1 -1  1 -1  1  1  0  0  0  0  0  0  0  0  0  0  0]
        [ 1  2 -1  1  1  0  0  1 -1  0  0  0  0  0  0  0  0  0]
        [-1 -1  2 -1  0  0  0  0  0 -1  0  0  1 -1  0  0  0  0]
        [ 1  1 -1  2  0  0  0  0  0  0  0  0  0 -1  0  0  1 -1]
        [-1  1  0  0  2 -1 -1  1 -1  0  0  0  0  0  0  0  0  0]
        [ 1  0  0  0 -1  2  1  0  0  0  1  0 -1  0 -1  0  0  0]
        [ 1  0  0  0 -1  1  2  0  0  0  0  0  0  0  1 -1  1  0]
        [ 0  1  0  0  1  0  0  2 -1  1 -1 -1  0  0  0  0  0  0]
        [ 0 -1  0  0 -1  0  0 -1  2  0  0 -1  0  0  0 -1  0  1]
        [ 0  0 -1  0  0  0  0  1  0  2 -1 -1 -1  1  0  0  0  0]
        [ 0  0  0  0  0  1  0 -1  0 -1  2  1 -1  0 -1  0  0  0]
        [ 0  0  0  0  0  0  0 -1 -1 -1  1  2  0  0  0  1  0 -1]
        [ 0  0  1  0  0 -1  0  0  0 -1 -1  0  2 -1  1  0  0  0]
        [ 0  0 -1 -1  0  0  0  0  0  1  0  0 -1  2  0  0 -1  1]
        [ 0  0  0  0  0 -1  1  0  0  0 -1  0  1  0  2 -1  1  0]
        [ 0  0  0  0  0  0 -1  0 -1  0  0  1  0  0 -1  2 -1 -1]
        [ 0  0  0  1  0  0  1  0  0  0  0  0  0 -1  1 -1  2 -1]
        [ 0  0  0 -1  0  0  0  0  1  0  0 -1  0  1  0 -1 -1  2]
\end{verbatim}


\section{\tt \ \ laplacian\_edge}

\textbf{Syntax:} \verb}laplacian_edge(Gamma, eo)}\\

Same as \verb+laplacian_edge_matrix+.


\section{\tt \ \ quotient\_adjacent}

\textbf{Syntax:} \verb}quotient_adjacent(orbits, VG, EG, i, j)}\\

\textbf{Description:}
\begin{verbatim}

This function returns 1 if orbit i is adcacent to orbit j in the quotient graph,
\verb+quotient_graph+.

\end{verbatim}


\section{\tt \ \ quotient\_graph}

\textbf{Syntax:} \verb}quotient_graph(Gamma, G, verbose=False, morphism=False)}\\

\textbf{Description:}
\begin{verbatim}
This function returns the graph quotient, and in verbose mode information on 
the vertex correspondence of the quotient map.

If morphism = True then it returns the associated quotient morphism in CL notation.

If Gamma = (V, E) is a graph and G a subgroup of its 
automorphism group acting faithfully on Gamma, we define 
the quotient graph Gamma/G as follows:
    
(1) The vertices of Gamma/G are the G-orbits in V. 
    
(2) Distinct vertices v1, v2 of Gamma/G are connected by an 
    edge of Gamma/G if and only if there is 
          a vertex v1' in V belonging to the orbit 
             of v1, and 
          a vertex v2' in V belonging to the orbit 
             of v2 (note v1', v2' belong to different $G$-orbits,
             by hypothesis), 
    for which (v1', v2') belongs to E.
    
Note: Gamma/G is simple.
    
    Input: Gamma - a graph in standard form,
           G     - a subgroup of its automorphism group
                   that acts faithfully on Gamma.
                   (for each non-trivial g in G, there is a v in V such that g(v)\not= v)
    Output: The quotient graph.
\end{verbatim}

\textbf{Example:}
\begin{verbatim}
    EXAMPLES:
        sage: Gamma = graphs.PaleyGraph(13)
        sage: G = Gamma.automorphism_group(); G
        Permutation Group with generators [(1,9,3)(2,5,6)(4,10,12)(7,11,8), 
         (1,10,9,12,3,4)(2,7,5,11,6,8), (0,1,2,3,4,5,6,7,8,9,10,11,12)]
        sage: quotient_graph(Gamma, G)
        Graph on 1 vertex
        sage: G2 = G.sylow_subgroup(2)
        sage: G2.order()
        2
        sage: quotient_graph(Gamma, G2)
        Graph on 7 vertices
        sage: G3 = G.sylow_subgroup(3)
        sage: G3.order()
        3
        sage: quotient_graph(Gamma, G3)
        Graph on 5 vertices
        sage: quotient_graph(Gamma, G3, verbose=True)
        (Graph on 5 vertices,
         [0, 1, 2, 3, 4],
         [({8, 11, 7}, {2, 5, 6}),
          ({8, 11, 7}, {1, 3, 9}),
          ({2, 5, 6}, {1, 3, 9}),
          ({0}, {1, 3, 9}),
          ({4, 10, 12}, {8, 11, 7}),
          ({4, 10, 12}, {1, 3, 9}),
          ({4, 10, 12}, {2, 5, 6}),
          ({4, 10, 12}, {0})],
         [{4, 10, 12}, {8, 11, 7}, {2, 5, 6}, {0}, {1, 3, 9}])
\end{verbatim}

\section{\tt \ \ quotient\_graph\_morphism}

\textbf{Syntax:} \verb}quotient_graph_morphism(Gamma, G)}\\

\textbf{Description:}
\begin{verbatim}
This function returns the graph quotient morphism, 
                   f: Gamma -> Gamma/G,
in CL format.

If Gamma = (V, E) is a graph and G a subgroup of its automorphism group 
acting faithfully on Gamma, we define the quotient graph Gamma/G as follows:
    (1) The vertices of Gamma/G are the G-orbits in V. 
    (2) Distinct vertices v1, v2 of Gamma/G are connected by an 
        edge of Gamma/G if and only if there is 
          a vertex v1' in V belonging to the orbit 
             of v1, and 
          a vertex v2' in V belonging to the orbit 
             of v2 (note v1', v2' belong to different $G$-orbits,
             by hypothesis), 
         for which (v1', v2') belongs to E.
    
Note: Gamma/G is simple.
    
    Input: Gamma - a graph in standard form,
           G     - a subgroup of its automorphism group
                   that acts faithfully on Gamma.
                   (for each non-trivial g in G, there is a v in V such that g(v)\not= v)
    Output: The quotient graph morphism.


\end{verbatim}
\textbf{Example:}
\begin{verbatim}
    EXAMPLES:
        sage: Gamma2 = graphs.HouseGraph()  
        sage: G2 = Gamma2.automorphism_group(); G2; G2.cardinality() 
         Permutation Group with generators [(0,1)(2,3)]
         2
        sage: Gamma1 = quotient_graph(Gamma2, G2)  ## P3  
        sage: f21_0 = quotient_graph_morphism(Gamma2, G2); f21_0
         [0, 0, 1, 1, 2]                                                        
        sage: V2 = Gamma2.vertices(sort=True)  
        sage: f21 = [V2, f21_0] 
        sage: is_harmonic_graph_morphism(Gamma1, Gamma2, f21, verbose=False) 
         True
        sage: [horizontal_multiplicity_of_harmonic_map(Gamma1, Gamma2, f21, v) for v in V2] 
         [1, 1, 1, 1, 2]
        sage: [vertical_multiplicity_of_harmonic_map(Gamma1, Gamma2, f21, v) for v in V2]
         [1, 1, 1, 1, 0]
        sage: Gamma2 = graphs.ButterflyGraph()  
        sage: G2 = Gamma2.automorphism_group(); G2; G2.cardinality() 
         Permutation Group with generators [(1,2), (0,1)(2,3)]
         8
        sage: Gamma1 = quotient_graph(Gamma2, G2)     ## P3       
        sage: V2 = Gamma2.vertices(sort=True)   
        sage: f21_0 = quotient_graph_morphism(Gamma2, G2); f21_0   
         [0, 0, 0, 0, 1]
        sage: f21 = [V2, f21_0] 
        sage: is_harmonic_graph_morphism(Gamma1, Gamma2, f21, verbose=False)
         True
        sage: [horizontal_multiplicity_of_harmonic_map(Gamma1, Gamma2, f21, v) for v in V2] 
         [1, 1, 1, 1, 4]
        sage: [vertical_multiplicity_of_harmonic_map(Gamma1, Gamma2, f21, v) for v in V2] 
         [1, 1, 1, 1, 0]
\end{verbatim}

\section{\tt \ \ zeta\_polynomial\_graph}

\textbf{Syntax:} \verb}zeta_polynomial_graph(Gamma)}\\

\textbf{Description:}
\begin{verbatim}
Returns the zeta polynomial of the graph Gamma, a special polynomial that contains 
information about how the graph is connected, using its cycles and spectrum. 
The reciprocal is the Bass-Sunada-Ihara zeta function of Gamma. 
\end{verbatim}

\textbf{Example:}
\begin{verbatim}
    EXAMPLES:
        sage: Gamma1 = graphs.CycleGraph(5)
        sage: zeta_polynomial_graph(Gamma1)
        (x^4 + x^3 + x^2 + x + 1)^2*(x - 1)^2
        sage: Gamma = graphs.DorogovtsevGoltsevMendesGraph(2)
        sage: zeta_polynomial_graph(Gamma)
        -(3*x^4 + x^3 + 3*x^2 + x + 1)^2*(3*x^3 + x^2 + x - 1)*(x^2 - 1)^3*(x - 1)
\end{verbatim}

\section{\tt \ \ star\_subgraph}

\textbf{Syntax:} \verb}star_subgraph(Gamma, v)}\\

\textbf{Description:}
\begin{verbatim}
Returns the star subgraph of Gamma based at the vertex v.
\end{verbatim}

\textbf{Example:}
\begin{verbatim}
    EXAMPLES:
        sage: Gamma = graphs.WheelGraph(4)
        sage: Gamma1 = star_subgraph(Gamma, 0)
        sage: Gamma1.vertices(sort=True)
        [0, 1, 2, 3]
        sage: Gamma1.edges(sort=True)
        [(0, 1, None), (0, 2, None), (0, 3, None)]
\end{verbatim}

\section{\tt \ \ is\_graph\_homomorphism}

\textbf{Syntax:} \verb}is_graph_homomorphism(Gamma1, Gamma2, f, verbose=False)}\\

\textbf{Description:}
\begin{verbatim}
Returns True if f defines a graph-theoretic homomorphism from Gamma2 to Gamma 1(ie, vertices to vertices, 
edges to edges), and False otherwise. 

NOTE:    Gamma1, Gamma2 must be in standard form.

Suppose Gamma2 has n vertices and m edges. A morphism 
    f: Gamma2 -> Gamma1
can represented by 2 pairs of lists [VL2, VL1] and [EL2, EL1],
where 

  - VL2 is the list of all n vertices of Gamma2,
  - VL1 is the list of length n of the vertices in Gamma1 that 
    form the corresponding image under the map f, 
  - EL2 is the list of all m edges of Gamma2,
  - EL1 is the list of length m of the edges in Gamma1 that form 
    the corresponding image under the map f (including None). 
    
However, the pair [EL2, EL1] can be deduced from the pair [VL2, VL1], 
so we drop [EL2, EL1].
\end{verbatim}

\textbf{Example:}
\begin{verbatim}
    EXAMPLES:
        sage: Gamma2 = graphs.WheelGraph(5)
        sage: Gamma1 = graphs.WheelGraph(2)
        sage: f = [[0,1,2,3,4],[0,1,1,1,1]]
	sage: L = [vertical_multiplicity_of_graph_morphism(Gamma1, Gamma2, f, v2) for v2 in V2]; L
        [0, 2, 2]
        sage: is_graph_homomorphism(Gamma1, Gamma2, f)
        False
        sage: Gamma1 = graphs.CycleGraph(4)
        sage: Gamma2 = graphs.PathGraph(3)
        sage: V2 = Gamma2.vertices(sort=True) 
        sage: f0 = [0,1,2,1]  
        sage: f = [V2, f0]
	sage: L = [vertical_multiplicity_of_graph_morphism(Gamma1, Gamma2, f, v2) for v2 in V2]; L
        [0, 0, 0]
        sage: is_graph_homomorphism(Gamma1, Gamma2, f) 
        True
\end{verbatim}

\section{\tt \ \ is\_graph\_morphism}

\textbf{Syntax:} \verb}is_graph_morphism(Gamma1, Gamma2, f, verbose=False)}\\

\textbf{Description:}
\begin{verbatim}
Returns True if f defines a graph-theoretic mapping from Gamma2 to Gamma 1 (ie, vertices to vertices, 
edges to edges or vertices), and False otherwise. 

NOTE:    Gamma1, Gamma2 must be in standard form.

Suppose Gamma2 has n vertices and m edges. A morphism 
    f: Gamma2 -> Gamma1
can represented by 2 pairs of lists [VL2, VL1] and [EL2, EL1],
where 

  - VL2 is the list of all n vertices of Gamma2,
  - VL1 is the list of length n of the vertices in Gamma1 that 
    form the corresponding image under the map f, 
  - EL2 is the list of all m edges of Gamma2,
  - EL1 is the list of length m of the edges in Gamma1 that form 
    the corresponding image under the map f (including None). 
\end{verbatim}

\textbf{Example:}
\begin{verbatim}
    EXAMPLES:
        sage: Gamma2 = graphs.WheelGraph(5)
        sage: Gamma1 = graphs.WheelGraph(2)
        sage: f = [[0,1,2,3,4],[0,1,1,1,1]]
        sage: is_graph_morphism(Gamma1, Gamma2, f)
        True
        sage: Gamma1 = graphs.WheelGraph(5)
        sage: Gamma2 = graphs.WheelGraph(6)
        sage: V2 = Gamma2.vertices(sort=True) 
        sage: f0 = [0,1,2,3,4,4]  
        sage: f = [V2, f0]   
        sage: is_graph_morphism(Gamma1, Gamma2, f) 
        True
        sage: f0 = [0,1,4, 2,3,4] 
        sage: f = [V2, f0]   
        sage: is_graph_morphism(Gamma1, Gamma2, f, verbose=True) 
        Gamma1 = (V1, E1) =  [0, 1, 2, 3, 4]
          [(0, 1), (0, 2), (0, 3), (0, 4), (1, 2), (1, 4), (2, 3), (3, 4)]
        Gamma2 = (V2, E2) =  [0, 1, 2, 3, 4, 5]
          [(0, 1), (0, 2), (0, 3), (0, 4), (0, 5), (1, 2), (1, 5),
           (2, 3), (3, 4), (4, 5)]
        f = [V2, f0], where f0 =  [0, 1, 4, 2, 3, 4]
        The edge e =  (2, 3)  has no image in Gamma1 under f.
        False
\end{verbatim}

\section{\tt \ \  edge\_image\_list}

\textbf{Syntax:} \verb}edge_image_list(Gamma1, Gamma2, f)}\\

\textbf{Description:}
\begin{verbatim}
Returns a list of the images of the edges under the graph morphism f:Gamma2->Gamma1 in color list notation.

\end{verbatim}

Represent the vertices of $\Gamma_1$ by $n_1$ distinct colors.
These colors can be ``pulled back'' to $\Gamma_2$ via $\phi$:
for each $j \in V_1=\{0,1,\dots,n_1-1\}$, we say
vertex $i$ of $\Gamma_2$ has {\it color} $j$
if $\phi(i)=j$. Call an $n_2$-tuple of $V_1$
a {\it color list} (CL) if it's attached to
a graph morphism $\phi:\Gamma_2\to\Gamma_1$ as above.

\textbf{Example:}
\begin{verbatim}
    EXAMPLES:
        sage: Gamma2 = graphs.WheelGraph(5)
        sage: Gamma1 = graphs.WheelGraph(2)
        sage: f = [[0,1,2,3,4],[0,1,1,1,1]]
        sage: is_graph_morphism(Gamma1, Gamma2, f)
        True
        sage: edge_image_list(Gamma1, Gamma2, f)
        [[(0, 1), (0, 2), (0, 3), (0, 4), (1, 2), (1, 4), (2, 3), (3, 4)],
         [(0, 1), (0, 1), (0, 1), (0, 1), 'None', 'None', 'None', 'None']]
\end{verbatim}

\section{\tt \ \  number\_of\_edges\_set\_in\_a\_walk}

\textbf{Syntax:} \verb}number_of_edges_set_in_a_walk(Gamma, W, verbose= False)}\\

\textbf{Description:}
\begin{verbatim}
INPUT:
    Gamma = (V,E) is a graph
    W     = a sequence of incident unlabeled, "directed" edges in E

OUTPUT: 
    The cardinality of the set of edges in W, if W is a walk,
    False, if W is not a walk in Gamma
\end{verbatim}

\textbf{Example:}
\begin{verbatim}
    EXAMPLES:
        sage: V = [0, 1, 2, 3, 4, 5, 6, 7, 8]
        sage: E = [(0,1),(0,2),(0,7),(0,8),(1,2),(1,3),(1,7),(2,3),(2,8), (3, 4), (3, 5), (4, 5), (4, 6), (4, 8), (5, 6), (5, 7), (6, 7), (6, 8)]
        sage: Gamma = Graph([V,E])
        sage: W = [(0,1), (1,7), (7,5), (5,7)]
        sage: number_of_edges_set_in_a_walk(Gamma, W)
        3
\end{verbatim}

\section{\tt \ \  edge\_set\_is\_a\_walk}
  
\textbf{Syntax:} \verb}edge_set_is_a_walk(Gamma, W, verbose = False)}\\

\textbf{Description:}
\begin{verbatim}
INPUT:
    Gamma = (V,E) is a graph
    W     = a sequence of incident unlabeled, "directed" edges in E

OUTPUT: 
    True if W is a walk in Gamma and False otherwise
\end{verbatim}

\textbf{Example:}
\begin{verbatim}
    EXAMPLES:
        sage: V2 = [0, 1, 2, 3, 4, 5, 6, 7, 8]
        sage: E2 = [(0,1),(0,2),(0,7),(0,8),(1,2),(1,3),(1,7),(2,3),(2,8), (3, 4), (3, 5), (4, 5), (4, 6), (4, 8), (5, 6), (5, 7), (6, 7), (6, 8)]
        sage: Gamma2 = Graph([V2,E2])
        sage: W = [(0,1), (1,7), (7,5), (5,0)] ## not a walk in Gamma2
        sage: edge_set_is_a_walk(Gamma2, W, verbose= True)
        Edge (5, 0) is not in the graph
        False
        sage: W = [(0,1), (1,7), (7,5)]
        sage: edge_set_is_a_walk(Gamma2, W, verbose= True)
        True
        sage: edge_set_is_a_walk(Gamma2, W, verbose= False)
        True
\end{verbatim}

\section{\tt \ \  image\_of\_walk\_under\_graph\_morphism}

\textbf{Syntax:} \verb}image_of_walk_under_graph_morphism(Gamma1, Gamma2, f, W)}\\

\textbf{Description:}
INPUT:
    Gamma1, Gamma2 -- graphs
          W        -- a sequence of incident "directed" edges in E2
          f        -- a graph morphism f: Gamma2 -> Gamma1 represented by a
                      color list, f = [V2,f(V2)]
     
OUTPUT:
    This function returns the walk in Gamma1 represented
    as a list of incident edges in E1 given by f(W).
    Returns False if W is not a walk.

\textbf{Example:}
\begin{verbatim}
EXAMPLES:
        sage: Gamma1 = graphs.CycleGraph(3)
        sage: V2 = [0, 1, 2, 3, 4, 5, 6, 7, 8]
        sage: E2 = [(0,1),(0,2),(0,7),(0,8),(1,2),(1,3),(1,7),(2,3),(2,8), (3, 4), (3, 5), (4, 5), (4, 6), (4, 8), (5, 6), (5, 7), (6, 7), (6, 8)]
        sage: Gamma2 = Graph([V2,E2])
        sage: f = [V2, [2, 1, 0, 2, 0, 1, 2, 0, 1]]
        sage: is_harmonic_graph_morphism(Gamma1, Gamma2, f)
        True
        sage: [vertical_multiplicity_of_harmonic_map(Gamma1, Gamma2, f, v) for v in Gamma2.vertices(sort=True)]
        [0, 0, 0, 0, 0, 0, 0, 0, 0]
        sage: [horizonal_multiplicity_of_harmonic_map(Gamma1, Gamma2, f, v) for v in Gamma2.vertices(sort=True)]
        [2, 2, 2, 2, 2, 2, 2, 2, 2]
        sage: W = [(0,1), (1,7), (7,5)]
        sage: image_of_walk_under_graph_morphism(Gamma1, Gamma2, f, W)
        [(2, 1), (1, 0), (0, 1)]
        sage: W = [(0,1), (1,7), (7,5), (5,6), (6,8), (8,0)]
        sage: image_of_walk_under_graph_morphism(Gamma1, Gamma2, f, W)
        [(2, 1), (1, 0), (0, 1), (1, 2), (2, 1), (1, 2)]
        sage: W = [(0,1), (1,7), (7,5), (5,0)] ## not a walk in Gamma2
        sage: image_of_walk_under_graph_morphism(Gamma1, Gamma2, f, W)
        False
\end{verbatim}

\section{\tt \ \  index\_of\_graph\_morphism\_edges}
\textbf{Syntax:} \verb}index_of_graph_morphism_edges(Gamma1, Gamma2, f, e0, s0, iota2, iota1)}\\

\textbf{Description:}
\begin{verbatim}
These are the entries of the signed edge matrix, matrix_of_graph_morphism_edges,
attached to the harmonic map f : Gamma2 -> Gamma1.

Here     s0 is an edge in Gamma1 (assumes s0[0]<s0[1]),
         e0 is an edge in Gamma2 (assumes e0[0]<e0[1]),
         iota1 is an orientation on Gamma1,
         iota2 is an orientation on Gamma2. 
    
Recall an orientation iota on the edges E of Gamma is with respect to the 
standard orientation (all edges are represented e = (x,y) with x<y, 
in which case we have iota(e)=+1). 
An orientation is represented by a list of 1, -1's of length |E|).
If an edge e=(x,y) has iota(e)=-1 then x>y (which can be represented as (y,x), if desired).
    
If f does not send e0 to s0 then the index is 0.
If f sends e0 to s0 and sends the head of e0 (wrt iota2) to the head of s0 (wrt iota1) 
then the index is +1.
If f sends e0 to s0 and sends head to tail then the index is +1.

For examples, see matrix_of_graph_morphism_edges.
\end{verbatim}

\section{\tt \ \  index\_of\_graph\_morphism\_edges\_unsigned}
  
\textbf{Syntax:} \verb}index_of_graph_morphism_edges_unsigned(Gamma1, Gamma2, f, e2, e1)}\\

\textbf{Description:}
\begin{verbatim}
    These are the entries of the signed edge matrix attached to the 
    harmonic map f : Gamma2 -> Gamma1.

    Here e1 is an edge in Gamma1 and e2 is an edge in Gamma2.

    This function helps answer a very specific question:

           “Does this edge e2 in the big graph Gamma2 get mapped to this 
            edge e1 in the smaller graph Gamma1 under the function f?”

    It returns:
      --  1 if yes (it matches — in either direction),
      --  0 if no,
    and raises an error if e2 isn't an edge in the domain of the morphism.

    Called the unsigned edge incidence — it doesn't care about direction, just whether the edge gets mapped or not.
\end{verbatim}

\textbf{Example:}
\begin{verbatim}
See the matrix entries in the example for graph_morphism_edges_matrix_unsigned(Gamma1, Gamma2, f).
\end{verbatim}

\section{\tt \ \ index\_of\_graph\_morphism\_vertices}
  
\textbf{Syntax:} \verb}index_of_graph_morphism_vertices(Gamma1, Gamma2, f, v2, v1)}\\

\textbf{Description:}
\begin{verbatim}
See the matrix entries in the example for matrix_of_graph_morphism_vertices(Gamma1, Gamma2, f, verbose = False)
\end{verbatim}

\section{\tt \ \  graph\_morphism\_vertices\_matrix}

\textbf{Syntax:} \verb}graph_morphism_vertices_matrix(Gamma1, Gamma2, f)}\\

\textbf{Description:}
\begin{verbatim}
See matrix_of_graph_morphism_vertices(Gamma1, Gamma2, f).
\end{verbatim}

\section{\tt \ \  matrix\_of\_graph\_morphism\_vertices}

\textbf{Syntax:} \verb}matrix_of_graph_morphism_vertices(Gamma1, Gamma2, f, verbose = False)}\\

\textbf{Description:}
\begin{verbatim}
Returns an n2 x n1 (0,1)-matrix representing Phi_V if f defines a graph-theoretic morphism
from Gamma2 to Gamma 1, and False otherwise.

Suppose Gamma2 has n vertices and m edges. A morphism 
    f: Gamma2 -> Gamma1
can represented by 2 pairs of lists [VL2, VL1] and [EL2, EL1],
where
    - VL2 is the list of all n vertices of Gamma2,
    - VL1 is the list of length n of the vertices
      in Gamma1 that form the corresponding image under
       the map f,
    - EL2 is the list of all m edges of Gamma2,
    - EL1 is the list of length m of the edges
      in Gamma1 that form the corresponding image under
      the map f (including None).

However, the pair [EL2, EL1] can be deduced from the pair [VL2, VL1].
\end{verbatim}

\textbf{Example:}
\begin{verbatim}
    EXAMPLES:
        sage: Gamma2 = graphs.CubeGraph(2)
        sage: Gamma1 = Gamma2.subgraph(vertices = ['00', '01'], edges = [('00', '01')])
        sage: f = [['00', '01', '10', '11'], ['00', '01', '00', '01']]
        sage: is_graph_morphism(Gamma1, Gamma2, f)
        True
        sage: Phi = matrix_of_graph_morphism_vertices(Gamma1, Gamma2, f); Phi
        [1 0 1 0]
        [0 1 0 1]
        sage: Phi.transpose()*Phi
        [1 0 1 0]
        [0 1 0 1]
        [1 0 1 0]
        [0 1 0 1]
        sage: Phi*Phi.transpose()
        [2 0]
        [0 2]
        sage: Gamma2 = graphs.CubeGraph(3)
        sage: Gamma1 = graphs.TetrahedralGraph()
        sage: f = [['000', '001', '010', '011', '100', '101', '110', '111'], [0, 1, 2, 3, 3, 2, 1, 0]]
        sage: Phi = matrix_of_graph_morphism_vertices(Gamma1, Gamma2, f); Phi
        [1 0 0 0 0 0 0 1]
        [0 1 0 0 0 0 1 0]
        [0 0 1 0 0 1 0 0]
        [0 0 0 1 1 0 0 0]
        sage: D = vector([-1,0,1,2,2,1,0,-1])
        sage: Phi*D
        (-2, 0, 2, 4)
        sage: Phi*Phi.transpose()
        [2 0 0 0]
        [0 2 0 0]
        [0 0 2 0]
        [0 0 0 2]
        sage: Phi.transpose()*Phi
        [1 0 0 0 0 0 0 1]
        [0 1 0 0 0 0 1 0]
        [0 0 1 0 0 1 0 0]
        [0 0 0 1 1 0 0 0]
        [0 0 0 1 1 0 0 0]
        [0 0 1 0 0 1 0 0]
        [0 1 0 0 0 0 1 0]
        [1 0 0 0 0 0 0 1]
        sage: Gamma2 = graphs.WheelGraph(5)
        sage: Gamma1 = graphs.WheelGraph(2)
        sage: f = [[0,1,2,3,4],[0,1,1,1,1]]
        sage: Phi = matrix_of_graph_morphism_vertices(Gamma1, Gamma2, f); Phi
        [1 0 0 0 0]
        [0 1 1 1 1]
        sage: Phi.transpose()*Phi
        [1 0 0 0 0]
        [0 1 1 1 1]
        [0 1 1 1 1]
        [0 1 1 1 1]
        [0 1 1 1 1]
        sage: Phi*Phi.transpose()
        [1 0]
        [0 4]
        sage: Q1 = Gamma1.laplacian_matrix()
        sage: Q2 = Gamma2.laplacian_matrix()
        sage: Q1*Phi
        [ 1 -1 -1 -1 -1]
        [-1  1  1  1  1]
        sage: Phi*Q2
        [ 4 -1 -1 -1 -1]
        [-4  1  1  1  1]
        sage: Gamma3 = graphs.CycleGraph(4)
        sage: Gamma3.add_edge([0,2])
        sage: f2 = [[0,1,2,3,4],[0,2,1,2,3]]
        sage: Phi2 = matrix_of_graph_morphism_vertices(Gamma3, Gamma2, f2); Phi2
        [1 0 0 0 0]
        [0 0 1 0 0]
        [0 1 0 1 0]
        [0 0 0 0 1]
        sage: Phi2.transpose()*Phi2
        [1 0 0 0 0]
        [0 1 0 1 0]
        [0 0 1 0 0]
        [0 1 0 1 0]
        [0 0 0 0 1]
        sage: Phi2*Phi2.transpose()
        [1 0 0 0]
        [0 1 0 0]
        [0 0 2 0]
        [0 0 0 1]
        sage: v = Q3.right_eigenspaces()[0][1].basis()[0]; v
        (0, 1, 0, -1)
        sage: Q3*v ## Does Phi2 send eigenspaces to eigenspaces?
        (0, 2, 0, -2)
        sage: Phi2.transpose()*Phi2*Q2*Phi2.transpose()*v
        (0, 0, 3, 0, -3)
        sage: w = Q3.right_eigenspaces()[2][1].basis()[0]; w
        (1, 0, -1, 0)
        sage: Q3*w # A: No.
        (4, 0, -4, 0)
        sage: Phi2.transpose()*Phi2*Q2*Phi2.transpose()*w
        (6, -8, 1, -8, 1)
\end{verbatim}

\section{\tt \ \ graph\_morphism\_edges\_matrix}

\textbf{Syntax:} \verb}graph_morphism_edges_matrix(Gamma1, Gamma2, f, iota2, iota1)}\\

\textbf{Description:}
\begin{verbatim}
Returns an m2 x m1 (0,1, -1)-matrix representing Phi_E if f defines a graph morphism
from Gamma2 to Gamma, and False otherwise. 

    Here iota1 is an orientation on Gamma1,
         iota2 is an orientation on Gamma2. 
\end{verbatim}

\textbf{Example:}
\begin{verbatim}
    EXAMPLES:
        sage: Gamma1 = graphs.PathGraph(3)
        sage: Gamma2 = graphs.HouseGraph()
        sage: V2 = Gamma2.vertices()
        sage: f0 = [0, 0, 1, 1, 2]
        sage: f = [V2, f0]
        sage: Phi_E_signed = matrix_of_graph_morphism_edges(Gamma1, Gamma2, f, iota2=[1]*6, iota1=[1,-1]); Phi_E_signed
	[ 0  0]
	[ 1  0]
	[ 1  0]
	[ 0  0]
	[ 0 -1]
	[ 0 -1]
        sage: Gamma2 = graphs.CycleGraph(4)
        sage: Gamma1 = graphs.PathGraph(2)
        sage: V2 = Gamma2.vertices()
        sage: f0 = [0, 1, 1, 0]
        sage: f = [V2, f0]
        sage: PhiE = matrix_of_graph_morphism_edges(Gamma1, Gamma2, f, iota2=[1]*4, iota1=[1]); PhiE
        [1]
        [0]
        [0]
        [-1]
        sage: PhiE.transpose()*PhiE
        [2]

        sage: Gamma2 = graph_convert_to_standard(graphs.PaleyGraph(9))
        sage: Gamma1 = graphs.CycleGraph(3)
        sage: E1 = Gamma1.edges(labels=None)
        sage: E2 = Gamma2.edges(labels=None)
        sage: V2 = Gamma2.vertices(sort=True)
        sage: iota2 = [1, 1, 1, 1, 1, 1, 1, 1, 1, 1, 1, 1, 1, 1, 1, 1, 1, 1]
        sage: iota1 = [1,1,1]
        sage: B1 = incidence_matrix(Gamma1, [E1,iota1])
        sage: B2 = incidence_matrix(Gamma2, [E2,iota2])
        sage: f = [V2, [0,0,0,2,2,2,1,1,1]]
        sage: is_harmonic_graph_morphism(Gamma1, Gamma2, f) 
        True
        sage: PhiV = matrix_of_graph_morphism_vertices(Gamma1, Gamma2, f)
        sage: PhiE = matrix_of_graph_morphism_edges(Gamma1, Gamma2, f, iota2, iota1)
        sage: PhiE*B1.transpose() == B2.transpose()*PhiV   
        True
\end{verbatim}

\section*{\tt matrix\_of\_graph\_morphism\_edges}

\textbf{Syntax:} \verb}matrix_of_graph_morphism_edges(Gamma1, Gamma2, f, iota2, iota1)}\\

\textbf{Description:}
\begin{verbatim}
Same as graph_morphism_edges_matrix(Gamma1, Gamma2, f, iota2, iota1)
\end{verbatim}

\section{\tt \ \ graph\_morphism\_edges\_matrix\_unsigned}

\textbf{Syntax:} \verb}graph_morphism_edges_matrix_unsigned(Gamma1, Gamma2, f)}\\

\textbf{Description:}
\begin{verbatim}
Returns an m2 x m1 (0,1)-matrix representing Phi_E if f defines a graph-theoretic mapping
from Gamma2 to Gamma 1 (ie, vertices to vertices, edges to edges), and False otherwise. 

    Here e0 is an orientation on Gamma1. 
\end{verbatim}

\textbf{Example:}
\begin{verbatim}
    EXAMPLES:
        sage: Gamma3 = graphs.CycleGraph(4)
        sage: Gamma2 = graphs.WheelGraph(5)
        sage: Gamma3.add_edge([0,2])
        sage: f2 = [[0,1,2,3,4],[0,2,1,2,3]]
        sage: PhiE2 = matrix_of_graph_morphism_edges_unsigned(Gamma3, Gamma2, f2); PhiE2.transpose()
        [ 0  1  0  0  0  0  0  0]
        [ 1  0  1  0  0  0  0  0]
        [ 0  0  0  1  0  0  0  0]
        [ 0  0  0  0  1  0  1  0]
        [ 0  0  0  0  0  1  0  1]
\end{verbatim}

\section{\tt \ \ matrix\_of\_graph\_morphism\_edges\_unsigned}

\textbf{Syntax:} \verb}matrix_of_graph_morphism_edges_unsigned(Gamma1, Gamma2, f)}\\

\textbf{Description:}
\begin{verbatim}
Same as graph_morphism_edges_matrix_unsigned(Gamma1, Gamma2, f).
\end{verbatim}


\section{\tt \ \ preimage\_of\_vertex\_under\_graph\_morphism}

\textbf{Syntax:} \verb}preimage_of_vertex_under_graph_morphism(Gamma1, Gamma2, f, v)}\\

\textbf{Description:}
\begin{verbatim}
Returns the (vertex) preimage(s) of v under f:Gamma2 -> Gamma 1, where v is an vertex in Gamma1.

A morphism  f: Gamma2 -> Gamma1 is represented as a color list by a pair of lists [L2, L1],
where 
       L2 is the list of all n vertices of Gamma2,
       L1 is the list of length n of the vertices in Gamma1 that form the corresponding image under the map f.
\end{verbatim}

\textbf{Example:}
\begin{verbatim}
    EXAMPLES:
        sage: Gamma2 = graphs.CubeGraph(2)
        sage: Gamma1 = Gamma2.subgraph(vertices = ['00', '01'], edges = [('00', '01')])
        sage: f = [['00', '01', '10', '11'], ['00', '01', '00', '01']]
        sage: Gamma1.vertices(sort=True)
        [('00', '01', None)]
        sage: v = Gamma1.vertices(sort=True)[0]
        sage: preimage_of_vertex_under_graph_morphism(Gamma1, Gamma2, f, v)
        ['00', '10']
\end{verbatim}

\section{\tt \ \ preimage\_of\_edge\_under\_graph\_morphism}
\textbf{Syntax:} \verb}preimage_of_edge_under_graph_morphism(Gamma1, Gamma2, f, e, labels=None)}\\

\textbf{Description:}
\begin{verbatim}
Returns the preimage(s) of e under f:Gamma2 -> Gamma1, where e is an edge in Gamma1. 
It is returned as a list, with the default being that the edges are unlabled.
\end{verbatim}

\textbf{Example:}
\begin{verbatim}
    EXAMPLES:
        sage: Gamma2 = graphs.CubeGraph(2)
        sage: Gamma1 = Gamma2.subgraph(vertices = ['00', '01'], edges = [('00', '01')])
        sage: f = [['00', '01', '10', '11'], ['00', '01', '00', '01']]
        sage: Gamma2 = graphs.CubeGraph(2)
        sage: Gamma1.edges(sort=True)
        [('00', '01', None)]
        sage: e = Gamma1.edges(sort=True)[0]
        sage: preimage_of_edge_under_graph_morphism(Gamma1, Gamma2, f, e, labels = True)
        [('00', '01', None), ('10', '11', None)]
        sage: Gamma2 = graphs.CubeGraph(3)
        sage: Gamma1 = graphs.TetrahedralGraph()
        sage: f = [['000', '001', '010', '011', '100', '101', '110', '111'], [0, 1, 2, 3, 3, 2, 1, 0]]
        sage: e = Gamma1.edges(sort=True)[1]; e
        (0, 2, None)
        sage: preimage_of_edge_under_graph_morphism(Gamma1, Gamma2, f, e, labels = True)
        [('000', '010', None), ('101', '111', None)]
        sage: preimage_of_edge_under_graph_morphism(Gamma1, Gamma2, f, e)   
         [('000', '010'), ('101', '111')]
        sage: preimage_of_edge_under_graph_morphism(Gamma1, Gamma2, f, e, labels = None) 
         [('000', '010'), ('101', '111')]
        sage: Gamma3 = graphs.PathGraph(2)             
        sage: Gamma4 = graphs.CycleGraph(4)             
        sage: V4 = Gamma4.vertices(sort=True)                   
        sage: g0 = [0,1,1,0]                            
        sage: g = [V4, g0]                             
        sage: is_harmonic_graph_morphism(Gamma3, Gamma4, g) 
        True
        sage: image_of_edge_under_graph_morphism(Gamma3, Gamma4, g, (0,3))
        0
        sage: image_of_edge_under_graph_morphism(Gamma3, Gamma4, g, (0,1))
        (0, 1)
\end{verbatim}

\section{\tt \ \ preimage\_of\_edges\_under\_graph\_morphism}

\textbf{Syntax:} \verb}preimage_of_edges_under_graph_morphism(Gamma1, Gamma2, f, S)}\\

\textbf{Description:}
\begin{verbatim}
Returns the preimage of S under f:Gamma2 -> Gamma 1, where S is a list of edges in Gamma1. 
It is returned as a list, without repetition.

\end{verbatim}

\textbf{Example:}
\begin{verbatim}
    EXAMPLES:
        sage: Gamma2 = graphs.CubeGraph(3)
        sage: V2 = Gamma2.vertices(sort=True); V2  
        ['000', '001', '010', '011', '100', '101', '110', '111']
        sage: E2 = Gamma2.edges(labels=None,sort=True); E2 
        [('000', '001'), ('000', '010'), ('000', '100'), ('001', '011'),
         ('001', '101'), ('010', '011'), ('010', '110'), ('011', '111'),
         ('100', '101'), ('100', '110'), ('101', '111'), ('110', '111')]
        sage: Gamma1 = graphs.TetrahedralGraph() 
        sage: E1 = Gamma1.edges(labels=None,sort=True); E1       
        [(0, 1), (0, 2), (0, 3), (1, 2), (1, 3), (2, 3)]
        sage: f  = [['000', '001', '010', '011', '100', '101', '110', '111'], [0, 1, 2, 3, 3, 2, 1, 0]] 
        sage: S = [(0, 2), (0, 3), (1, 2), (1, 3)]         
        sage: preimage_of_edges_under_graph_morphism(Gamma1, Gamma2, f, S) 
        [('100', '110'),
         ('000', '010'),
         ('001', '101'),
         ('011', '111'),
         ('101', '111'),
         ('010', '110'),
         ('001', '011'),
         ('000', '100')]
\end{verbatim}

\section{\tt \ \ graph\_preimage\_of\_edges\_under\_graph\_morphism}

\textbf{Syntax:} \verb}graph_preimage_of_edges_under_graph_morphism(Gamma1, Gamma2, f, S)}\\

\textbf{Description:}
\begin{verbatim}
Returns the preimage of S under f:Gamma2 -> Gamma 1, where S is a list of edges in Gamma1. 
It is returned as a subgraph of Gamma2.
\end{verbatim}

\textbf{Example:}
\begin{verbatim}
    EXAMPLES:
        sage: Gamma2 = graphs.CubeGraph(3)
        sage: V2 = Gamma2.vertices(sort=True); V2  
        ['000', '001', '010', '011', '100', '101', '110', '111']
        sage: E2 = Gamma2.edges(labels=None,sort=True); E2 
        [('000', '001'), ('000', '010'), ('000', '100'), ('001', '011'),
         ('001', '101'), ('010', '011'), ('010', '110'), ('011', '111'),
         ('100', '101'), ('100', '110'), ('101', '111'), ('110', '111')]
        sage: Gamma1 = graphs.TetrahedralGraph() 
        sage: E1 = Gamma1.edges(labels=None,sort=True); E1       
        [(0, 1), (0, 2), (0, 3), (1, 2), (1, 3), (2, 3)]
        sage: f  = [['000', '001', '010', '011', '100', '101', '110', '111'], [0, 1, 2, 3, 3, 2, 1, 0]] 
        sage: S = [(0, 2), (0, 3), (1, 2), (1, 3)]       
        sage: graph_preimage_of_edges_under_graph_morphism(Gamma1, Gamma2, f, S) 
         Subgraph of (3-Cube): Graph on 8 vertices
\end{verbatim}

\section{\tt \ \ image\_of\_edge\_under\_graph\_morphism}

\textbf{Syntax:} \verb}image_of_edge_under_graph_morphism(Gamma1, Gamma2, f, e, verbose=False)}\\

\textbf{Description:}
\begin{verbatim}
Returns the image of e under f:Gamma2 -> Gamma 1, where e is an edge in Gamma2.
\end{verbatim}

\textbf{Example:}
\begin{verbatim}
    EXAMPLES:
        sage: Gamma2 = graphs.CubeGraph(2)
        sage: Gamma1 = Gamma2.subgraph(vertices = ['00', '01'], edges = [('00', '01')])
        sage: f = [['00', '01', '10', '11'], ['00', '01', '00', '01']]
        sage: Gamma2 = graphs.CubeGraph(2)
        sage: Gamma1.edges(sort=True)
        [('00', '01', None)]
        sage: e = Gamma1.edges(sort=True)[0]; e
        ('00', '01', None)
        sage: image_of_edge_under_graph_morphism(Gamma1, Gamma2, f, e)
        ('00', '01', None)
        sage: e = Gamma2.edges(sort=True)[1]; e
        ('00', '10', None)
        sage: image_of_edge_under_graph_morphism(Gamma1, Gamma2, f, e)
        ('00', '00')
\end{verbatim}

\section{\tt \ \ image\_of\_vertex\_under\_graph\_morphism}

\textbf{Syntax:} \verb}image_of_vertex_under_graph_morphism(Gamma1, Gamma2, f, v)}\\

\textbf{Description:}
\begin{verbatim}
Returns the image of v under f:Gamma2 -> Gamma 1, where v is an vertex in Gamma2.
\end{verbatim}

\textbf{Example:}
\begin{verbatim}
    EXAMPLES:
        sage: Gamma2 = graphs.CubeGraph(2)
        sage: Gamma1 = Gamma2.subgraph(vertices = ['00', '01'], edges = [('00', '01')])
        sage: f = [['00', '01', '10', '11'], ['00', '01', '00', '01']]
        sage: v = Gamma2.vertices(sort=True)[0]
        sage: image_of_vertex_under_graph_morphism(Gamma1, Gamma2, f, v)
        '00'
\end{verbatim}

\section{\tt \ \ local\_horizontal\_multiplicity\_of\_harmonic\_map}

\textbf{Syntax:} \verb}local_horizontal_multiplicity_of_harmonic_map(Gamma1, Gamma2, f, v2, e)}\\

\textbf{Description:}
\begin{verbatim}
Returns mu_{phi}(v2, e)
        = \mu_{\phi}(v,e)
        = \sum_{u\in N_{\Gamma_2}(v), \phi(u)=y} 1,
if e=(\phi(u), \phi(v)).

NOTE:  Gamma1 is a graph in standard form and so is Gamma2.

    f : defines a graph-theoretic mapping
         from Gamma2 to Gamma1 represented by a color list as usual,
    v2 : a vertex of Gamma2
    e : an edge of Gamma1 incident to phi(v2)
\end{verbatim}

\textbf{Example:}
\begin{verbatim}
    EXAMPLES:
        sage: V2 = [0,1,2,3]
        sage: E2 = [(0, 1), (1, 2), (0, 2), (0, 3)]
        sage: Gamma2 = Graph([V2,E2])
        sage: V1 = [0,1,2]
        sage: E1 = [(0, 1), (1, 2), (0, 2)]
        sage: Gamma1 = Graph([V1, E1])
        sage: f = [[0,1,2,3], [0, 1, 2, 0]]
        sage: e = (1, 2)
        sage: [local_horizontal_multiplicity_of_harmonic_map(Gamma1, Gamma2, f, v, e) for v in V2]
         [0, 1, 1, 0]
\end{verbatim}

\section{\tt \ \ weighted\_local\_horizontal\_multiplicity\_of\_harmonic\_map}

\textbf{Syntax:} \verb}weighted_local_horizontal_multiplicity_of_harmonic_map(woGamma1, woGamma2, f, v2, e)}\\

\textbf{Description:}
\begin{verbatim}
Returns mu_{wt1, wt2, phi}(v2, e)
        = \mu_{w,\phi}(v,e) = \sum_{u\in N_{\Gamma_2}(v), \phi(u)=y} w_2(u,v),
        if e=(\phi(u), \phi(v)), divided by the wt1(e) (wt1 the edge weights of woGamma1).

For the sum, first fix a vertex v2 on Gamma2 and an edge s of Gamma1. 
If s is not incident to phi(v) then discard it and pick a new one that is.
    
Let Nb_{v2, s} denote the set of neighbors u of v2 in Gamma2 such that
    (1) phi(v2) != phi(u), 
    (2) s = (phi(u), phi(v2)).
The sum in  
          mu_{wt1, wt2, phi}(v2, e) 
runs over all such neighbors u. It may seem that (u,v2) runs over the preimages of e, 
but really (u,v2) only runs over those preimages that are incident to v2 (which has been 
fixed once and for all).

Here:
    - woGamma1 is a triple: [Gamma1 = (V1, E1), epsilon1, wt1], where 
    - epsilon1 is a pair [E1,e1], where e1 is a list of length |E1| of edge orientations,
    - wt1 is a pair [E1,w1], where w1 is a list of length |E1| of edge weights > 0.
    - woGamma2 is similar.
    - f : defines a graph-theoretic mapping
      from Gamma2 to Gamma1 represented by a pair of lists [L1, L2],
      where L1 is the list of all n2 vertices of Gamma2,
      and L2 is the list of length n2 of the vertices in Gamma1: L2 forms 
      the corresponding image of L1 under the map f)
    - v2 : a vertex of Gamma2
    -  e : an edge of Gamma1 incident to phi(v2)
\end{verbatim}

\textbf{Example:}
\begin{verbatim}
    EXAMPLES:
        sage: V2 = [0,1,2,3]
        sage: E2 = [(0, 1), (1, 2), (0, 2), (0, 3)]
        sage: Gamma2 = Graph([V2,E2])
        sage: eps2 = [1,1,1,1]
        sage: wt2 = [E2, [1/2, 2/3, 3/5, 4/3]]
        sage: woGamma2 = [Gamma2, eps2, wt2]
        sage: V1 = [0,1,2]
        sage: E1 = [(0, 1), (1, 2), (0, 2)]
        sage: Gamma1 = Graph([V1, E1])
        sage: eps1 = [1,1,1]
        sage: wt1 = [E1, [1/2, 2/3, 3/5]]
        sage: woGamma1 = [Gamma1, eps1, wt1]
        sage: f = [[0,1,2,3], [0, 1, 2, 0]]
        sage: e = (1, 2)
        sage: [weighted_local_horizontal_multiplicity_of_harmonic_map(woGamma1, woGamma2, f, v, e) for v in V2]
         [0, 1, 1, 0]
\end{verbatim}

\section{\tt \ \ weighted\_vertical\_multiplicity\_of\_harmonic\_map}
\textbf{Syntax:} \verb}weighted_vertical_multiplicity_of_harmonic_map(woGamma1, woGamma2, f, v2)}\\

\textbf{Description:}
\begin{verbatim}
returns the weighted vertical multiplicity of f at v2
\end{verbatim}

\textbf{Example:}
\begin{verbatim}
    EXAMPLES:
        sage: V2 = [0,1,2,3]
        sage: E2 = [(0, 1), (1, 2), (0, 2), (0, 3)]
        sage: Gamma2 = Graph([V2,E2])
        sage: eps2 = [1,1,1,1]
        sage: wt2 = [E2, [1/2, 2/3, 3/5, 4/3]]
        sage: woGamma2 = [Gamma2, eps2, wt2]
        sage: V1 = [0,1,2]
        sage: E1 = [(0, 1), (1, 2), (0, 2)]
        sage: Gamma1 = Graph([V1, E1])
        sage: eps1 = [1,1,1]
        sage: wt1 = [E1, [1/2, 2/3, 3/5]]
        sage: woGamma1 = [Gamma1, eps1, wt1]
        sage: f = [[0,1,2,3], [0, 1, 2, 0]]
        sage: weighted_vertical_multiplicity_of_harmonic_map(woGamma1, woGamma2, f, 0)
         4/3
\end{verbatim}

\section{\tt \ \ weighted\_horizontal\_multiplicity\_of\_harmonic\_map}

\textbf{Syntax:} \verb}weighted_horizontal_multiplicity_of_harmonic_map(woGamma1, woGamma2, f, v2, verbose=False)}\\

\textbf{Description:}

Returns the pair (True, mu\_{wt2, phi}(v2)) if  mu\_{wt2, phi}(v2, e)
    is independent of e incident to phi(v2), and the pair
    (False, [mu\_{wt2, phi}(v2, e) for e incident to phi(v2)])
    otherwise.

    woGamma1 is a triple: [Gamma1 = (V1, E1), epsilon1, wt1], where 
     epsilon1 is a pair [E1,e1], where e1 is a list of length |E1| of edge orientations,
     wt1 is a pair [E1,w1], where w1 is a list of length |E1| of edge weights > 0.
    woGamma2 is similar.
    f : defines a graph-theoretic mapping
     from Gamma2 to Gamma1 represented by a pair of lists [L1, L2],
     where L1 is the list of all n2 vertices of Gamma2,
     and L2 is the list of length n2 of the vertices in Gamma1: L2 forms 
     the corresponding image of L1 under the map f)
    v2 : a vertex of Gamma2
    e : an edge of Gamma1 incident to phi(v2)

\textbf{Example:}
\begin{verbatim}
EXAMPLES:
        sage: V2 = [0,1,2,3]
        sage: E2 = [(0, 1), (1, 2), (0, 2), (0, 3)]
        sage: Gamma2 = Graph([V2,E2])
        sage: eps2 = [1,1,1,1]
        sage: wt2 = [E2, [1/2, 2/3, 3/5, 4/3]]
        sage: woGamma2 = [Gamma2, eps2, wt2]
        sage: V1 = [0,1,2]
        sage: E1 = [(0, 1), (1, 2), (0, 2)]
        sage: Gamma1 = Graph([V1, E1])
        sage: eps1 = [1,1,1]
        sage: wt1 = [E1, [1/2, 2/3, 3/5]]
        sage: woGamma1 = [Gamma1, eps1, wt1]
        sage: f = [[0,1,2,3], [0, 1, 2, 0]]
        sage: weighted_horizontal_multiplicity_of_harmonic_map(woGamma1, woGamma2, f, 0)
         (True, 1)
        sage: weighted_horizontal_multiplicity_of_harmonic_map(woGamma1, woGamma2, f, 1)
         (True, 1)
        sage: weighted_horizontal_multiplicity_of_harmonic_map(woGamma1, woGamma2, f, 2)
         (True, 1)
        sage: weighted_horizontal_multiplicity_of_harmonic_map(woGamma1, woGamma2, f, 3)
         (True, 0)
        sage: weighted_horizontal_multiplicity_of_harmonic_map(woGamma1, woGamma2, f, 0, verbose=True)
         The image of  0  under phi is  0 .   The (mu(phi, v, s), v, s), s incident to phi(v):   [(1, 0, (0, 1)), (1, 0, (0, 2))] .
         (True, 1)
\end{verbatim}

\section{\tt \ \ is\_weighted\_harmonic\_graph\_morphism}

\textbf{Syntax:} \verb}is_weighted_harmonic_graph_morphism(woGamma1, woGamma2, f, verbose = False)}\\

\textbf{Description:}
\begin{verbatim}

Returns true if f defines a weighted graph-theoretic mapping
    with a "constant" local horizontal multiplicity and False otherwise. 

    If verbose=True then more information will be printed out when the
    map f is not a harmonic morphism.

    woGamma1 is a triple: [Gamma1 = (V1, E1), epsilon1, wt1], where 
     epsilon1 is a pair [E1,e1], where e1 is a list of length |E1| of edge orientations,
     wt1 is a pair [E1,w1], where w1 is a list of length |E1| of edge weights >0.
    woGamma2 is similar.

    Recall, f is given as a map from an edge-weighted, oriented graph Gamma2 to 
    an edge-weighted, oriented graph Gamma1.
    Suppose Gamma2 has n vertices. A morphism 
              f: Gamma2 -> Gamma1
    is represented by a pair of lists [L2, L1],
    where L2=V2 is the list of all n vertices of Gamma2,
    and L1 is the list of length n of the vertices
    in Gamma1 that form the corresponding image under
    the map f.

    Slightly different output format than
      is\_harmonic\_graph\_morphism(Gamma1, Gamma2, f, verbose = False)

\end{verbatim}

\textbf{Example:}
\begin{verbatim}

    EXAMPLES:
        sage: V2 = [0,1,2,3]
        sage: E2 = [(0, 1), (1, 2), (0, 2), (0, 3)]
        sage: Gamma2 = Graph([V2,E2])
        sage: eps2 = [E2,[1,1,1,1]]
        sage: wt2 = [E2, [1/2, 2/3, 3/5, 4/3]]
        sage: woGamma2 = [Gamma2, eps2, wt2]
        sage: V1 = [0,1,2]
        sage: E1 = [(0, 1), (1, 2), (0, 2)]
        sage: Gamma1 = Graph([V1, E1])
        sage: eps1 = [E1, [1,1,1]]
        sage: wt1 = [E1, [1/2, 2/3, 3/5]]
        sage: woGamma1 = [Gamma1, eps1, wt1]
        sage: f = [[0,1,2,3], [0, 1, 2, 0]]       
        sage: is_weighted_harmonic_graph_morphism(woGamma1, woGamma2, f, verbose = False)
         False
        sage: is_weighted_harmonic_graph_morphism(woGamma1, woGamma2, f, verbose = True)
         morphism phi has a non-constant local horizontal mult at 1
         False
    
        sage: V2 = [0,1,2,3]
        sage: E2 = [(0, 1), (0, 2), (0, 3), (1, 2), (2, 3)]
        sage: Gamma2 = Graph([V2,E2])
        sage: eps2 = [E2, [1,1,1,1,1]]
        sage: wt2 = [E2, [1/2, 2/3, 1/2, 1/3, 2/3]]       
        sage: woGamma2 = [Gamma2, eps2, wt2]
        sage: V1 = [0,1,2]
        sage: E1 = [(0, 1), (1, 2)]
        sage: Gamma1 = Graph([V1,E1])
        sage: eps1 = [E1, [1,1]]
        sage: wt1 = [E1, [1/2, 1/3]]     
        sage: woGamma1 = [Gamma1, eps1, wt1]
        sage: f = [[0,1,2,3], [1, 0, 1, 2]]       
        sage: is_weighted_harmonic_graph_morphism(woGamma1, woGamma2, f, verbose = True)
          morphism phi has a non-constant local horizontal mult at 0
          False
        sage: weighted_horizontal_multiplicity_of_harmonic_map(woGamma1, woGamma2, f, 0)
         (False, [1, 3/2])
        sage: weighted_horizontal_multiplicity_of_harmonic_map(woGamma1, woGamma2, f, 1)
         (True, 5/3)
        sage: weighted_horizontal_multiplicity_of_harmonic_map(woGamma1, woGamma2, f, 2)
         (False, [2/3, 2])
        sage: weighted_horizontal_multiplicity_of_harmonic_map(woGamma1, woGamma2, f, 3)
         (True, 7/2)

        sage: V2 = [0,1,2,3]
        sage: E2 = [(0, 1), (0, 2), (0, 3), (1, 2), (2, 3)]
        sage: Gamma2 = Graph([V2,E2])
        sage: eps2 = [E2, [1,1,1,1,1]]
        sage: wt2 = [E2, [1/2, 2/3, 1/2, 1/3, 1/3]]  ## slightly different than above
        sage: woGamma2 = [Gamma2, eps2, wt2]
        sage: V1 = [0,1,2]
        sage: E1 = [(0, 1), (1, 2)]
        sage: Gamma1 = Graph([V1,E1])
        sage: eps1 = [E1, [1,1]]
        sage: wt1 = [E1, [1/3, 1/3]]     ## slightly different than above
        sage: woGamma1 = [Gamma1, eps1, wt1]
        sage: f = [[0,1,2,3], [1, 0, 1, 2]]       
        sage: is_weighted_harmonic_graph_morphism(woGamma1, woGamma2, f, verbose = False)
         True
        sage: weighted_horizontal_multiplicity_of_harmonic_map(woGamma1, woGamma2, f, 0)   
         (True, 3/2)
        sage: weighted_horizontal_multiplicity_of_harmonic_map(woGamma1, woGamma2, f, 1)   
         (True, 5/2)
        sage: weighted_horizontal_multiplicity_of_harmonic_map(woGamma1, woGamma2, f, 2)    
         (True, 1)
        sage: weighted_horizontal_multiplicity_of_harmonic_map(woGamma1, woGamma2, f, 3)  
         (True, 5/2)
\end{verbatim}

\section{\tt \ \ is\_harmonic\_graph\_morphism}

\textbf{Syntax:} \verb}is_harmonic_graph_morphism(Gamma1, Gamma2, f, verbose = False)}\\

\textbf{Description:}
\begin{verbatim}
   Returns True if f defines a graph-theoretic mapping
    from Gamma2 to Gamma1 that is harmonic, and False otherwise. 

    Assumes Gamma1, Gamma2 are in standard form.

    Suppose Gamma2 has n vertices. A morphism 
              f: Gamma2 -> Gamma1
    is represented by a pair of lists [L2, L1],
    where L2 is the list of all n vertices of Gamma2,
    and L1 is the list of length n of the vertices
    in Gamma1 that form the corresponding image under
    the map f.
\end{verbatim}

\textbf{Example:}
\begin{verbatim}
    EXAMPLES:
        sage: Gamma2 = graphs.TetrahedralGraph()
        sage: Gamma1 = graphs.CycleGraph(3)
        sage: f = [[0,1,2,3],[0,1,2,0]]
        sage: is_harmonic_graph_morphism(Gamma1, Gamma2, f)
        False
        sage: is_harmonic_graph_morphism(Gamma1, Gamma2, f, verbose=True)
        This [<vertex>, <intersection sizes>]] passes the check: [0, [1, 1]]
        This [<vertex>, <intersection sizes>]] fails the check: [1, [2, 1]]
        This [<vertex>, <intersection sizes>]] fails the check: [2, [2, 1]]
        False
        sage: V1 = [0,1,2]
        sage: V2 = [0,1,2,3]
        sage: E2 = [(0,1),(1,2),(2,0),(2,3)]
        sage: E1 = [(0,1),(1,2),(2,0)]
        sage: Gamma2 = Graph([V2,E2],format='vertices_and_edges')
        sage: Gamma1 = Graph([V1,E1],format='vertices_and_edges')
        sage: f = [V2, [0,1,2,2]]    # thi is a retract
        sage: is_harmonic_graph_morphism(Gamma1, Gamma2, f, verbose=True)
        This [<vertex>, <intersection sizes>]] passes the check: [0, [1, 1]]
        This [<vertex>, <intersection sizes>]] passes the check: [1, [1, 1]]
        This [<vertex>, <intersection sizes>]] passes the check: [2, [1, 1]]
        This [<vertex>, <intersection sizes>]] passes the check: [3, [0, 0]]
        True
        sage: Gamma1 = graphs.TetrahedralGraph() 
        sage: Gamma2 = graphs.CubeGraph(3)
        sage: is_harmonic_graph_morphism(Gamma1,Gamma2,f)
         'Error! -- Gamma2 must be in standard form.'
\end{verbatim}

\section{\tt \ \ find\_all\_harmonic\_graph\_morphism}

\textbf{Syntax:} \verb}find_all_harmonic_graph_morphism(Gamma1, Gamma2, out_filename = "")}\\

\textbf{Description:}
\begin{verbatim}
Runs over all possible f:V1 -> V2 and checks if
is_harmonic_graph_morphism(Gamma1, Gamma2, f) is true.
If so, writes f to filename (unless out_filename = "", in
which case it prints it to file).
\end{verbatim}

Potentially very long output, possibly several Mb in size.

\section{\tt \ \ vertical\_multiplicity\_of\_graph\_morphism}

\textbf{Syntax:} \verb}vertical_multiplicity_of_graph_morphism(Gamma1, Gamma2, f, v2)}\\

\textbf{Description:}
\begin{verbatim}
INPUT:
    - Gamma1 : a graph, 
    - Gamma2 : a graph, 
    - f : defines a graph morphism
      from Gamma2 to Gamma1 
      (represented by a pair of lists [L1, L2],
      where L1 is the list of all n2 vertices of Gamma2,
      and L2 is the list of length n2 of the vertices
      in Gamma1: L2 forms the corresponding image of L1 under the map f)
    - v2 : a vertex of Gamma2.

OUTPUT:
    Returns nu_phi(v2) - the number of edges of Gamma2 incident to v2 that
    have the same (vertex) image as v2 under f.
\end{verbatim}

\textbf{Example:}
\begin{verbatim}
    EXAMPLES:
        sage: V2 = [0,1,2,3]
        sage: E2 = [(0, 1), (1, 2), (1, 3)]
        sage: Gamma2 = Graph([V2,E2])
        sage: Gamma1 = graphs.PathGraph(3)
        sage: f = [[0,1,2,3],[2, 1, 1, 0]]
        sage: [vertical_multiplicity_of_graph_morphism(Gamma1, Gamma2, f, v) for v in V2]
        [0, 1, 1, 0]
        sage: Gamma2 = graphs.CycleGraph(6)
        sage: Gamma1 = graphs.PathGraph(3)
        sage: f = [[0,1,2,3,4,5],[2, 2, 1, 0, 0, 1]]
        sage: [vertical_multiplicity_of_graph_morphism(Gamma1, Gamma2, f, v) 
                                                     for v in Gamma2.vertices(sort=True)]
        [1, 1, 0, 1, 1, 0]
\end{verbatim}

\section{\tt \ \ vertical\_multiplicity\_of\_harmonic\_map}

\textbf{Syntax:} \verb}vertical_multiplicity_of_harmonic_map(Gamma1, Gamma2, f, v2)}\\

If f is not harmonic, this function works anyway.

\textbf{Description:}
\begin{verbatim}
INPUT:
    - Gamma1 : a graph, 
    - Gamma2 : a graph, 
    - f : defines a graph-theoretic mapping
      from Gamma2 to Gamma1 that (presumably) is harmonic
      (represented by a pair of lists [L1, L2],
      where L1 is the list of all n2 vertices of Gamma2,
      and L2 is the list of length n2 of the vertices
      in Gamma1: L2 forms the corresponding image of L1 under the map f)
    - v2 : a vertex of Gamma2.
    
OUTPUT:
    Returns nu_phi(v1)
\end{verbatim}

\textbf{Example:}
\begin{verbatim}
    EXAMPLES:
        sage: V2 = [0,1,2,3]
        sage: E2 = [(0, 1), (1, 2), (1, 3)]
        sage: Gamma2 = Graph([V2,E2])
        sage: Gamma1 = graphs.PathGraph(3)
        sage: f = [[0,1,2,3],[2, 1, 1, 0]]
        sage: [vertical_multiplicity_of_harmonic_map(Gamma1, Gamma2, f, v) for v in V1]
        [0, 1, 1, 0]
        sage: Gamma2 = graphs.CycleGraph(6)
        sage: Gamma1 = graphs.PathGraph(3)
        sage: f = [[0,1,2,3,4,5],[2, 2, 1, 0, 0, 1]]
        sage: [vertical_multiplicity_of_harmonic_map(Gamma1, Gamma2, f, v) for v in Gamma2.vertices(sort=True)]
        [1, 1, 0, 1, 1, 0]
\end{verbatim}

\section{\tt \ \ horizontal\_multiplicity\_of\_harmonic\_map}

\textbf{Syntax:} \verb}horizontal_multiplicity_of_harmonic_map(Gamma1, Gamma2, f, v2)}\\

\textbf{Description:}
\begin{verbatim}
INPUT:
    - Gamma1 : a graph, 
    - Gamma2 : a graph, 
    - f : defines a graph-theoretic mapping
      from Gamma2 to Gamma1 that (presumably) is harmonic
      (represented by a pair of lists [L1, L2],
      where L1 is the list of all n2 vertices of Gamma2,
      and L2 is the list of length n2 of the vertices
      in Gamma1: L2 forms the corresponding image of L1 under the map f)
    - v2 : a vertex of Gamma2.
    
OUTPUT:
    Returns mu_phi(v1)
\end{verbatim}

\textbf{Example:}
\begin{verbatim}
    EXAMPLES:
        sage: V2 = [0,1,2,3]
        sage: E2 = [(0, 1), (1, 2), (1, 3)]
        sage: Gamma2 = Graph([V2,E2])
        sage: Gamma1 = graphs.PathGraph(3)
        sage: f = [[0,1,2,3],[2, 1, 1, 0]]
        sage: [horizonal_multiplicity_of_harmonic_map(Gamma1, Gamma2, f, v) for v in V2]
\end{verbatim}

\section{\tt \ \ horizonal\_multiplicity\_of\_harmonic\_map}

\textbf{Syntax:} \verb}horizonal_multiplicity_of_harmonic_map(Gamma1, Gamma2, f, v2)}\\

I mis-spell horizontal all the time:-)

\section{\tt \ \ genus}

\textbf{Syntax:} \verb}genus(Gamma)}\\

\textbf{Description:}
\begin{verbatim}
If Gamma = (V, E) then this computes |E|-|V|+1. 
Not the same as Gamma.genus()
\end{verbatim}

\textbf{Example:}
\begin{verbatim}
    EXAMPLES:
        sage: Gamma2 = graphs.WheelGraph(5)
        sage: Gamma2.genus()
        0
        sage: genus(Gamma2)
        4
\end{verbatim}

\section{\tt \ \ perm\_number\_action}

\textbf{Syntax:} \verb}perm_number_action(G, g0)}\\

This is a function for permutation groups, used in some of the Cayley graph functions
in this module.

\textbf{Description:}
\begin{verbatim}
INPUT:
  G - a permutation group
 g0 - an element of G

OUTPUT:
Returns list of pairs, indexed by the number of elements of the group G, 
where first number represents the group element and the second is where that 
group element is sent by right multiplication by g0.
\end{verbatim}

\textbf{Example:}
\begin{verbatim}
        sage: G = SymmetricGroup(3)
        sage: g0 = G.random_element(); g0
        (1,2,3)
        sage: perm_number_action(G, g0)
        [(0, 2), (1, 0), (2, 1), (3, 5), (4, 3), (5, 4)]
\end{verbatim}

\section{\tt \ \ lifted\_vertex\_to\_cover}

\textbf{Syntax:} \verb}lifted_vertex_to_cover(G, Gamma, v, g0)}\\

The $G$-cover of $\Gamma$ builds a copy of your graph for every group element in $G$.

Each vertex $v$ in the original graph becomes multiple vertices in the cover - one for each group element.

This function tells you ``Which of those copies of $v$ corresponds to a specific group label g0?''

It returns the index in the covering graph’s vertex list for the pair 
$(g_0, v)$.

\textbf{Description:}
\begin{verbatim}
In a (to be constructed) G-covering of Gamma, call it Gamma' = (V', E'), 
this function returns the vertex in V' that v lifts to associated to
g0 under the inverse of the projection map phi: Gamma' -> Gamma.

    INPUT:
        G - a non-trivial subgroup of a symmetric group,
        Gamma - a graph with vertices V and edges E
        v - a vertex in V
        g0 - a non-trivial element of G

\end{verbatim}

\textbf{Example:}
\begin{verbatim}
    EXAMPLES:
        sage: Gamma = graphs.PathGraph(3)
        sage: V = Gamma.vertices(sort=True); V
         [0, 1, 2]
        sage: G = PermutationGroup([[(0,1)]])  # order 2 group: identity and (0 1)
        sage: G.list()
         [(), (0,1)]
        sage: g0 = G.list()[1]  # the nontrivial permutation (0 1)
        sage: lifted_vertex_to_cover(G, Gamma, 0, g0)
         1
        sage: lifted_vertex_to_cover(G, Gamma, 1, g0)
         0
        sage: lifted_vertex_to_cover(G, Gamma, 2, g0)
         2

    EXPLANATION:

        - The group G = {(), (0 1)} acts on the vertex set [0, 1, 2].
        - The cover has 2 x 3 = 6 vertices (one copy of [0,1,2] for each group element).
        - The lifted vertex (g0, v) is assigned an index in range(6) corresponding to 
          the action of g0 on v, computed via perm_number_action.
\end{verbatim}

\section{\tt \ \ lifted\_vertices\_to\_cover}

\textbf{Syntax:} \verb}lifted_vertices_to_cover(G, Gamma, v, g0 = 1)}\\

This returns the fiber, not just a single pre-image.

\textbf{Description:}
\begin{verbatim}
Returns the fiber over a vertex v in the G-covering of Gamma.

INPUT:
   - G is a finite group of Symm_n acting on V (n = |V|),
   - Gamma is a graph with vertex set V in standard form.
   - v, a vertex in V
   - g0, optional group element (defaults to identity)

OUTPUT:
    A list of length |G| giving the indices of the lifted vertices
     in the flat vertex list of the cover graph.

The associated G-covering graph has vertex set 
                   V' = G x V. 
This function returns the list of vertices in V' that map down to v under the projection 
               phi: Gamma' -> Gamma. 
These are the "lifted vertices" of v.

If a group element g0 != 1 is provided, the fiber is rotated by applying g0 to 
the index of v before lifting.

\end{verbatim}

\textbf{Example:}
\begin{verbatim}
    EXAMPLES:
        sage: Gamma = graphs.PathGraph(3)
        sage: V = Gamma.vertices(sort=True); V
        [0, 1, 2]
        sage: G = PermutationGroup([[(0,1)]])  # group of order 2
        sage: G.list()
        [(), (0,1)]
        sage: lifted_vertices_to_cover(G, Gamma, 0)
        [0, 3]
        sage: lifted_vertices_to_cover(G, Gamma, 1)
        [1, 4]
        sage: lifted_vertices_to_cover(G, Gamma, 2)
        [2, 5]
        sage: g0 = G.list()[1]  # (0 1)
        sage: lifted_vertices_to_cover(G, Gamma, 0, g0)
        [1, 4]

    The total cover has 2 x 3 = 6 vertices. Each call returns
    the indices of the 2 vertices in the fiber over v.
\end{verbatim}

\section{\tt \ \ lifted\_edges\_to\_cover}

\textbf{Syntax:} \verb}lifted_edges_to_cover(G, Gamma, e)}\\

\textbf{Description:}
\begin{verbatim}
In a (to be constructed) G-covering of Gamma, call it Gamma' = (V', E'), 
this function returns the edges in E' that e lifts to associated to
the inverse of the projection map
                  phi: Gamma' -> Gamma.

    INPUT:
        G - a non-trivial subgroup of a symmetric group,
        Gamma - a graph with vertices V and edges E
        e - a vertex in E
\end{verbatim}

\textbf{Example:}
\begin{verbatim}
    EXAMPLES:
        sage: Gamma = graphs.PathGraph(3)
        sage: G = PermutationGroup([[(0,1)]])
        sage: e = (0, 1)
        sage: lifted_edges_to_cover(G, Gamma, e)
        [[0, 4], [3, 1]]

    The base graph Gamma has 3 vertices.
    The cover has 2 x 3 = 6 vertices:
        ((), 0), ((), 1), ((), 2),
        ((0,1), 0), ((0,1), 1), ((0,1), 2)
    The edge (0, 1) in Gamma lifts to two edges
    connecting lifted vertices: [0, 4] and [3, 1].

    EXAMPLES: ## (continued)
        sage: Gamma = graphs.PathGraph(4)
        sage: G = PermutationGroup([[(0,1,2)]])  # a 3-cycle
        sage: lifted_edges_to_cover_full(G, Gamma, (1, 2))
        [[1, 8], [4, 2], [7, 5]]

\end{verbatim}

\section{\tt \ \ lifted\_graph\_to\_cover}

\textbf{Syntax:} \verb}lifted_graph_to_cover0(G, Gamma)}\\

\textbf{Description:}
\begin{verbatim}
    Constructs a G-covering of Gamma.

    INPUT:
        G - a non-trivial subgroup of a symmetric group,
        Gamma - a graph with vertices V and edges E

\end{verbatim}

\textbf{Example:}
\begin{verbatim}
    EXAMPLES:
        sage: Gamma = graphs.CycleGraph(3)
        sage: S3 = SymmetricGroup(3)
        sage: Gamma2 = lifted_graph_to_cover(S3, Gamma)
        sage: Gamma2.show()
        sage: Gamma2.automorphism_group().cardinality()
        768
        sage: quotient_graph(Gamma2, AG.sylow_subgroup(3))
        Graph on 6 vertices
\end{verbatim}


\section{\tt \ \ cayley\_expander\_graph\_from\_gen\_mat}

\textbf{Syntax:} \verb}cayley_expander_graph_from_gen_mat(G)}\\

This function takes a kxn generator matrix G of such a code C = RowSpan(G)
and uses it to build a Cayley graph, Gamma = Cay(G).


Start with a vector space C (the linear code associated to the generator matrix)
over GF(2) — this gives the vertices, V.

Connect two vertices u and v if their difference (or sum, since it’s over GF(2))
is a row of the generator matrix G - this gives the edges E.


\textbf{Description:}
\begin{verbatim}
INPUT:
  G - a rank k kxn matrix over GF(2).

Assumes k < n  

OUTPUT:
  Cay(G) = (V, E)
\end{verbatim}

\textbf{Example:}
\begin{verbatim}
    EXAMPLES:
        sage: C = codes.HammingCode(GF(2), 3); C
        [7, 4] Hamming Code over GF(2)
        sage: G = C.generator_matrix()
        sage: cayley_expander_graph_from_gen_mat(G)
        Graph on 16 vertices
\end{verbatim}

\section{\tt \ \ cayley\_expander\_graph\_from\_code}

\textbf{Syntax:} \verb}cayley_expander_graph_from_code(C)}\\

Builds a Cayley expander graph from a binary linear code C.

Let G be the generator matrix returned by SageMath's 
    {\tt generator\_matrix} method for the LinearCode class C belongs to.

    The vertices V of the graph Cay(C) correspond to all vectors in C. 
    Two vertices u and v are connected if u + v is a row of G.
    
\textbf{Description:}
\begin{verbatim}
    INPUT:
        C -- a linear code over GF(2)

    OUTPUT:
        A undirected graph Gamma with 2^k vertices
\end{verbatim}

\textbf{Example:}
\begin{verbatim}
    EXAMPLES:
        sage: C = codes.HammingCode(GF(2), 3); C
        [7, 4] Hamming Code over GF(2)
        sage: G = C.generator_matrix()
        sage: C1 = cayley_expander_graph_from_gen_mat(G); C1
        Graph on 16 vertices
        sage: C2 = cayley_expander_graph_from_code(C); C2
        Graph on 16 vertices
        sage: C1 == C2
        True
        sage: C1 == C
        False
        sage: C2 == C
        False
\end{verbatim}

\section{\tt \ \ edge\_expansion}

\textbf{Syntax:} \verb}edge_expansion(Gamma)}\\

\textbf{Description:}
\begin{verbatim}
Returns the edge expansion of Gamma.
The edge expansion (also isoperimetric number or Cheeger constant) 
         h(G) 
of a graph is defined in [0]. Basically, edge expansion measures 
how hard it is to isolate a small part of the graph. If this number 
is big, Gamma is hard to cut apart. 

    REFERENCES:
        [0] https://en.wikipedia.org/wiki/Expander_graph#Edge_expansion
\end{verbatim}

\textbf{Example:}
\begin{verbatim}
    EXAMPLE:
        sage: G = graphs.CycleGraph(6)
        sage: edge_expansion(G)
        1
\end{verbatim}

\section{\tt \ \ laplacian\_pseudoinverse}

\textbf{Syntax:} \verb}laplacian_pseudoinverse(Gamma)}\\

This function simply implements a known formula for the Moore–Penrose
pseudoinverse of the Laplacian, $Q^+$.

\textbf{Description:}
\begin{verbatim}
Returns Q^+
\end{verbatim}

\textbf{Example:}
\begin{verbatim}
    EXAMPLES:
        sage: laplacian_pseudoinverse(graphs.CycleGraph(4))
         [ 5/16 -1/16 -3/16 -1/16]
         [-1/16  5/16 -1/16 -3/16]
         [-3/16 -1/16  5/16 -1/16]
         [-1/16 -3/16 -1/16  5/16]
        sage: Gamma = graphs.CycleGraph(5); Gamma
        Cycle graph: Graph on 5 vertices
        sage: laplacian_pseudoinverse(Gamma)
         [ 2/5    0 -1/5 -1/5    0]
         [   0  2/5    0 -1/5 -1/5]
         [-1/5    0  2/5    0 -1/5]
         [-1/5 -1/5    0  2/5    0]
         [   0 -1/5 -1/5    0  2/5]
\end{verbatim}

\section{\tt \ \ hamiltonian\_paths0}

\textbf{Syntax:} \verb}hamiltonian_paths0(Gamma)}\\

\textbf{Description:}
\begin{verbatim}
This function collects all spanning trees and filters those with
max degree <= 2.
\end{verbatim}

\textbf{Example:}
\begin{verbatim}
    EXAMPLES:
        sage: Gamma = graphs.GridGraph([3,3])
        sage: HP = hamiltonian_paths(Gamma)
        sage: len(HP)
        20
\end{verbatim}

\section{\tt \ \ hamiltonian\_paths}

\textbf{Syntax:} \verb}hamiltonian_paths(Gamma, signed_adjacency_matrix = [])}\\

\textbf{Description:}
\begin{verbatim}
When a signed adjacency matrix is passed into hamiltonian\_paths
returns a list of triples:
                [Graph, vertex_ordering, signature]
the vertex_ordering describing a Hamiltonian path in Gamma.
\end{verbatim}

\textbf{Example:}
\begin{verbatim}
    EXAMPLES:
        sage: Gamma = graphs.GridGraph([3,3])
        sage: HP = hamiltonian_paths(Gamma)
        sage: len(HP)
        20
        sage: A = matrix(QQ,[
        [0 , -1 , 1  , -1 , -1 , -1 ],
        [1,   0 ,  -1,  1,  1,   -1  ],
        [-1 , 1 ,  0 ,  1 , 1  , -1  ],
        [1 , -1 , -1,  0 ,  -1 , -1  ],
        [1 , - 1 , - 1 , 1 , 0 , - 1  ],
        [1 ,  1  ,  1  , 1  , 1  , 0 ]
        ])
        sage: Gamma = Graph(A, format='weighted_adjacency_matrix')
        sage: HP = hamiltonian_paths(Gamma, signed_adjacency_matrix = A)
        sage: L = [sum(x[2]) for x in HP]; max(L)
        5
        sage: L.index(5)
        21
        sage: HP[21]                                 
        [Graph on 6 vertices,
         [0, 5, 2, 1, 3, 4],
         [-1, 1, -1, -1, -1, 1, 1, 1, 1, 1, 1, 1, 1, 1, -1]]
        sage: L.count(5)
        1
\end{verbatim}

\section{\tt \ \ random\_walk}

\textbf{Syntax:} \verb}random_walk(Gamma, m)}\\

\textbf{Description:}
\begin{verbatim}
Returns a random walk of length n in Gamma.
\end{verbatim}

\textbf{Example:}
\begin{verbatim}
    EXAMPLES:
        sage: Gamma = graphs.HouseGraph()
        sage: random_walk(Gamma, 3)
        [1, 3, 2, 3]
\end{verbatim}

\section{\tt \ \ random\_trek}

\textbf{Syntax:} \verb}random_trek(Gamma, m)}\\

\textbf{Description:}
\begin{verbatim}
Returns a random trek of length n in Gamma.
\end{verbatim}

\textbf{Example:}
\begin{verbatim}
sage: Gamma = graphs.HouseGraph()
sage: random_trek(Gamma, 10)
[4, 2, 0, 1, 3, 2, 4, 3, 1, 0, 2]

\end{verbatim}

\section{\tt \ \ codifferential\_basis}

\textbf{Syntax:} \verb}codifferential_basis(v, e, Gamma)}\\

If the graph Gamma is regarded a network of water pipes (edges) and junctions (vertices),
this function tells you the effect on a particular junction v when only one edge e is
turned on (set to 1), and all other edges are set to 0.

\textbf{Description:}
\begin{verbatim}
INPUT: 
   - Gamma a graph
   - v a vertex of Gamma
   - e an edge of Gamma.

Returns the value of the codifferential
          delta:C^0(E) -> C^0(V)
on the characteristic functions of E, chi_e: E -> RR  
(the function on E that is 1 at e and 0 elsewhere), 
for each v in V.
\end{verbatim}

\textbf{Example:}
\begin{verbatim}
EXAMPLES:
        sage: Gamma = graphs.DiamondGraph()  
        sage: V = Gamma.vertices(sort=True)   
        sage: E = list(Gamma.edges(labels=None,sort=True)); E
        [(0, 1), (0, 2), (1, 2), (1, 3), (2, 3)]
        sage: delta_matrix = matrix([[codifferential_basis(v, e, Gamma) for e in E] for v in V])        
        sage: delta_matrix  
        [ 1/2  1/2    0    0    0]
        [-1/3    0  1/3  1/3    0]
        [   0 -1/3 -1/3    0  1/3]
        [   0    0    0 -1/2 -1/2]
        sage: delta_matrix.right_kernel()
        Vector space of degree 5 and dimension 2 over Rational Field
        Basis matrix:
        [ 1 -1  0  1 -1]
        [ 0  0  1 -1  1]
        sage: cycle_space(Gamma, [E, [1,1,1,1,1]])
        Vector space of degree 5 and dimension 2 over Rational Field
        Basis matrix:
        [ 1 -1  0  1 -1]
        [ 0  0  1 -1  1]
\end{verbatim}

\section{\tt \ \ codifferential\_matrix}

\textbf{Syntax:} \verb}codifferential_matrix(Gamma)}\\


\textbf{Description:}
\begin{verbatim}
Returns the |V|x|E| matrix of the codifferential
    delta:C^0(E) -> C^0(V)
that is, the matrix of values of codifferential_basis(v, e, Gamma).

\end{verbatim}

\textbf{Example:}
\begin{verbatim}
EXAMPLES:
        sage: Gamma = graphs.DiamondGraph() 
        sage: codifferential_matrix(Gamma) 
        [ 1/2  1/2    0    0    0]
        [-1/3    0  1/3  1/3    0]
        [   0 -1/3 -1/3    0  1/3]
        [   0    0    0 -1/2 -1/2]
\end{verbatim}

\section{\tt \ \ pushforward\_matrix\_vertex\_functions}

\textbf{Syntax:} \verb}pushforward_matrix_vertex_functions(Gamma1, Gamma2, f)}\\

\textbf{Description:}
\begin{verbatim}

If we assume there is a graph morphism
         phi:Gamma2 -> Gamma1,
denoted f in the usual color list notation, this function
find the matrix repn M of
        phi_*:C^0(Gamma2) -> C^0(Gamma1)
with respect to the standard characteristic function basis

We want a |V1|x|V2| matrix M for which e_i in RR^|V1| (0<=i<=|V1|)
is the sum of e_j in RR^|V2|, where phi(j)=i.
\end{verbatim}

\textbf{Example:}
\begin{verbatim}
EXAMPLES:
        sage: Gamma1 = graphs.PathGraph(3)
        sage: Gamma2 = graphs.DiamondGraph()
        sage: V2 = Gamma2.vertices(sort=True); V2
        [0, 1, 2, 3]
        sage: f = [V2, [0, 1, 1, 2]]
        sage: P = pushforward_matrix_vertex_functions(Gamma1, Gamma2, f); P
        [1 0 0 0]
        [0 1 1 0]
        [0 0 0 1]
        sage: P.rank() 
        3
        sage: matrix_of_graph_morphism_vertices(Gamma1, Gamma2, f)
        [1 0 0]
        [0 1 0]
        [0 1 0]
        [0 0 1]
\end{verbatim}

\section{\tt \ \ pushforward\_matrix\_edge\_functions}

\textbf{Syntax:} \verb}pushforward_matrix_edge_functions(Gamma1, Gamma2, f)}\\

\textbf{Description:}
\begin{verbatim}
If we assume there is a graph morphism
    phi:Gamma2 -> Gamma1, denoted f in the usual color list notation,
find the matrix repn of
    phi\_*:C^1(Gamma2) -> C^1(Gamma1)
with respect to the standard characteristic function basis

We want a |E1|x|E2| matrix P for which ei in RR^|E1| (0<=i<=|E1|)
is the sum of ej in RR^|E2|, where phi(E2[j])=E1[i].
\end{verbatim}

\textbf{Example:}
\begin{verbatim}
EXAMPLES:
        sage: Gamma1 = graphs.PathGraph(3)
        sage: Gamma2 = graphs.DiamondGraph()
        sage: V2 = Gamma2.vertices(sort=True); V2
        [0, 1, 2, 3]
        sage: f = [V2, [0, 1, 1, 2]]
        sage: P = pushforward_matrix_edge_functions(Gamma1, Gamma2, f); P
        [1 1 0 0 0]
        [0 0 0 1 1]
        sage: matrix_of_graph_morphism_edges_unsigned(Gamma1, Gamma2, f)
        [1 0]
        [1 0]
        [0 0]
        [0 1]
        [0 1]
\end{verbatim}

\section{\tt \ \ pullback\_matrix\_vertex\_functions}

\textbf{Syntax:} \verb}pullback_matrix_vertex_functions(Gamma1, Gamma2, f)}\\

\textbf{Description:}
\begin{verbatim}
If we assume there is a graph morphism
    phi:Gamma2 -> Gamma1, denoted f in the usual color list notation,
find the matrix repn of
    phi^*:C^0(Gamma1) -> C^0(Gamma2)
with respect to the standard characteristic function basis.

This is the transpose of pushforward\_matrix\_vertex\_functions.

We want a |V2|x|V1| matrix P for which ej in RR^|V1| (0<=j<=|V1|)
is the ei in RR^|V2|, where phi(i)=j.
\end{verbatim}

\textbf{Example:}
\begin{verbatim}
EXAMPLES:
        sage: Gamma1 = graphs.PathGraph(3)
        sage: Gamma2 = graphs.DiamondGraph()
        sage: V2 = Gamma2.vertices(sort=True); V2
        [0, 1, 2, 3]
        sage: f = [V2, [0, 1, 1, 2]]
        sage: P = pullback_matrix_vertex_functions(Gamma1, Gamma2, f); P
        [1 0 0]
        [0 1 0]
        [0 1 0]
        [0 0 1]
\end{verbatim}

\section{\tt \ \ proj\_blowup}

\textbf{Syntax:} \verb}proj_blowup(x, Gamma, r)}\\

\textbf{Description:}
\begin{verbatim}
This returns the projection map referenced in the function
    blowup\_graph(Gamma, r).

    Assumes 0 <= x <= nr-1. If V = [x\_0, ..., x\_{n-1}]
    then this function sends 
      0, ..., r-1 to x\_0 
      r, ..., 2r-1 to x\_1, 
    and so on.
\end{verbatim}

\textbf{Example:}
\begin{verbatim}
EXAMPLES:
        sage: Gamma = graphs.DiamondGraph()
        sage: [proj_blowup(x, Gamma, 2) for x in range(8)]
         [0, 0, 1, 1, 2, 2, 3, 3]
\end{verbatim}

\section{\tt \ \ blowup\_balanced\_graph}

\textbf{Syntax:} \verb}blowup_balanced_graph(Gamma, r)}\\

\textbf{Description:}
\begin{verbatim}
This function returns the r-blowup of Gamma.

Defn: The (balanced) $r$-blowup \Gamma^{(r)} of the graph $\Gamma$ is defined as follows: 
    *) For each vertex $x_i$ of of $\Gamma$, the graph \Gamma^{(r)}$ has $r$ vertices, 
       $ x_{i,1}, x_{i,2}, \dots , x_{i,r}$.
    *) There is an edge between vertices $x_{i,k}$ and $x_{j, \ell}$ in $\Gamma^{(r)}$
       if and only if there is an edge between vertices $x_i$ and $x_j$ in $\Gamma$.

Note: The projection map 
      $\phi : \Gamma^{(r)} \to \Gamma$ 
given on vertices by 
      $\phi(x_{i,j}) = x_i$ 
is a harmonic morphism (in fact, it's an r-multicover).    
\end{verbatim}

\textbf{Example:}
\begin{verbatim}
    EXAMPLES:
        sage: Gamma = graphs.DiamondGraph()  
        sage: Gamma.degree_sequence()        
         [3, 3, 2, 2]
        sage: Gammabu = blowup_balance_graph(Gamma, 2)
        sage: Gammabu.degree_sequence()      
         [6, 6, 6, 6, 4, 4, 4, 4]

        sage: Gamma = graphs.CycleGraph(4)   
        sage: Gammabu = blowup_balance_graph(Gamma, 2)
        sage: K44 = graphs.CompleteBipartiteGraph(4,4)
        sage: Gammabu.is_isomorphic(K44)      
         True
\end{verbatim}

\section{\tt \ \ blowup\_weighted\_graph}

\textbf{Syntax:} \verb}blowup_weighted_graph(wGamma, r)}\\

\textbf{Description:}
\begin{verbatim}
INPUT:
    - wGamma is a weighted graph given as a pair: 
    - wGamma = [(V, E), wt], where wt is a pair [E,w], 
      where w is a list of length |E| of edge weights >0.
OUTPUT:
    This function returns the weighted r-blowup of Gamma.

Defn: The (weighted) $r$-blowup $\Gamma^{(r)}$ of the graph $\Gamma$ is defined as follows: 
    *) For each vertex $x_i$ of of $\Gamma$, the graph $\Gamma^{(r)}$ has $r$ vertices, 
       $ x_{i,1}, x_{i,2}, \dots , x_{i,r}$.
    *) There is an edge between vertices $x_{i,k}$ and $x_{j, \ell}$ in $\Gamma^{(r)}$
       if and only if there is an edge between vertices $x_i$ and $x_j$ in $\Gamma$.
    *) Define $\pi^{(r)} : \Gamma^{(r)}\to \Gamma$ by
                  $\pi^{(r)}(x_{i,j}) = x_i$.
       If the weights of $\Gamma$ are defined by $w: E \to \rrr$ then 
       the weights of the $r$-blowup are defined by 
                  $w^{(r)}: E^{(r)} \to \rrr$
       where 
            $w^{(r)}(v_1,v_2) = w(\pi^{(r)}(v_1), \pi^{(r)}(v_2))$,
       if $(v_1, v_2)\in E$.

Note: The projection map 
      $\phi : \Gamma^{(r)} \to \Gamma$ 
given on vertices by 
      $\phi(x_{i,j}) = x_i$ 
is a weighted harmonic morphism. 
\end{verbatim}

\textbf{Example:}
\begin{verbatim}
    EXAMPLES:
        sage: V = [0,1,2,3]
        sage: E = [(0, 1), (1, 2), (0, 2), (0, 3)]
        sage: Gamma = Graph([V, E])
        sage: wt = [E, [1, 2, 3, 4]] 
        sage: wGamma = [Gamma, wt]
        sage: blowup_weighted_graph(wGamma, 2)
        [Graph on 8 vertices,
         [[(0, 2), (0, 3), (0, 4), (0, 5), (0, 6), (0, 7), (1, 2), (1, 3),
           (1, 4), (1, 5), (1, 6), (1, 7), (2, 4), (2, 5), (3, 4), (3, 5)],
          [1, 1, 3, 3, 4, 4, 1, 1, 3, 3, 4, 4, 2, 2, 2, 2]]] 
\end{verbatim}

\section{\tt \ \ find\_cover\_graph}

\textbf{Syntax:} \verb}find_cover_graph(Gamma, r, verbose=False)}\\

\textbf{Description:}
\begin{verbatim}
Returns a degree r covering graph
             pi: Gammar -> Gamma,
found by random search over possible edge sets, and the 
morphism f returned as an rn-tuple if entries from V,
where V = [0, ..., n-1].

Warning: Gammar can be disconnected.

Assumes that Gamma is connected and in standard form.
\end{verbatim}

\textbf{Example:}
\begin{verbatim}
EXAMPLES:
        sage: Gamma1 = graphs.CycleGraph(4)   
        sage: Gamma2, f0 = find_cover_graph(Gamma1, 3)  
        sage: Gamma2.is_connected()                   
         True
        sage: V2 = Gamma2.vertices(sort=True); V2               
         [0, 1, 2, 3, 4, 5, 6, 7, 8, 9, 10, 11]  
        sage: is_graph_morphism(Gamma1, Gamma2, f)   
         True
        sage: is_harmonic_graph_morphism(Gamma1, Gamma2, f) 
         False
        sage: Gamma2.degree_sequence() 
         [4, 3, 2, 2, 2, 2, 2, 2, 2, 1, 1, 1]
        sage: [vertical_multiplicity_of_harmonic_map(Gamma1, Gamma2, f, v) for v in V2] 
         [0, 0, 0, 0, 0, 0, 0, 0, 0, 0, 0, 0]
\end{verbatim}

\section{\tt \ \ find\_random\_incidence\_matrix}

\textbf{Syntax:} \verb}find_random_incidence_matrix(m,n)}\\

\textbf{Description:}
\begin{verbatim}
Returns either
    an mxn (0,1)-matrix M such that
    (a) each column has exactly two $1$'s, 
    (b) no row is all $0$s, 
    (c) no columns are repeated,
or
    returns False
if the random search (or 100n calls to random matrix)
did not find such a matrix.
\end{verbatim}

\textbf{Example:}
\begin{verbatim}
EXAMPLES:
   sage: find_random_incidence_matrix(5, 3)
    False
   sage: find_random_incidence_matrix(5, 3)
    [1 0 0]
    [1 0 0]
    [0 1 1]
    [0 0 1]
    [0 1 0]
\end{verbatim}

\section{\tt \ \ flower\_graph}

\textbf{Syntax:} \verb}flower_graph(N)}\\

\textbf{Description:}
\begin{verbatim}

Returns the flower graph of cycle graphs 
          C_n1, C_n2, ..., C_nk,
where k is the number of *petals* and
where N = [n1,n2,...,nk] is a list of integers ni>2.

If k=1 then it returns the cycle graph C_n1.
\end{verbatim}

\textbf{Example:}
\begin{verbatim}
EXAMPLES:
        sage: N = [3,4,5,6]  
        sage: Gamma = flower_graph(N)  
        sage: Gamma.vertices(sort=True)     
        [0, 1, 2, 3, 4, 5, 6, 7, 8, 9, 10, 11, 12, 13, 14]
        sage: Gamma.edges(labels=None,sort=True)    
        [(0, 1), (0, 2), (0, 3), (0, 5), (0, 6), (0, 9), (0, 10), (0, 14),
         (1, 2), (3, 4), (4, 5), (6, 7), (7, 8), (8, 9), (10, 11),
         (11, 12), (12, 13), (13, 14)]
        sage: Gamma.degree_sequence()    
        [8, 2, 2, 2, 2, 2, 2, 2, 2, 2, 2, 2, 2, 2, 2]
\end{verbatim}

\section{\tt \ \ vertex\_matrix\_map\_to\_CL}

\textbf{Syntax:} \verb}vertex_matrix_map_to_CL(Gamma1, Gamma2, Phi, verbose = False)}\\

\textbf{Description:}

Returns the color list associated to the n2xn1 matrix, if it is ``legal.''
Otherwise, it returns an empty list.

\textbf{Example:}
\begin{verbatim}
EXAMPLES:
        sage: Gamma1 = graphs.CycleGraph(4)     
        sage: Gamma2 = graphs.CycleGraph(8)      
        sage: f0 = [0,1,2,3,0,1,2,3]    
        sage: V2 = Gamma2.vertices(sort=True)    
        sage: f = [V2, f0]         
        sage: is_harmonic_graph_morphism(Gamma1, Gamma2, f) 
        True
        sage: Phi = matrix_of_graph_morphism_vertices(Gamma1, Gamma2, f); Phi 
        [1 0 0 0]
        [0 1 0 0]
        [0 0 1 0]
        [0 0 0 1]
        [1 0 0 0]
        [0 1 0 0]
        [0 0 1 0]
        [0 0 0 1]
        sage: vertex_matrix_map_to_CL(Gamma1, Gamma2, Phi, verbose = False)  
        [0, 1, 2, 3, 0, 1, 2, 3]
        sage: vertex_matrix_map_to_CL(Gamma1, Gamma2, Phi, verbose = True) 
        ([0, 1, 2, 3, 0, 1, 2, 3], True)
\end{verbatim}

\section{\tt \ \ vertical\_edges\_list}

\textbf{Syntax:} \verb}vertical_edges_list(Gamma1, Gamma2, f)}\\

\textbf{Description:}

Returns the list of vertical edges, if any, of the graph morphism represented by the CL f.

\textbf{Example:}
\begin{verbatim}
EXAMPLES:
        sage: Gamma2 = graphs.CycleGraph(4)    
        sage: Gamma1 = graphs.PathGraph(2)       
        sage: f0 = [0,0,1,1]          
        sage: V2 = Gamma2.vertices(sort=True) 
        sage: f = [V2, f0]    
        sage: PhiE = matrix_of_graph_morphism_edges_unsigned(Gamma1, Gamma2, f); PhiE 
        [0]
        [1]
        [1]
        [0]
        sage: vertical_edges_list(Gamma1, Gamma2, f) 
        [(0, 1), (2, 3)]
        sage: Gamma2 = graphs.FriendshipGraph(2)                                           
        sage: Gamma1 = graphs.PathGraph(3)                                                 
        sage: V2 = Gamma2.vertices(sort=True)                                                       
        sage: f0 = [0,0,2,2,1]                                                             
        sage: f = [V2, f0]                                                                 
        sage: is_harmonic_graph_morphism(Gamma1, Gamma2, f)                                
        True
        sage: PhiE = matrix_of_graph_morphism_edges_unsigned(Gamma1, Gamma2, f); PhiE      
        [0 0]
        [1 0]
        [1 0]
        [0 0]
        [0 1]
        [0 1]
        sage: vertical_edges_list(Gamma1, Gamma2, f)                                       
        [(0, 1), (2, 3)]
\end{verbatim}

\section{\tt \ \ adjacency\_matrix\_decomposition}

\textbf{Syntax:} \verb}adjacency_matrix_decomposition(Gamma1, Gamma2, f)}\\

\textbf{Description:}
Returns the decomposition A2 = A\_h+A\_v, where A\_v is supported on the edge positions
associated to vertical edges.

\textbf{Example:}
\begin{verbatim}
EXAMPLES:
        sage: Gamma2 = graphs.FriendshipGraph(2)                                           
        sage: Gamma1 = graphs.PathGraph(3)                                                 
        sage: V2 = Gamma2.vertices(sort=True)                                                       
        sage: f0 = [0,0,2,2,1]                                                             
        sage: f = [V2, f0]                                                                 
        sage: vertical_edges_list(Gamma1, Gamma2, f)                                       
        [(0, 1), (2, 3)]
        sage: adjacency_matrix_decomposition(Gamma1, Gamma2, f)                            
        (
        [0 0 0 0 1]  [0 1 0 0 0]
        [0 0 0 0 1]  [1 0 0 0 0]
        [0 0 0 0 1]  [0 0 0 1 0]
        [0 0 0 0 1]  [0 0 1 0 0]
        [1 1 1 1 0], [0 0 0 0 0]
        )
        sage: A,B = adjacency_matrix_decomposition(Gamma1, Gamma2, f)                      
        sage: A2 = Gamma2.adjacency_matrix()
        sage: A2 == A+B                                                                    
        True
\end{verbatim}

\section{\tt \ \ harmonic\_edge\_irreducible\_subgraph}

\textbf{Syntax:} \verb}harmonic_edge_irreducible_subgraph(Gamma1, Gamma2, f)}\\

\textbf{Description:}
\begin{verbatim}
We say a harmonic morphism phi: Gamma2 -> Gamma1 is *edge irreducible*
if every vertical edge is a bridge. A harmonic morphism with no vertical edges is
automatically edge irreducible.

This function returns the smallest subgraph Gamma of Gamma2 having the same
vertices such that phi|Gamma is a harmonic morphism with no vertical edges.
\end{verbatim}

\textbf{Example:}
\begin{verbatim}
EXAMPLES:
        sage: Gamma2 = graphs.FriendshipGraph(2)                                           
        sage: Gamma1 = graphs.PathGraph(3)                                                 
        sage: Gamma2.edges(labels=None,sort=True)                                                    
        [(0, 1), (0, 4), (1, 4), (2, 3), (2, 4), (3, 4)]
        sage: f0 = [0,0,2,2,1]                                                             
        sage: V2 = Gamma2.vertices(sort=True)                                                       
        sage: f = [V2, f0]  
        sage: harmonic_edge_irreducible_subgraph(Gamma1, Gamma2, f).edges(labels=None,sort=True)              
        [(0, 4), (1, 4), (2, 4), (3, 4)]
\end{verbatim}

\section{\tt \ \ is\_harmonic\_edge\_irreducible}

\textbf{Syntax:} \verb}is_harmonic_edge_irreducible(Gamma1, Gamma2, f)}\\

\textbf{Description:}
\begin{verbatim}

We say a harmonic morphism phi: Gamma2 -> Gamma1 is *edge irreducible* if every vertical edge is a
bridge. A harmonic morphism with no vertical edges is automatically edge irreducible.

This function returns True if every vertical edge is a bridge. Otherwise, it returns False.
\end{verbatim}

\textbf{Example:}
\begin{verbatim}
EXAMPLES:
        sage: Gamma2 = graphs.FriendshipGraph(2)
        sage: Gamma1 = graphs.PathGraph(3)
        sage: f0 = [0,0,2,2,1] 
        sage: V2 = Gamma2.vertices(sort=True) 
        sage: f = [V2, f0]  
        sage: is_harmonic_edge_irreducible(Gamma1, Gamma2, f) 
        False
\end{verbatim}

\section{\tt \ \ coalescence\_vertex}

\textbf{Syntax:} \verb}coalescence_vertex(Gamma1, Gamma2)}\\

\textbf{Description:}
\begin{verbatim}

Returns the graph obtained by identifying the vertex v=0
in Gamma1 with the vertex v=0 in Gamma2, called the 
*vertex coalescence* at 0.

NOTE: This function assumes Gamma1 and Gamma2 are in "standard form",
that is V1 and V2 are integer ranges of the form [0, 1, ..., n-1],
for some n>0.
\end{verbatim}

\textbf{Example:}
\begin{verbatim}
EXAMPLES:
        sage: Gamma1 = graphs.CycleGraph(3)                                                       
        sage: Gamma2 = graphs.CycleGraph(4)                                                       
        sage: Gamma = coalescence_vertex(Gamma2, Gamma1)                                       
        sage: Gamma.vertices(sort=True)                                                                    
        [0, 1, 2, 3, 4, 5]
        sage: Gamma.edges(labels=None,sort=True)                                                            
        [(0, 1), (0, 3), (0, 4), (0, 5), (1, 2), (2, 3), (4, 5)]
\end{verbatim}

\section{\tt \ \ coalescence\_edge}

\textbf{Syntax:} \verb}coalescence_edge(Gamma1, Gamma2)}\\

\textbf{Description:}

Returns the graph obtained by identifying the edge e=(0,1) in Gamma1
with the edge e=(0,1) in Gamma2, called the 
*edge coalescence* at e.

NOTE: This function assumes Gamma1 and Gamma2 are in ``standard form.''

\textbf{Example:}
\begin{verbatim}
EXAMPLES:
        sage: Gamma1 = graphs.CycleGraph(3)                                                       
        sage: Gamma2 = graphs.CycleGraph(4)                                                       
        sage: Gamma = coalescence_edge(Gamma2, Gamma1)                                       
        sage: Gamma.vertices(sort=True)
        [0, 1, 2, 3, 4]
        sage: Gamma.edges(labels=None,sort=True)                                                            
        [(0, 1), (0, 3), (0, 4), (1, 2), (1, 4), (2, 3)]
\end{verbatim}

\section{\tt \ \ barycentric\_subdivision\_graph}

\textbf{Syntax:} \verb}barycentric_subdivision_graph(Gamma)}\\

\textbf{Description:}
Returns the subdivision graph of Gamma obtained by subdividing each edge, 
but leaves Gamma itself unchanged. This always returns a bipartite graph.

\textbf{Example:}
\begin{verbatim}
EXAMPLES:
        sage: Gamma = graphs.ButterflyGraph()            
        sage: subdivision_graph(Gamma).vertices(sort=True)              
         [0, 1, 2, 3, 4, 5, 6, 7, 8, 9, 10]
        sage: Gamma.vertices(sort=True) # to show Gamma is unchanged    
         [0, 1, 2, 3, 4]
        sage: Gamma = graphs.CycleGraph(3)
        sage: Gamma.is_bipartite()
         False
        sage: subdivision_graph(Gamma).is_bipartite()
         True
\end{verbatim}

\section{\tt \ \ which\_edges\_can\_be\_added}

\textbf{Syntax:} \verb}which_edges_can_be_added(Gamma1, Gamma2,f)}\\

\textbf{Description:}
Returns the list of edges that can be added to
    Gamma2 for which f is still harmonic when extended
    to the larger graph.

\textbf{Example:}
\begin{verbatim}
EXAMPLES:
        sage: Gamma2 = graphs.CycleGraph(4)  
        sage: Gamma1 = graphs.PathGraph(3)   
        sage: f0 = [0, 1, 2, 1]    
        sage: V2 = Gamma2.vertices(sort=True)   
        sage: f = [V2, f0]          
        sage: is_harmonic_graph_morphism(Gamma1, Gamma2, f)
        True
        sage: which_edges_can_be_added(Gamma1, Gamma2,f) 
        [(1, 3)]
        sage:
        sage: Gamma2 = graphs.CycleGraph(6) 
        sage: Gamma1 = graphs.CycleGraph(3) 
        sage: f0 = [0, 1, 2, 0, 1, 2] 
        sage: V2 = Gamma2.vertices(sort=True)  
        sage: f = [V2, f0]          
        sage: is_harmonic_graph_morphism(Gamma1, Gamma2, f) 
        True
        sage: which_edges_can_be_added(Gamma1, Gamma2, f) 
        [(0, 3), (1, 4), (2, 5)]
        sage: [horizonal_multiplicity_of_harmonic_map(Gamma1, Gamma2, f, v) for v in V2] 
        [1, 1, 1, 1, 1, 1]
        sage: [vertical_multiplicity_of_harmonic_map(Gamma1, Gamma2, f, v) for v in V2]  
        [0, 0, 0, 0, 0, 0]
        sage: E2 = Gamma2.edges(labels=None,sort=True) 
        sage: E3 = E2 + [(0, 3), (1, 4), (2, 5)] 
        sage: Gamma3 = Graph([V2, E3])     
        sage: is_harmonic_graph_morphism(Gamma1, Gamma3, f)  
        True
        sage: [horizonal_multiplicity_of_harmonic_map(Gamma1, Gamma3, f, v) for v in V2]
        [1, 1, 1, 1, 1, 1]
        sage: [vertical_multiplicity_of_harmonic_map(Gamma1, Gamma3, f, v) for v in V2] 
        [1, 1, 1, 1, 1, 1]
        sage: Gamma2 = graphs.StarGraph(4) # 5-vertex star graph
        sage: Gamma1 = graphs.PathGraph(3)                                              
        sage: f0 = [1, 0, 0, 2, 2]                                                      
        sage: V2 = Gamma2.vertices(sort=True)                                                    
        sage: f = [V2, f0]                                                              
        sage: is_harmonic_graph_morphism(Gamma1, Gamma2, f)                             
        True
        sage: which_edges_can_be_added(Gamma1, Gamma2, f) # adding these creates the butterfly graph      
        [(1, 2), (3, 4)]
        sage: f0 = [1, 0, 1, 1, 2]                                                      
        sage: f = [V2, f0]                                                              
        sage: is_harmonic_graph_morphism(Gamma1, Gamma2, f)                             
        True
        sage: which_edges_can_be_added(Gamma1, Gamma2, f)                               
        [(2, 3)]
        sage: [horizonal_multiplicity_of_harmonic_map(Gamma1, Gamma2, f, v) for v in V2] 
        [1, 1, 0, 0, 1]
        sage: [vertical_multiplicity_of_harmonic_map(Gamma1, Gamma2, f, v) for v in V2] 
        [2, 0, 1, 1, 0]
        sage: E2 = Gamma2.edges(labels=None,sort=True)  
        sage: E3 = E2 + [(2, 3)] 
        sage: Gamma3 = Graph([V2,E3]) # St4+edge
        sage: which_edges_can_be_added(Gamma1, Gamma3, f) 
        []
        sage: is_harmonic_graph_morphism(Gamma1, Gamma3, f)
        True
        sage: [horizonal_multiplicity_of_harmonic_map(Gamma1, Gamma3, f, v) for v in V2]
        [1, 1, 0, 0, 1]
        sage: [vertical_multiplicity_of_harmonic_map(Gamma1, Gamma3, f, v) for v in V2]
        [2, 0, 2, 2, 0]
\end{verbatim}

\section{\tt \ \ one\_vertex\_corona}

\textbf{Syntax:} \verb}one_vertex_corona(Gamma)}\\

\textbf{Description:}
Returns the graph obtained by adding a new vertex 0 to Gamma and 
    an edge from 0 to each vertex in Gamma.

\textbf{Example:}
\begin{verbatim}
EXAMPLES:
        sage: Gamma = graphs.CycleGraph(4)    # cycle graph
        sage: E = Gamma.edges(labels=None,sort=True); E    
        [(0, 1), (0, 3), (1, 2), (2, 3)]
        sage: Gamma0 = one_vertex_corona(Gamma)   # wheel graph
        sage: E0 = Gamma0.edges(labels=None,sort=True); E0  
        [(0, 1), (0, 2), (0, 3), (0, 4), (1, 2), (1, 4), (2, 3), (3, 4)]
\end{verbatim}

\section{\tt \ \ corona}

\textbf{Syntax:} \verb}corona(Gamma1, Gamma2)}\\

\textbf{Description:}
\begin{verbatim}
Returns the graph obtained by taking one copy of Gamma1 and n1 copies
of Gamma2, and then joining the i-th vertex of Gamma1 to every vertex
in the i-th copy of Gamma2.

Assumes both Gamma1 and Gamma2 are in standard form.
  
Here, the 0th copy (start counting at 0) of Gamma2 has vertices
    [n1, n1+1, ..., n1+n2-1],
the 1st copy of Gamma2 has vertices
    [n1+n2, n1+n2+1, ..., n1+2*n2-1], ...,
the jth copy of Gamma2 has vertices
    [n1+j*n2, n1+j*n2+1, ..., n1+(j+1)*n2-1], ...
 
NOTE: The corona graph construction was first introduced by
Fruct and Harary in
   "On the corona of two graphs," Aequationes Mathematicae (1970).
\end{verbatim}

\textbf{Example:}
\begin{verbatim}
EXAMPLES:
        sage: Gamma1 = graphs.PathGraph(2)       
        sage: Gamma2 = graphs.CycleGraph(3)       
        sage: Gamma = corona(Gamma1, Gamma2)      
        sage: Gamma.vertices(sort=True)                    
        [0, 1, 2, 3, 4, 5, 6, 7]
        sage: Gamma.edges(labels=None,sort=True)           
        [(0, 1), (0, 2), (0, 3), (0, 4), (1, 5), (1, 6), (1, 7), (2, 3), (2, 4),
         (3, 4), (5, 6), (5, 7), (6, 7)]
\end{verbatim}

\section{\tt \ \ CL\_of\_join\_of\_morphisms}

\textbf{Syntax:} \verb}CL_of_join_of_morphisms(Gamma1, Gamma2, f, Gamma3, Gamma4, g)}\\

\textbf{Description:}
\begin{verbatim}
All these graphs use the standard ordering, so Gamma_i has vertices
0, 1, ... , n_i - 1, and
       f: Gamma2 -> Gamma1 and  g: Gamma4 -> Gamma3
are morphisms in CL notation.

The graph join of Gamma1 and Gamma2, denoted
                 join(Gamma1, Gamma2),
has vertex set V1 cup V2 (disjoint union), identified with
    {(0, v1) | v1 in V1} cup {(1,v2) | v2 in V2},
and edge set
    E1 cup E2 cup {(v1,v2) | v1 in V1, v2 in V2},
identified with
    {((0, v1), (0,v1')) | v1, v1' in V1} cup
     {((1, v2), (1,v2')) | v2, v2' in V2} cup
      {((0, v1), (1,v2)) | v1 in V1, v2 in V2}.

This is then rewritten in lexicographic order.

Except for the labeling, this is the same as 
       graph_convert_to_standard(Gamma1.join(Gamma2))

Returns the CL associated to join(f,g) going from
          join(Gamma2, Gamma4) -> join(Gamma1, Gamma3).
\end{verbatim}

\textbf{Example:}
\begin{verbatim}
    EXAMPLES:
        sage: Gamma1 = graphs.PathGraph(3)    
        sage: Gamma2 = graphs.PathGraph(4)     
        sage: f0 = [0,1,2,2]     
        sage: V2 = Gamma2.vertices(sort=True)  
        sage: f = [V2, f0]         
        sage: is_harmonic_graph_morphism(Gamma1, Gamma2, f)    
        True
        sage: [horizontal_multiplicity_of_harmonic_map(Gamma1, Gamma2, f, v) for v in V2] 
        [1, 1, 1, 0]
        sage: [vertical_multiplicity_of_harmonic_map(Gamma1, Gamma2, f, v) for v in V2]  
        [0, 0, 1, 1]
        sage: Gamma3 = graphs.PathGraph(2)   
        sage: Gamma4 = graphs.CycleGraph(4) 
        sage: g0 = [0,1,1,0]   
        sage: V4 = Gamma4.vertices(sort=True) 
        sage: g = [V4, g0]   
        sage: is_harmonic_graph_morphism(Gamma3, Gamma4, g)  
        True
        sage: [horizontal_multiplicity_of_harmonic_map(Gamma3, Gamma4, g, v) for v in V4]  
        [1, 1, 1, 1]
        sage: [vertical_multiplicity_of_harmonic_map(Gamma3, Gamma4, g, v) for v in V4]  
        [1, 1, 1, 1]
        sage: CL_of_join_of_morphisms(Gamma1, Gamma2, f, Gamma3, Gamma4, g)  
         [0, 1, 2, 2, 3, 4, 4, 3]
        sage: Gamma_2join4 = graph_convert_to_standard(Gamma2.join(Gamma4))    
        sage: Gamma_1join3 = graph_convert_to_standard(Gamma1.join(Gamma3))   
        sage: V_2join4 = Gamma_2join4.vertices(sort=True)    
        sage: f_join_g_0 = CL_of_join_of_morphisms(Gamma1, Gamma2, f, Gamma3, Gamma4, g) 
        sage: f_join_g = [V_2join4, f_join_g_0]    
        sage: is_harmonic_graph_morphism(Gamma_1join3, Gamma_2join4, f_join_g, verbose=True)   
        Gamma1 = (V1, E1) =  [0, 1, 2, 3, 4] [(0, 1), (0, 3), (0, 4), (1, 2), (1, 3), (1, 4), (2, 3), (2, 4), (3, 4)]
        Gamma2 = (V2, E2) =  [0, 1, 2, 3, 4, 5, 6, 7] [(0, 1), (0, 4), (0, 5), (0, 6), (0, 7), (1, 2), (1, 4), (1, 5), 
                                                       (1, 6), (1, 7), (2, 3), (2, 4), (2, 5), (2, 6), (2, 7), (3, 4), 
                                                       (3, 5), (3, 6), (3, 7), (4, 5), (4, 7), (5, 6), (6, 7)]
        f = [V2, f0], where f0 =  [0, 1, 2, 2, 3, 4, 4, 3]
        The edge e =  (2, 3)  is vertical under f.
        The edge e =  (4, 7)  is vertical under f.
        The edge e =  (5, 6)  is vertical under f.
        This [<vertex>, <intersection sizes>]] fails the check: [0, [1, 2, 2]]
        This [<vertex>, <intersection sizes>]] fails the check: [1, [1, 1, 2, 2]]
        This [<vertex>, <intersection sizes>]] fails the check: [2, [1, 2, 2]]
        This [<vertex>, <intersection sizes>]] fails the check: [3, [0, 2, 2]]
        This [<vertex>, <intersection sizes>]] fails the check: [4, [1, 1, 2, 1]]
        This [<vertex>, <intersection sizes>]] fails the check: [5, [1, 1, 2, 1]]
        This [<vertex>, <intersection sizes>]] fails the check: [6, [1, 1, 2, 1]]
        This [<vertex>, <intersection sizes>]] fails the check: [7, [1, 1, 2, 1]]
        False
        sage: [vertical_multiplicity_of_harmonic_map(Gamma_1join3, Gamma_2join4, f_join_g, v) for v in V_2join4]   
         [0, 0, 1, 1, 1, 1, 1, 1]
        sage: table_of_local_horizontal_multiplicities(Gamma_1join3, Gamma_2join4, f_join_g)                                                       
        [[(0, (0, 1), 1), (0, (0, 3), 2), (0, (0, 4), 2)],
         [(1, (0, 1), 1), (1, (1, 2), 1), (1, (1, 3), 2), (1, (1, 4), 2)],
         [(2, (1, 2), 1), (2, (2, 3), 2), (2, (2, 4), 2)],
         [(3, (1, 2), 0), (3, (2, 3), 2), (3, (2, 4), 2)]]
        sage: Gamma1 = graphs.PathGraph(2)  
        sage: Gamma2 = graphs.CycleGraph(6)   
        sage: V2 = Gamma2.vertices(sort=True) 
        sage: f21_0 = [0,0,0,1,1,1]         
        sage: f21 = [V2, f21_0]              
        sage: Gamma3 = graphs.PathGraph(2)    
        sage: Gamma4 = graphs.CycleGraph(4)   
        sage: f43_0 = [0,1,1,0]     
        sage: f43 = [V4, f43_0]     
        sage: V4 = Gamma4.vertices(sort=True)   
        sage: f43 = [V4, f43_0]   
        sage: [horizontal_multiplicity_of_harmonic_map(Gamma1, Gamma2, f21, v) for v in V2]  
         [1, 0, 1, 1, 0, 1]
        sage: [horizontal_multiplicity_of_harmonic_map(Gamma3, Gamma4, f43, v) for v in V4]            
         [1, 1, 1, 1]
        sage: Gamma_2join4 = graph_convert_to_standard(Gamma2.join(Gamma4))  
        sage: Gamma_1join3 = graph_convert_to_standard(Gamma1.join(Gamma3))   
        sage: V_2join4 = Gamma_2join4.vertices(sort=True)    
        sage: f21_join_f43_0 = CL_of_join_of_morphisms(Gamma1, Gamma2, f21, Gamma3, Gamma4, f43)  
        sage: f21_join_f43 = [V_2join4, f21_join_f43_0]; f21_join_f43       
        [[0, 1, 2, 3, 4, 5, 6, 7, 8, 9], [0, 0, 0, 1, 1, 1, 2, 3, 3, 2]]
        sage: is_harmonic_graph_morphism(Gamma_1join3, Gamma_2join4, f21_join_f43, verbose=True)  
         Gamma1 = (V1, E1) =  [0, 1, 2, 3]
                              [(0, 1), (0, 2), (0, 3), (1, 2), (1, 3), (2, 3)]
         Gamma2 = (V2, E2) =  [0, 1, 2, 3, 4, 5, 6, 7, 8, 9]
                              [(0, 1), (0, 5), (0, 6), (0, 7), (0, 8), (0, 9), (1, 2),
                               (1, 6), (1, 7), (1, 8), (1, 9), (2, 3), (2, 6), (2, 7),
                               (2, 8), (2, 9), (3, 4), (3, 6), (3, 7), (3, 8), (3, 9), (4, 5),
                               (4, 6), (4, 7), (4, 8), (4, 9), (5, 6), (5, 7), (5, 8), (5, 9),
                               (6, 7), (6, 9), (7, 8), (8, 9)]
        f = [V2, f0], where f0 =  [0, 0, 0, 1, 1, 1, 2, 3, 3, 2]
        The edge e =  (0, 1)  is vertical under f.
        The edge e =  (1, 2)  is vertical under f.
        The edge e =  (3, 4)  is vertical under f.
        The edge e =  (4, 5)  is vertical under f.
        The edge e =  (6, 9)  is vertical under f.
        The edge e =  (7, 8)  is vertical under f.
        This [<vertex>, <intersection sizes>]] fails the check: [0, [1, 2, 2]]
        This [<vertex>, <intersection sizes>]] fails the check: [1, [0, 2, 2]]
        This [<vertex>, <intersection sizes>]] fails the check: [2, [1, 2, 2]]
        This [<vertex>, <intersection sizes>]] fails the check: [3, [1, 2, 2]]
        This [<vertex>, <intersection sizes>]] fails the check: [4, [0, 2, 2]]
        This [<vertex>, <intersection sizes>]] fails the check: [5, [1, 2, 2]]
        This [<vertex>, <intersection sizes>]] fails the check: [6, [3, 3, 1]]
        This [<vertex>, <intersection sizes>]] fails the check: [7, [3, 3, 1]]
        This [<vertex>, <intersection sizes>]] fails the check: [8, [3, 3, 1]]
        This [<vertex>, <intersection sizes>]] fails the check: [9, [3, 3, 1]]
        False
        sage: table_of_local_horizontal_multiplicities(Gamma_1join3, Gamma_2join4, f21_join_f43)                                                   
        [[(0, (0, 1), 1), (0, (0, 2), 2), (0, (0, 3), 2)],
         [(1, (0, 1), 0), (1, (0, 2), 2), (1, (0, 3), 2)],
         [(2, (0, 1), 1), (2, (0, 2), 2), (2, (0, 3), 2)],
         [(3, (0, 1), 1), (3, (1, 2), 2), (3, (1, 3), 2)],
         [(4, (0, 1), 0), (4, (1, 2), 2), (4, (1, 3), 2)],
         [(5, (0, 1), 1), (5, (1, 2), 2), (5, (1, 3), 2)]]
\end{verbatim}

\section{\tt \ \ CL\_of\_product\_of\_morphisms}

\textbf{Syntax:} \verb}CL_of_product_of_morphisms(Gamma1, Gamma2, f, Gamma3, Gamma4, g)}\\

\textbf{Description:}
\begin{verbatim}
Returns the CL associated to fxg going from
          Gamma2 x Gamma4 -> Gamma1 x Gamma3
in the lexicographical ordering.

 NOTE: Does not depend on which product (Cartesian, tensor, etc) is used.

Notation:
          - f represents a harmonic morphism Gamma2 -> Gamma1
          - g represents a harmonic morphism Gamma4 -> Gamma3
            (all these graphs use the standard ordering, so Gamma\_i has vertices
             0, 1, ...., n\_i - 1). Order the vertices of Gamma2 x Gamma4 
             and of Gamma1 x Gamma3 in the lexicographical ordering.
\end{verbatim}

The CL for fxg is returned using this ordering.

\textbf{Example:}
\begin{verbatim}
EXAMPLES:
        sage: Gamma1 = graphs.PathGraph(3)    
        sage: Gamma2 = graphs.PathGraph(4)     
        sage: f0 = [0,1,2,2]     
        sage: V2 = Gamma2.vertices(sort=True)  
        sage: f = [V2, f0]         
        sage: is_harmonic_graph_morphism(Gamma1, Gamma2, f)    
        True
        sage: Gamma3 = graphs.PathGraph(2)   
        sage: Gamma4 = graphs.CycleGraph(4) 
        sage: g0 = [0,1,1,0]   
        sage: V4 = Gamma4.vertices(sort=True) 
        sage: g = [V4, g0]   
        sage: is_harmonic_graph_morphism(Gamma3, Gamma4, g)  
        True
        sage: CL_of_product_of_morphisms(Gamma1, Gamma2, f, Gamma3, Gamma4, g)    
        [0, 1, 1, 0, 2, 3, 3, 2, 4, 5, 5, 4, 4, 5, 5, 4]
\end{verbatim}

\section{\tt \ \ matrix\_of\_graph\_of\_product\_of\_morphisms\_vertices}

\textbf{Syntax:} \verb}matrix_of_graph_of_product_of_morphisms_vertices(Gamma1, Gamma2, f, Gamma3, Gamma4, g)}\\

\textbf{Description:}
\begin{verbatim}

Returns the vertex matrix map associated to fxg going from
          Gamma2 x Gamma4 -> Gamma1 x Gamma3
in the lexicographical ordering.

Notation: f represents a harmonic morphism Gamma2 -> Gamma1
          g represents a harmonic morphism Gamma4 -> Gamma3

    (all these graphs use the standard ordering, so Gamma\_i has vertices
    0, 1, ...., n\_i - 1). Order the vertices of Gamma2 x Gamma4 
    and of Gamma1 x Gamma3 in the lexicographical ordering. 
\end{verbatim}

\textbf{Example:}
\begin{verbatim}
EXAMPLES:
        sage: Gamma1 = graphs.PathGraph(2)
        sage: Gamma2 = graphs.PathGraph(3)
        sage: f0 = [0,1,1]       
        sage: V2 = Gamma2.vertices(sort=True) 
        sage: f = [V2, f0]     
        sage: Phi = matrix_of_graph_morphism_vertices(Gamma1, Gamma2, f)
        [1 0]
        [0 1]
        [0 1]
        sage: Gamma3 = graphs.PathGraph(2) 
        sage: Gamma4 = graphs.CycleGraph(4) 
        sage: V4 = Gamma4.vertices(sort=True) 
        sage: g0 = [0,1,1,0]     
        sage: g = [V4, g0]  
        sage: Phi_g = matrix_of_graph_morphism_vertices(Gamma3, Gamma4, g)
        [1 0]
        [0 1]
        [0 1]
        [1 0]
        sage: matrix_of_graph_of_product_of_morphisms_vertices(Gamma1, Gamma2, f, Gamma3, Gamma4, g)              
        [1 0 0 0]
        [0 1 0 0]
        [0 1 0 0]
        [1 0 0 0]
        [0 0 1 0]
        [0 0 0 1]
        [0 0 0 1]
        [0 0 1 0]
        [0 0 1 0]
        [0 0 0 1]
        [0 0 0 1]
        [0 0 1 0]
\end{verbatim}

\section{\tt \ \ CL\_product\_projection1\_harmonic}

\textbf{Syntax:} \verb}CL_product_projection1_harmonic(Gamma1, Gamma2)}\\

\textbf{Description:}
\begin{verbatim}
Returns CL in the form [V2, f0] of pr_1 (proj of product onto 1st factor) from 
Gamma1xGamma2 to Gamma1, and False otherwise, where V2 is the list of vertices
of Gamma1xGamma2 in standard form.

NOTE: Assumes Gamma1, Gamma2 are in standard form.

NOTE: While Cartesian product is used below, this CL does not depend on
which product (Cartesian, tensor, strong, etc) is used.
\end{verbatim}

\textbf{Example:}
\begin{verbatim}
    EXAMPLES:
        sage: Gamma1 = graphs.PathGraph(2)
        sage: Gamma2 = graphs.PathGraph(3)
        sage: CL_product_projection1_harmonic(Gamma1, Gamma2) 
        [[0, 1, 2, 3, 4, 5], [0, 0, 0, 1, 1, 1]]
	sage: f = CL_product_projection1_harmonic(Gamma1, Gamma2)
	sage: Gamma12 = Gamma1.cartesian_product(Gamma2)
	sage: Gamma12_0 = graph_convert_to_standard(Gamma12)  
        sage: is_harmonic_graph_morphism(Gamma1, Gamma12_0, f) 
        True

        sage: is_cartesian_product_projection1_harmonic(Gamma1, Gamma2, f) 

        sage: Gamma1 = graphs.PathGraph(2) 
        sage: Gamma2 = graphs.CycleGraph(4)
	sage: CL_product_projection1_harmonic(Gamma1, Gamma2) 
        [[0, 1, 2, 3, 4, 5, 6, 7], [0, 0, 0, 0, 1, 1, 1, 1]]
	sage: f = CL_product_projection1_harmonic(Gamma1, Gamma2)
	sage: Gamma12 = Gamma1.cartesian_product(Gamma2)
	sage: Gamma12_0 = graph_convert_to_standard(Gamma12)  
        sage: is_harmonic_graph_morphism(Gamma1, Gamma12_0, f) 
        True
\end{verbatim}

\section{\tt \ \ CL\_product\_projection2\_harmonic}

\textbf{Syntax:} \verb}CL_product_projection2_harmonic(Gamma1, Gamma2)}\\

\textbf{Description:}
\begin{verbatim}
Returns CL in the form [V2, f0] of pr_2 (proj of product onto 2nd factor) from 
Gamma1xGamma2 to Gamma1, and False otherwise, where V2 is the list of vertices
of Gamma1xGamma2 in standard form.

NOTE: Assumes Gamma1, Gamma2 are in standard form.

NOTE: While Cartesian product is used below, this CL does not depend on
which product (Cartesian, tensor, strong, etc) is used.
\end{verbatim}

\textbf{Example:}
\begin{verbatim}
   EXAMPLES:
        sage: Gamma1 = graphs.PathGraph(2)
        sage: Gamma2 = graphs.PathGraph(3)
        sage: CL_product_projection2_harmonic(Gamma1, Gamma2) 
        [[0, 1, 2, 3, 4, 5], [0, 1, 2, 0, 1, 2]]
	sage: f = CL_product_projection2_harmonic(Gamma1, Gamma2)
	sage: Gamma12 = Gamma1.cartesian_product(Gamma2)
	sage: Gamma12_0 = graph_convert_to_standard(Gamma12)  
        sage: is_harmonic_graph_morphism(Gamma2, Gamma12_0, f) 
        True
        sage: Gamma1 = graphs.PathGraph(2) 
        sage: Gamma2 = graphs.CycleGraph(4)
	sage: CL_product_projection2_harmonic(Gamma1, Gamma2) 
        [[0, 1, 2, 3, 4, 5, 6, 7], [0, 1, 2, 3, 0, 1, 2, 3]]
	sage: f = CL_product_projection2_harmonic(Gamma1, Gamma2)
	sage: Gamma12 = Gamma1.cartesian_product(Gamma2)
	sage: Gamma12_0 = graph_convert_to_standard(Gamma12)  
        sage: is_harmonic_graph_morphism(Gamma2, Gamma12_0, f) 
        True
\end{verbatim}

\section{\tt \ \ is\_cartesian\_product\_projection1\_harmonic}

\textbf{Syntax:} \verb}is_cartesian_product_projection1_harmonic(Gamma1, Gamma2, verbose=False)}\\

\textbf{Description:}
\begin{verbatim}
Returns True if pr_1 (proj of cartesian product onto 1st factor) defines a 
harmonic morphism from Gamma1xGamma2 to Gamma1, and False otherwise. 

Assumes Gamma1, Gamma2 are in standard form.
\end{verbatim}

\textbf{Example:}
\begin{verbatim}
    EXAMPLES:
        sage: Gamma1 = graphs.PathGraph(2)
        sage: Gamma2 = graphs.PathGraph(3)
        sage: is_cartesian_product_projection1_harmonic(Gamma1, Gamma2) 
        True
        sage: is_cartesian_product_projection1_harmonic(Gamma1, Gamma2, verbose=True) 
        This [<vertex>, <intersection sizes>]] passes the check: [0, [1]]
        This [<vertex>, <intersection sizes>]] passes the check: [1, [1]]
        This [<vertex>, <intersection sizes>]] passes the check: [2, [1]]
        This [<vertex>, <intersection sizes>]] passes the check: [3, [1]]
        This [<vertex>, <intersection sizes>]] passes the check: [4, [1]]
        This [<vertex>, <intersection sizes>]] passes the check: [5, [1]]
        True
\end{verbatim}

\section{\tt \ \ is\_cartesian\_product\_of\_graph\_morphisms\_harmonic}

\textbf{Syntax:} \verb}is_cartesian_product_of_graph_morphisms_harmonic(Gamma1, Gamma2, f, Gamma3, Gamma4, g, verbose=False)}\\

\textbf{Description:}
\begin{verbatim}
Returns True if fxg (cartesian product) defines a graph-theoretic mapping
from Gamma2xGamma4 to Gamma1xGamma3 that is harmonic, and False otherwise. 

Assumes Gamma1, Gamma2, Gamma3, Gamma4 are in standard form.
\end{verbatim}

\textbf{Example:}
\begin{verbatim}
    EXAMPLES:
        sage: Gamma1 = graphs.PathGraph(2)
        sage: Gamma2 = graphs.PathGraph(3)
        sage: f0 = [0,1,1]       
        sage: V2 = Gamma2.vertices(sort=True) 
        sage: f = [V2, f0]     
        sage: Gamma3 = graphs.PathGraph(2) 
        sage: Gamma4 = graphs.CycleGraph(4) 
        sage: V4 = Gamma4.vertices(sort=True) 
        sage: g0 = [0,1,1,0]     
        sage: g = [V4, g0]  
        sage: [vertical_multiplicity_of_harmonic_map(Gamma3, Gamma4, g, v) for v in V4]  
        [1, 1, 1, 1]
        sage: [horizonal_multiplicity_of_harmonic_map(Gamma3, Gamma4, g, v) for v in V4] 
        [1, 1, 1, 1]
        sage: is_cartesian_product_of_graph_morphisms_harmonic(Gamma1, Gamma2, f, Gamma3, Gamma4, g) 
        False
        sage: Gamma1 = graphs.PathGraph(2) 
        sage: Gamma2 = graphs.CycleGraph(4) 
        sage: f0 = [0,1,1, 0]    
        sage: V2 = Gamma2.vertices(sort=True)    
        sage: f = [V2, f0]                         
        sage: is_cartesian_product_of_graph_morphisms_harmonic(Gamma1, Gamma2, f, Gamma3, Gamma4, g)  
        True
        sage: Gamma13 = Gamma1.cartesian_product(Gamma3) 
        sage: Gamma1x3 = graph_convert_to_standard(Gamma13)
        sage: Gamma24 = Gamma2.cartesian_product(Gamma4)   
        sage: Gamma2x4 = graph_convert_to_standard(Gamma24) 
        sage: V2x4 = Gamma2x4.vertices(sort=True)
        sage: fxg0 = CL_of_product_of_morphisms(Gamma1, Gamma2, f, Gamma3, Gamma4, g) 
        sage: fxg = [V2x4, fxg0]  
        sage: [vertical_multiplicity_of_harmonic_map(Gamma1x3, Gamma2x4, fxg, v) for v in V2x4]  
        [2, 2, 2, 2, 2, 2, 2, 2, 2, 2, 2, 2, 2, 2, 2, 2]
        sage: [horizonal_multiplicity_of_harmonic_map(Gamma1x3, Gamma2x4, fxg, v) for v in V2x4] 
        [1, 1, 1, 1, 1, 1, 1, 1, 1, 1, 1, 1, 1, 1, 1, 1]
\end{verbatim}

\section{\tt \ \ is\_cartesian\_product\_of\_graph\_morphisms\_a\_morphism}

\textbf{Syntax:} \verb}is_cartesian_product_of_graph_morphisms_a_morphism(Gamma1, Gamma2, f, Gamma3, Gamma4, g, verbose=False)}\\

\textbf{Description:}
\begin{verbatim}
Returns True if fxg (Cartesian product map) defines a graph-theoretic mapping
from Gamma2xGamma4 to Gamma1xGamma3 that is a morphism, and False otherwise. 

Assumes Gamma1, Gamma2, Gamma3, Gamma4 are in standard form.
\end{verbatim}

\textbf{Example:}
\begin{verbatim}
    EXAMPLES:
        sage: Gamma1 = graphs.PathGraph(2)
        sage: Gamma2 = graphs.PathGraph(3)
        sage: f0 = [0,1,1]       
        sage: V2 = Gamma2.vertices(sort=True) 
        sage: f = [V2, f0]     
        sage: [vertical_multiplicity_of_harmonic_map(Gamma1, Gamma2, f, v) for v in V2] 
        [0, 1, 1]
        sage: [horizonal_multiplicity_of_harmonic_map(Gamma1, Gamma2, f, v) for v in V2]
        [1, 1, 0]
        sage: Gamma3 = graphs.PathGraph(2) 
        sage: Gamma4 = graphs.CycleGraph(4) 
        sage: V4 = Gamma4.vertices(sort=True) 
        sage: g0 = [0,1,1,0]     
        sage: g = [V4, g0]  
        sage: [vertical_multiplicity_of_harmonic_map(Gamma3, Gamma4, g, v) for v in V4] 
        [1, 1, 1, 1]
        sage: [horizonal_multiplicity_of_harmonic_map(Gamma3, Gamma4, g, v) for v in V4]  
        [1, 1, 1, 1]
        sage: is_cartesian_product_of_graph_morphisms_a_morphism(Gamma1, Gamma2, f, Gamma3, Gamma4, g, verbose=True)  
        (0, 2) [(0, 1), (0, 3), (1, 2), (1, 4), (2, 5), (3, 4), (4, 5)]
        Does not send edges to {edges} union {vertices}
        False
        sage: is_cartesian_product_of_graph_morphisms_a_morphism(Gamma1, Gamma2, f, Gamma3, Gamma4, g) 
        False
\end{verbatim}

\section{\tt \ \ is\_tensor\_product\_projection1\_harmonic}

\textbf{Syntax:} \verb}is_tensor_product_projection1_harmonic(Gamma1, Gamma2, verbose=False)}\\

\textbf{Description:}
\begin{verbatim}
Imagine you have two graphs: Gamma1 and Gamma2. Now suppose you build a new graph where:
The vertices are pairs like 
                         (v1, v2),
one from Gamma1, one from Gamma1. There's an 
           edge between (v1, v2) and (v1', v2') 
if both v1 is adjacent to v1' in Gamma1 and v2 is adjacent to v2' in Gamma2.

This is called the tensor product of graphs, often written as 
           Gamma1 \otimes Gamma2
or as 
           Gamma1 x_{tensor} Gamma2

Returns 
        True, if pr_1 (proj of tensor product onto 1st factor) defines a 
                 harmonic morphism from Gamma1xGamma2 to Gamma1, 
and 
        False, otherwise. 

Assumes Gamma1, Gamma2 are in standard form.
\end{verbatim}

\textbf{Example:}
\begin{verbatim}
    EXAMPLES:
        sage: Gamma1 = graphs.PathGraph(2)
        sage: Gamma2 = graphs.PathGraph(3)
        sage: is_tensor_product_projection1_harmonic(Gamma1, Gamma2) 
        True
        sage: is_tensor_product_projection1_harmonic(Gamma1, Gamma2, verbose=True) 
        This [<vertex>, <intersection sizes>]] passes the check: [0, [1]]
        This [<vertex>, <intersection sizes>]] passes the check: [1, [2]]
        This [<vertex>, <intersection sizes>]] passes the check: [2, [1]]
        This [<vertex>, <intersection sizes>]] passes the check: [3, [1]]
        This [<vertex>, <intersection sizes>]] passes the check: [4, [2]]
        This [<vertex>, <intersection sizes>]] passes the check: [5, [1]]
        True
\end{verbatim}

\section{\tt \ \ is\_tensor\_product\_of\_graph\_morphisms\_harmonic}

\textbf{Syntax:} \verb}is_tensor_product_of_graph_morphisms_harmonic(Gamma1, Gamma2, f, Gamma3, Gamma4, g, verbose=False)}\\

\textbf{Description:}
\begin{verbatim}
Returns
     True, if fxg (tensor product) defines a graph-theoretic mapping
           from Gamma2xGamma4 to Gamma1xGamma3 that is harmonic, and
     False, otherwise. 

Assumes Gamma1, Gamma2, Gamma3, Gamma4 are in standard form.
\end{verbatim}

\textbf{Example:}
\begin{verbatim}
    EXAMPLES:
       sage: V2 = [0, 1, 2, 3, 4, 5]                                                                            
       sage: E2 = [(0, 1), (0, 2), (0, 4), (0, 5), (1, 2), (1, 3), (1, 4), (2, 3), (2, 5), (3, 4), (3, 5), (4, 5)] 
       sage: Gamma2 = Graph([V2,E2]) # d4r6a                                                                     
       sage: Gamma2.chromatic_number()                                                                           
       3
       sage: Gamma1 = graphs.CycleGraph(3)                                                                       
       sage: f0 = [2, 0, 1, 2, 1, 0]                                                                             
       sage: f = [V2, f0]   
       sage: [horizonal_multiplicity_of_harmonic_map(Gamma1, Gamma2, f, v) for v in V2]                          
       [2, 2, 2, 2, 2, 2]
       sage: [vertical_multiplicity_of_harmonic_map(Gamma1, Gamma2, f, v) for v in V2]                           
       [0, 0, 0, 0, 0, 0]                                                                 
       sage: is_tensor_product_of_graph_morphisms_a_morphism(Gamma1, Gamma2, f, Gamma1, Gamma2, f) 
       True
       sage: is_tensor_product_of_graph_morphisms_harmonic(Gamma1, Gamma2, f, Gamma1, Gamma2, f)   
       True
       sage: Gamma13 = Gamma1.tensor_product(Gamma1)                                                             
       sage: Gamma24 = Gamma2.tensor_product(Gamma2)                                                             
       sage: V24 = Gamma24.vertices(sort=True)                                                                            
       sage: fxg0 = CL_of_product_of_morphisms(Gamma1, Gamma2, f, Gamma1, Gamma2, f)                             
       sage: Gamma2x4 = graph_convert_to_standard(Gamma24)                                                       
       sage: Gamma1x3 = graph_convert_to_standard(Gamma13)                                                       
       sage: V2x4 = Gamma2x4.vertices(sort=True)                                                                          
       sage: fxg = [V2x4, fxg0]                                                                                  
       sage: is_harmonic_graph_morphism(Gamma1, Gamma2, f)                                                       
       True
       sage: is_harmonic_graph_morphism(Gamma1x3, Gamma2x4, fxg)                                                 
       True
\end{verbatim}

\section{\tt \ \ is\_tensor\_product\_of\_graph\_morphisms\_a\_morphism}

\textbf{Syntax:} \verb}is_tensor_product_of_graph_morphisms_a_morphism(Gamma1, Gamma2, f, Gamma3, Gamma4, g, verbose=False)}\\

\textbf{Description:}
\begin{verbatim}
Returns 
         True if fxg (tensor product map) defines a graph-theoretic mapping
                 from Gamma2xGamma4 to Gamma1xGamma3 that is a morphism, 
and 
         False otherwise. 

Assumes Gamma1, Gamma2, Gamma3, Gamma4 are in standard form.
\end{verbatim}

\textbf{Example:}
\begin{verbatim}
    EXAMPLES:
        sage: Gamma1 = graphs.PathGraph(2)
        sage: Gamma2 = graphs.PathGraph(3)
        sage: f0 = [0,1,1]       
        sage: V2 = Gamma2.vertices(sort=True) 
        sage: f = [V2, f0]     
        sage: [vertical_multiplicity_of_harmonic_map(Gamma1, Gamma2, f, v) for v in V2] 
        [0, 1, 1]
        sage: [horizonal_multiplicity_of_harmonic_map(Gamma1, Gamma2, f, v) for v in V2]
        [1, 1, 0]
        sage: Gamma3 = graphs.PathGraph(2) 
        sage: Gamma4 = graphs.CycleGraph(4) 
        sage: V4 = Gamma4.vertices(sort=True) 
        sage: g0 = [0,1,1,0]     
        sage: g = [V4, g0]  
        sage: [vertical_multiplicity_of_harmonic_map(Gamma3, Gamma4, g, v) for v in V4] 
        [1, 1, 1, 1]
        sage: [horizonal_multiplicity_of_harmonic_map(Gamma3, Gamma4, g, v) for v in V4]  
        [1, 1, 1, 1]
        sage: is_tensor_product_of_graph_morphisms_a_morphism(Gamma1, Gamma2, f, Gamma3, Gamma4, g, verbose=True)  
        The Gamma2 edge (0, 2)  does not map to Gamma1 edge set 
          [(0, 1), (0, 3), (1, 2), (1, 4), (2, 5), (3, 4), (4, 5)]  or vertex set  [0, 1, 2, 3, 4, 5]
        Error: Does not send edges to {edges} union {vertices}
        False
        sage: is_tensor_product_of_graph_morphisms_a_morphism(Gamma1, Gamma2, f, Gamma3, Gamma4, g) 
        False
\end{verbatim}

\section{\tt \ \ is\_strong\_product\_projection1\_harmonic}

\textbf{Syntax:} \verb}is_strong_product_projection1_harmonic(Gamma1, Gamma2, verbose=False)}\\

\textbf{Description:}
\begin{verbatim}
Imagine you have two graphs: Gamma1 and Gamma2. Now suppose you build a new graph where:
The vertices are pairs like 
                         (v1, v2),
one from Gamma1, one from Gamma1. There's an 
           edge between (v1, v2) and (v1', v2') 
if either 
   (a) both v1 = v1' in Gamma1 and v2 connects to v2' in Gamma2,
or
   (b) both v1 connects to v1' in Gamma1 and v2 = v2' in Gamma2,
or
   (a) both v1 and v1' are adjacent in Gamma1 and v2 is adjacent to v2' in Gamma2.

This is called the strong product of graphs, often written as 
           Gamma1 x_{strong} Gamma2

Returns True if pr_1 (proj of strong product onto 1st factor) defines a 
harmonic morphism from Gamma1xGamma2 to Gamma1, and False otherwise. 

    Assumes Gamma1, Gamma2 are in standard form.
\end{verbatim}

\textbf{Example:}
\begin{verbatim}
    EXAMPLES:
        sage: Gamma1 = graphs.PathGraph(2)
        sage: Gamma2 = graphs.PathGraph(3)
        sage: is_strong_product_projection1_harmonic(Gamma1, Gamma2) 
        True
        sage: is_strong_product_projection1_harmonic(Gamma1, Gamma2, verbose=True) 
        This [<vertex>, <intersection sizes>]] passes the check: [0, [2]]
        This [<vertex>, <intersection sizes>]] passes the check: [1, [3]]
        This [<vertex>, <intersection sizes>]] passes the check: [2, [2]]
        This [<vertex>, <intersection sizes>]] passes the check: [3, [2]]
        This [<vertex>, <intersection sizes>]] passes the check: [4, [3]]
        This [<vertex>, <intersection sizes>]] passes the check: [5, [2]]
        True
\end{verbatim}

\section{\tt \ \ is\_strong\_product\_of\_graph\_morphisms\_harmonic}

\textbf{Syntax:} \verb}is_strong_product_of_graph_morphisms_harmonic(Gamma1, Gamma2, f, Gamma3, Gamma4, g, verbose=False)}\\

\textbf{Description:}
\begin{verbatim}
Returns True if fxg (strong product) defines a graph-theoretic mapping
from Gamma2 to Gamma1 that is harmonic, and False otherwise. 

    Assumes Gamma1, Gamma2, Gamma3, Gamma4 are in standard form.
\end{verbatim}

\textbf{Example:}
\begin{verbatim}
    EXAMPLES:
        sage: Gamma1 = graphs.PathGraph(2)
        sage: Gamma2 = graphs.PathGraph(3)
        sage: Gamma3 = graphs.CycleGraph(4)
        sage: Gamma4 = graphs.CycleGraph(4) 
        sage: V2 = Gamma2.vertices(sort=True)
        sage: V4 = Gamma4.vertices(sort=True)
        sage: f0 = [0, 1, 1]
        sage: g0 = [0,1,2,3] 
        sage: f = [V2, f0] 
        sage: g = [V4, g0]
        sage: is_harmonic_graph_morphism(Gamma1, Gamma2, f)
        True
        sage: is_harmonic_graph_morphism(Gamma3, Gamma4, g)
        True
        sage: [horizonal_multiplicity_of_harmonic_map(Gamma1, Gamma2, f, v) for v in V2]
        [1, 1, 0]
        sage: [vertical_multiplicity_of_harmonic_map(Gamma1, Gamma2, f, v) for v in V2]
        [0, 1, 1]
        sage: [horizonal_multiplicity_of_harmonic_map(Gamma3, Gamma4, g, v) for v in V4]
        [1, 1, 1, 1]
        sage: [vertical_multiplicity_of_harmonic_map(Gamma3, Gamma4, g, v) for v in V4]
        [0, 0, 0, 0]
        sage: is_strong_product_of_graph_morphisms_harmonic(Gamma1, Gamma2, f, Gamma3, Gamma4, g, verbose=False)  
        False
        sage: is_strong_product_of_graph_morphisms_harmonic(Gamma1, Gamma2, f, Gamma3, Gamma4, g, verbose=True)   
        This [<vertex>, <intersection sizes>]] passes the check: [0, [1, 1, 1, 1, 1]]
        This [<vertex>, <intersection sizes>]] passes the check: [1, [1, 1, 1, 1, 1]]
        This [<vertex>, <intersection sizes>]] passes the check: [2, [1, 1, 1, 1, 1]]
        This [<vertex>, <intersection sizes>]] passes the check: [3, [1, 1, 1, 1, 1]]
        This [<vertex>, <intersection sizes>]] fails the check: [4, [1, 1, 1, 2, 2]]
        This [<vertex>, <intersection sizes>]] fails the check: [5, [1, 1, 1, 2, 2]]
        This [<vertex>, <intersection sizes>]] fails the check: [6, [1, 1, 1, 2, 2]]
        This [<vertex>, <intersection sizes>]] fails the check: [7, [1, 1, 1, 2, 2]]
        This [<vertex>, <intersection sizes>]] fails the check: [8, [0, 0, 0, 2, 2]]
        This [<vertex>, <intersection sizes>]] fails the check: [9, [0, 0, 0, 2, 2]]
        This [<vertex>, <intersection sizes>]] fails the check: [10, [0, 0, 0, 2, 2]]
        This [<vertex>, <intersection sizes>]] fails the check: [11, [0, 0, 0, 2, 2]]
        False
        sage: is_strong_product_of_graph_morphisms_harmonic(Gamma3, Gamma4, g, Gamma3, Gamma4, g)                 
        True
        sage: Gamma1 = graphs.CycleGraph(3)                                                                       
        sage: Gamma2 = graphs.CycleGraph(6)                                                                       
        sage: V2 = Gamma2.vertices(sort=True)                                                                              
        sage: f0 = [0, 1, 2, 0, 1, 2]                                                                             
        sage: f = [V2, f0]                                                                                        
        sage: is_harmonic_graph_morphism(Gamma1, Gamma2, f)                                                       
        True
        sage: [horizonal_multiplicity_of_harmonic_map(Gamma1, Gamma2, f, v) for v in V2] 
        [1, 1, 1, 1, 1, 1]
        sage: [vertical_multiplicity_of_harmonic_map(Gamma1, Gamma2, f, v) for v in V2]   
        [0, 0, 0, 0, 0, 0]
        sage: Gamma3 = graphs.CycleGraph(4)                                                                       
        sage: Gamma4 = graphs.CycleGraph(8)                                                                       
        sage: V4 = Gamma4.vertices(sort=True)                                                                              
        sage: g0 = [0,1,2,3,0,1,2,3]                                                                              
        sage: g = [V4, g0]                                                                                        
        sage: is_harmonic_graph_morphism(Gamma3, Gamma4, g)                                                       
        True
        sage: [horizonal_multiplicity_of_harmonic_map(Gamma1, Gamma2, f, v) for v in V2]                          
        [1, 1, 1, 1, 1, 1]
        sage: [vertical_multiplicity_of_harmonic_map(Gamma1, Gamma2, f, v) for v in V2]                           
        [0, 0, 0, 0, 0, 0]
        sage: is_strong_product_of_graph_morphisms_harmonic(Gamma3, Gamma4, g, Gamma3, Gamma4, g)                 
        True
        sage: is_strong_product_of_graph_morphisms_harmonic(Gamma1, Gamma2, f, Gamma3, Gamma4, g)   
        True
\end{verbatim}

\section{\tt \ \ is\_strong\_product\_of\_graph\_morphisms\_a\_morphism}

\textbf{Syntax:} \verb}is_strong_product_of_graph_morphisms_a_morphism(Gamma1, Gamma2, f, Gamma3, Gamma4, g, verbose=False)}\\

\textbf{Description:}
\begin{verbatim}
Returns True if fxg (strong product map) defines a graph-theoretic mapping
from Gamma2 to Gamma1 that is a morphism, and False otherwise. 

Assumes Gamma1, Gamma2, Gamma3, Gamma4 are in standard form.
\end{verbatim}

\textbf{Example:}
\begin{verbatim}
    EXAMPLES:
        sage: Gamma1 = graphs.PathGraph(2)
        sage: Gamma2 = graphs.PathGraph(3)
        sage: f0 = [0,1,1]       
        sage: V2 = Gamma2.vertices(sort=True) 
        sage: f = [V2, f0]     
        sage: [vertical_multiplicity_of_harmonic_map(Gamma1, Gamma2, f, v) for v in V2] 
        [0, 1, 1]
        sage: [horizonal_multiplicity_of_harmonic_map(Gamma1, Gamma2, f, v) for v in V2]
        [1, 1, 0]
        sage: Gamma3 = graphs.PathGraph(2) 
        sage: Gamma4 = graphs.CycleGraph(4) 
        sage: V4 = Gamma4.vertices(sort=True) 
        sage: g0 = [0,1,1,0]     
        sage: g = [V4, g0]  
        sage: [vertical_multiplicity_of_harmonic_map(Gamma3, Gamma4, g, v) for v in V4] 
        [1, 1, 1, 1]
        sage: [horizonal_multiplicity_of_harmonic_map(Gamma3, Gamma4, g, v) for v in V4]  
        [1, 1, 1, 1]
        sage: is_strong_product_of_graph_morphisms_a_morphism(Gamma1, Gamma2, f, Gamma3, Gamma4, g) 
        False
        sage: is_strong_product_of_graph_morphisms_a_morphism(Gamma1, Gamma2, f, Gamma3, Gamma4, g, verbose=True)
        The edge e =  (0, 3)  has no image in Gamma1 under f.
        False
\end{verbatim}

\section{\tt \ \ is\_lexicographic\_product\_of\_graph\_morphisms\_harmonic}

  \textbf{Syntax:} \verb}is_lexicographic_product_of_graph_morphisms_harmonic(Gamma1, Gamma2, f, Gamma3, Gamma4, g, verbose=False)}
  
\textbf{Description:}
\begin{verbatim}
Imagine you have two graphs: Gamma1 and Gamma2. Now suppose you build a new graph where:
The vertices are pairs like 
                         (v1, v2),
one from Gamma1, one from Gamma1. There's an 
           edge between (v1, v2) and (v1', v2') 
if either 
   (a) both v1 = v1' in Gamma1 and v2 is adjacent to v2' in Gamma2,
or
   (b) v1 is adjacent to v1' in Gamma1 (regardless of v2, v2' in Gamma2).

This is called the lexicographical product of graphs, often written as 
           Gamma1 x_{lex} Gamma2

Returns True if fxg (lexicographic product) defines a graph-theoretic mapping
from Gamma2xGamma4 to Gamma1xGamma3 that is harmonic, and False otherwise. 

The lex product is not symmetric.

Assumes Gamma1, Gamma2, Gamma3, Gamma4 are in standard form.
\end{verbatim}

\textbf{Example:}
\begin{verbatim}
    EXAMPLES:
        sage: Gamma1 = graphs.PathGraph(2)                                                                        
        sage: Gamma2 = graphs.PathGraph(3)                                                                        
        sage: Gamma3 = graphs.CycleGraph(4)                                                                       
        sage: Gamma4 = graphs.CycleGraph(4)                                                                       
        sage: V2 = Gamma2.vertices(sort=True)                                                                              
        sage: V4 = Gamma4.vertices(sort=True)                                                                              
        sage: f0 = [0, 1, 1]                                                                                      
        sage: g0 = [0,1,2,3]                                                                                      
        sage: f = [V2, f0]                                                                                        
        sage: g = [V4, g0]                                                                                        
        sage: is_harmonic_graph_morphism(Gamma1, Gamma2, f)                                                       
        True
        sage: is_harmonic_graph_morphism(Gamma3, Gamma4, g)                                                       
        True
        sage: [horizonal_multiplicity_of_harmonic_map(Gamma1, Gamma2, f, v) for v in V2]                          
        [1, 1, 0]
        sage: [vertical_multiplicity_of_harmonic_map(Gamma1, Gamma2, f, v) for v in V2]                           
        [0, 1, 1]
        sage: [horizonal_multiplicity_of_harmonic_map(Gamma3, Gamma4, g, v) for v in V4]                          
        [1, 1, 1, 1]
        sage: [vertical_multiplicity_of_harmonic_map(Gamma3, Gamma4, g, v) for v in V4]                           
        [0, 0, 0, 0]
        sage: is_lexicographic_product_of_graph_morphisms_harmonic(Gamma1, Gamma2, f, Gamma3, Gamma4, g, verbose=False)  
        False
        sage: is_lexicographic_product_of_graph_morphisms_harmonic(Gamma1, Gamma2, f, Gamma3, Gamma4, g, verbose=True)   
        This [<vertex>, <intersection sizes>]] passes the check: [0, [1, 1, 1, 1, 1, 1]]
        This [<vertex>, <intersection sizes>]] passes the check: [1, [1, 1, 1, 1, 1, 1]]
        This [<vertex>, <intersection sizes>]] passes the check: [2, [1, 1, 1, 1, 1, 1]]
        This [<vertex>, <intersection sizes>]] passes the check: [3, [1, 1, 1, 1, 1, 1]]
        This [<vertex>, <intersection sizes>]] fails the check: [4, [1, 1, 1, 1, 2, 2]]
        This [<vertex>, <intersection sizes>]] fails the check: [5, [1, 1, 1, 1, 2, 2]]
        This [<vertex>, <intersection sizes>]] fails the check: [6, [1, 1, 1, 1, 2, 2]]
        This [<vertex>, <intersection sizes>]] fails the check: [7, [1, 1, 1, 1, 2, 2]]
        This [<vertex>, <intersection sizes>]] fails the check: [8, [0, 0, 0, 0, 2, 2]]
        This [<vertex>, <intersection sizes>]] fails the check: [9, [0, 0, 0, 0, 2, 2]]
        This [<vertex>, <intersection sizes>]] fails the check: [10, [0, 0, 0, 0, 2, 2]]
        This [<vertex>, <intersection sizes>]] fails the check: [11, [0, 0, 0, 0, 2, 2]]
        False
        sage: is_lexicographic_product_of_graph_morphisms_harmonic(Gamma3, Gamma4, g, Gamma3, Gamma4, g)                 
        True
        sage: Gamma1 = graphs.CycleGraph(3)                                                                       
        sage: Gamma2 = graphs.CycleGraph(6)                                                                       
        sage: V2 = Gamma2.vertices(sort=True)                                                                              
        sage: f0 = [0, 1, 2, 0, 1, 2]                                                                             
        sage: f = [V2, f0]                                                                                        
        sage: is_harmonic_graph_morphism(Gamma1, Gamma2, f)                                                       
        True
        sage: [horizonal_multiplicity_of_harmonic_map(Gamma1, Gamma2, f, v) for v in V2] 
        [1, 1, 1, 1, 1, 1]
        sage: [vertical_multiplicity_of_harmonic_map(Gamma1, Gamma2, f, v) for v in V2]   
        [0, 0, 0, 0, 0, 0]
        sage: Gamma3 = graphs.CycleGraph(4)                                                                       
        sage: Gamma4 = graphs.CycleGraph(8)                                                                       
        sage: V4 = Gamma4.vertices(sort=True)                                                                              
        sage: g0 = [0,1,2,3,0,1,2,3]                                                                              
        sage: g = [V4, g0]                                                                                        
        sage: is_harmonic_graph_morphism(Gamma3, Gamma4, g)                                                       
        True
        sage: [horizonal_multiplicity_of_harmonic_map(Gamma1, Gamma2, f, v) for v in V2]                          
        [1, 1, 1, 1, 1, 1]
        sage: [vertical_multiplicity_of_harmonic_map(Gamma1, Gamma2, f, v) for v in V2]                           
        [0, 0, 0, 0, 0, 0]
        sage: is_lexicographic_product_of_graph_morphisms_harmonic(Gamma3, Gamma4, g, Gamma3, Gamma4, g)                 
        False
        sage: is_lexicographic_product_of_graph_morphisms_harmonic(Gamma1, Gamma2, f, Gamma3, Gamma4, g)   
        False
        sage: Gamma1 = graphs.CompleteGraph(4)
        sage: V2 = [0, 1, 2, 3, 4, 5, 6, 7] 
        sage: E2 = [(0, 1), (0, 2), (0, 7), (1, 2), (1, 3), (2, 3), (3, 4), (4, 5), (4, 6), (5, 6), (5, 7), (6, 7)]
        sage: Gamma2 = Graph([V2,E2]) # d3reg8a
        sage: Gamma3 = graphs.WheelGraph(4)
        sage: V3 = Gamma3.vertices(sort=True)
        sage: g0 = [1,2,3,0]
        sage: g = [V3,g0]
        sage: f0 = [0,1,2,3,0,2,1,3]
        sage: f = [V2, f0]
        sage: is_harmonic_graph_morphism(Gamma1, Gamma2, f)
        True
        sage: is_harmonic_graph_morphism(Gamma3, Gamma1, g)
        True                   
        sage: is_lexicographic_product_of_graph_morphisms_harmonic(Gamma1, Gamma2, f, Gamma1, Gamma3, g, verbose = False)
        True 
\end{verbatim}

\section{\tt \ \ is\_lexicographic\_product\_projection1\_harmonic}

\textbf{Syntax:} \verb}is_lexicographic_product_projection1_harmonic(Gamma1, Gamma2, verbose=False)}\\

\textbf{Description:}
\begin{verbatim}
Returns 
       True, if pr_1 (proj of lexicographic product onto 1st factor) defines a 
                harmonic morphism from Gamma1xGamma2 to Gamma1, 
and 
       False, otherwise. 

Assumes Gamma1, Gamma2 are in standard form.
\end{verbatim}

\textbf{Example:}
\begin{verbatim}
    EXAMPLES:
        sage: Gamma1 = graphs.PathGraph(2)
        sage: Gamma2 = graphs.PathGraph(3)
        sage: is_lexicographic_product_projection1_harmonic(Gamma1, Gamma2) 
        True
        sage: is_lexicographic_product_projection1_harmonic(Gamma1, Gamma2, verbose=True) 
        This [<vertex>, <intersection sizes>]] passes the check: [0, [3]]
        This [<vertex>, <intersection sizes>]] passes the check: [1, [3]]
        This [<vertex>, <intersection sizes>]] passes the check: [2, [3]]
        This [<vertex>, <intersection sizes>]] passes the check: [3, [3]]
        This [<vertex>, <intersection sizes>]] passes the check: [4, [3]]
        This [<vertex>, <intersection sizes>]] passes the check: [5, [3]]
        True
\end{verbatim}

\section{\tt \ \ is\_lexicographic\_product\_projection2\_harmonic}

\textbf{Syntax:} \verb}is_lexicographic_product_projection2_harmonic(Gamma1, Gamma2, verbose=False)}\\

\textbf{Description:}
\begin{verbatim}
Returns
    True, if pr_2 (proj of lexicographic product onto 2nd factor) defines a 
            harmonic morphism from Gamma1xGamma2 to Gamma2,
and
    False, otherwise. 


Assumes Gamma1, Gamma2 are in standard form.
\end{verbatim}

\textbf{Example:}
\begin{verbatim}
    EXAMPLES:
        sage: Gamma1 = graphs.PathGraph(2)
        sage: Gamma2 = graphs.PathGraph(3)
        sage: is_lexicographic_product_projection2_harmonic(Gamma1, Gamma2) 
        True
        sage: is_lexicographic_product_projection2_harmonic(Gamma1, Gamma2, verbose=True) 
        This [<vertex>, <intersection sizes>]] passes the check: [0, [2]]
        This [<vertex>, <intersection sizes>]] passes the check: [1, [2, 2]]
        This [<vertex>, <intersection sizes>]] passes the check: [2, [2]]
        This [<vertex>, <intersection sizes>]] passes the check: [3, [2]]
        This [<vertex>, <intersection sizes>]] passes the check: [4, [2, 2]]
        This [<vertex>, <intersection sizes>]] passes the check: [5, [2]]
        True
\end{verbatim}

\section{\tt \ \ is\_disjunctive\_product\_of\_graph\_morphisms\_a\_morphism}

\textbf{Syntax:} \verb}is_disjunctive_product_of_graph_morphisms_a_morphism(Gamma1, Gamma2, f, Gamma3, Gamma4, g, verbose = False)}\\

\textbf{Description:}
\begin{verbatim}
Returns
      True, if fxg (disjunctive product map) defines a graph-theoretic mapping
               from Gamma2xGamma4 to Gamma1xGamma3 that is a morphism,
and
      False, otherwise. 

Assumes Gamma1, Gamma2, Gamma3, Gamma4 are in standard form.

Converts Gamma2xGamma4 and Gamma1xGamma3 to standard form (eg, for verbose=True output).
\end{verbatim}

\textbf{Example:}
\begin{verbatim}
    EXAMPLES:
        sage: Gamma1 = graphs.PathGraph(2)
        sage: Gamma2 = graphs.PathGraph(3)
        sage: f0 = [0,1,1]       
        sage: V2 = Gamma2.vertices(sort=True) 
        sage: f = [V2, f0]     
        sage: [vertical_multiplicity_of_harmonic_map(Gamma1, Gamma2, f, v) for v in V2] 
        [0, 1, 1]
        sage: [horizonal_multiplicity_of_harmonic_map(Gamma1, Gamma2, f, v) for v in V2]
        [1, 1, 0]
        sage: Gamma3 = graphs.PathGraph(2) 
        sage: Gamma4 = graphs.CycleGraph(4) 
        sage: V4 = Gamma4.vertices(sort=True) 
        sage: g0 = [0,1,1,0]     
        sage: g = [V4, g0]  
        sage: [vertical_multiplicity_of_harmonic_map(Gamma3, Gamma4, g, v) for v in V4] 
        [1, 1, 1, 1]
        sage: [horizonal_multiplicity_of_harmonic_map(Gamma3, Gamma4, g, v) for v in V4]  
        [1, 1, 1, 1]
        sage: is_disjunctive_product_of_graph_morphisms_a_morphism(Gamma1, Gamma2, f, Gamma3, Gamma4, g) 
        False
        sage: is_disjunctive_product_of_graph_morphisms_a_morphism(Gamma1, Gamma2, f, Gamma3, Gamma4, g, verbose = True) 
        The edge e =  (0, 3)  has no image in Gamma1 under f.
        False
        sage: Gamma1 = graphs.CompleteGraph(4)
        sage: V2 = [0, 1, 2, 3, 4, 5, 6, 7] 
        sage: E2 = [(0, 1), (0, 2), (0, 7), (1, 2), (1, 3), (2, 3), (3, 4), (4, 5), (4, 6), (5, 6), (5, 7), (6, 7)]
        sage: Gamma2 = Graph([V2,E2]) # d3reg8a
        sage: Gamma3 = graphs.WheelGraph(4)
        sage: V3 = Gamma3.vertices(sort=True)
        sage: g0 = [1,2,3,0]
        sage: g = [V3,g0]
        sage: f0 = [0,1,2,3,0,2,1,3]
        sage: f = [V2, f0]
        sage: is_harmonic_graph_morphism(Gamma1, Gamma2, f)
        True
        sage: is_harmonic_graph_morphism(Gamma3, Gamma1, g)
        True
        sage: is_disjunctive_product_of_graph_morphisms_a_morphism(Gamma1, Gamma2, f, Gamma1, Gamma3, g, verbose = False)
        True
\end{verbatim}

\section{\tt \ \ is\_disjunctive\_product\_projection1\_morphism}

\textbf{Syntax:} \verb}is_disjunctive_product_projection1_morphism(Gamma1, Gamma2, verbose=False)}\\

\textbf{Description:}
\begin{verbatim}
Returns True if pr_1 (proj of disjunctive product onto 1st factor) defines a 
graph morphism from Gamma1xGamma2 to Gamma1, and False otherwise. 

Assumes Gamma1, Gamma2 are in standard form.

Converts Gamma1xGamma2 to standard form (eg, for verbose=True output).
\end{verbatim}

\textbf{Example:}
\begin{verbatim}
    EXAMPLES:
        sage: Gamma1 = graphs.PathGraph(2)
        sage: Gamma2 = graphs.PathGraph(3)
        sage: is_disjunctive_product_projection1_morphism(Gamma1, Gamma2) 
        False
        sage: is_disjunctive_product_projection1_morphism(Gamma1, Gamma2, verbose=True) 
        The edge e =  (0, 1)  has no image in Gamma1 under f.
        False
\end{verbatim}

\section{\tt \ \ is\_disjunctive\_product\_projection2\_morphism}

\textbf{Syntax:} \verb}is_disjunctive_product_projection2_morphism(Gamma1, Gamma2, verbose=False)}\\

\textbf{Description:}
\begin{verbatim}
Returns
        True, if pr_2 (proj of disjunctive product onto 2nd factor) defines a 
                 graph morphism from Gamma1xGamma2 to Gamma2,
and
        False, otherwise. 

Assumes Gamma1, Gamma2 are in standard form.

Converts Gamma1xGamma2 to standard form (eg, for verbose=True output).
\end{verbatim}

\textbf{Example:}
\begin{verbatim}
    EXAMPLES:
        sage: Gamma1 = graphs.PathGraph(2)
        sage: Gamma2 = graphs.PathGraph(3)
        sage: is_disjunctive_product_projection2_morphism(Gamma1, Gamma2) 
        False
        sage: is_disjunctive_product_projection2_morphism(Gamma1, Gamma2, verbose=True) 
        The edge e =  (0, 3)  has no image in Gamma1 under f.
        False
\end{verbatim}

\section{\tt \ \ is\_disjunctive\_product\_projection1\_harmonic}

\textbf{Syntax:} \verb}is_disjunctive_product_projection1_harmonic(Gamma1, Gamma2, verbose=False)}\\

\textbf{Description:}
\begin{verbatim}
Returns
        True, if pr_1 (proj of disjunctive product onto 1st factor) defines a 
                 harmonic morphism from Gamma1xGamma2 to Gamma1,
and
        False. otherwise. 

NOTE:  Assumes Gamma1, Gamma2 are in standard form.

Converts Gamma1xGamma2 to standard form (eg, for verbose=True output).
\end{verbatim}

\textbf{Example:}
\begin{verbatim}
    EXAMPLES:
        sage: Gamma1 = graphs.PathGraph(2)
        sage: Gamma2 = graphs.PathGraph(3)
        sage: is_disjunctive_product_projection1_harmonic(Gamma1, Gamma2) 
        False
        sage: is_disjunctive_product_projection1_harmonic(Gamma1, Gamma2, verbose=True) 
        False
\end{verbatim}

\section{\tt \ \ is\_disjunctive\_product\_projection2\_harmonic}

\textbf{Syntax:} \verb}is_disjunctive_product_projection2_harmonic(Gamma1, Gamma2, verbose=False)}\\

\textbf{Description:}
\begin{verbatim}
Returns
        True, if pr_2 (proj of disjunctive product onto 2nd factor) defines a 
                 harmonic morphism from Gamma1xGamma2 to Gamma2,
and
        False. otherwise. 

NOTE:  Assumes Gamma1, Gamma2 are in standard form.

Converts Gamma1xGamma2 to standard form (eg, for verbose=True output).
\end{verbatim}

\textbf{Example:}
\begin{verbatim}
    EXAMPLES:
        sage: Gamma1 = graphs.PathGraph(2)
        sage: Gamma2 = graphs.PathGraph(3)
        sage: is_disjunctive_product_projection2_harmonic(Gamma1, Gamma2) 
        False
        sage: is_disjunctive_product_projection2_harmonic(Gamma1, Gamma2, verbose=True) 
        False
\end{verbatim}

\section{\tt \ \ morphism\_of\_line\_graphs}

\textbf{Syntax:} \verb}morphism_of_line_graphs(Gamma1, Gamma2, f, verbose=False)}\\

\textbf{Description:}
\begin{verbatim}

Returns L(f), the line graph morphism
    L(f): L(Gamma2) -> L(Gamma1) 
associated to the graph morphism 
    f: Gamma2 -> Gamma1
given as a color list.

Argument verbose=True to returns (L(f),  L(Gamma1), L(Gamma2)).

Internally, converts both L(Gamma2), L(Gamma1) to standard form. 

NOTE:    Gamma1, Gamma2 must be in standard form.
\end{verbatim}

\textbf{Example:}
\begin{verbatim}
EXAMPLES:
        sage: Gamma1 = graphs.PathGraph(2)
        sage: Gamma2 = graphs.PathGraph(3) 
        sage: f0 = [0,1,1]  
        sage: V2 = Gamma2.vertices(sort=True) 
        sage: f = [V2, f0]    
        sage: morphism_of_line_graphs(Gamma1, Gamma2, f, verbose=True)  
        ([[0, 1], [0, 0]], Graph on 1 vertex, Graph on 2 vertices)
        sage: morphism_of_line_graphs(Gamma1, Gamma2, f)  
        [[0, 1], [0, 0]]
        sage: Gamma1 = graphs.PathGraph(3)   
        sage: Gamma2 = graphs.CycleGraph(4)  
        sage: f0 = [0, 1, 2, 1]    
        sage: V2 = Gamma2.vertices(sort=True)   
        sage: f = [V2, f0]          
        sage: is_harmonic_graph_morphism(Gamma1, Gamma2, f)
        True
        sage: morphism_of_line_graphs(Gamma1, Gamma2, f, verbose=True) 
        ([[0, 1, 2, 3], [0, 0, 1, 0]], Graph on 2 vertices, Graph on 4 vertices)
        sage: Gamma1 = graphs.CycleGraph(3) 
        sage: Gamma2 = graphs.CycleGraph(6) 
        sage: f0 = [0, 1, 2, 0, 1, 2] 
        sage: V2 = Gamma2.vertices(sort=True)  
        sage: f = [V2, f0]          
        sage: is_harmonic_graph_morphism(Gamma1, Gamma2, f) 
        True
        sage: [horizonal_multiplicity_of_harmonic_map(Gamma1, Gamma2, f, v) for v in V2] 
        [1, 1, 1, 1, 1, 1]
        sage: [vertical_multiplicity_of_harmonic_map(Gamma1, Gamma2, f, v) for v in V2]  
        [0, 0, 0, 0, 0, 0]
        sage: morphism_of_line_graphs(Gamma1, Gamma2, f, verbose=True) 
        ([[0, 1, 2, 3, 4, 5], [0, 1, 2, 0, 0, 2]],
         Graph on 3 vertices,
         Graph on 6 vertices)
        sage: ans = morphism_of_line_graphs(Gamma1, Gamma2, f, verbose=True) 
        sage: Lf = ans[0] 
        sage: LGamma2_st = ans[2]  
        sage: LGamma1_st = ans[1] 
        sage: is_graph_morphism(LGamma1_st, LGamma2_st, Lf) 
        True
        sage: is_harmonic_graph_morphism(LGamma1_st, LGamma2_st, Lf) 
        True
        sage: is_harmonic_graph_morphism(LGamma1_st, LGamma2_st, Lf, verbose=True) 
        Gamma1 = (V1, E1) =  [0, 1, 2] [(0, 1), (0, 2), (1, 2)]
        Gamma2 = (V2, E2) =  [0, 1, 2, 3, 4, 5] [(0, 1), (0, 2), (1, 5), (2, 3), (3, 4), (4, 5)]
        f = [V2, f0], where f0 =  [0, 1, 2, 1, 0, 2]
        This [<vertex>, <intersection sizes>]] passes the check: [0, [1, 1]]
        This [<vertex>, <intersection sizes>]] passes the check: [1, [1, 1]]
        This [<vertex>, <intersection sizes>]] passes the check: [2, [1, 1]]
        This [<vertex>, <intersection sizes>]] passes the check: [3, [1, 1]]
        This [<vertex>, <intersection sizes>]] passes the check: [4, [1, 1]]
        This [<vertex>, <intersection sizes>]] passes the check: [5, [1, 1]]
        True
        sage: VLG2_st = LGamma2_st.vertices(sort=True) 
        sage: [horizonal_multiplicity_of_harmonic_map(LGamma1_st, LGamma2_st, Lf, v) for v in VLG2_st]             
        [1, 1, 1, 1, 1, 1]
        sage: [vertical_multiplicity_of_harmonic_map(LGamma1_st, LGamma2_st, Lf, v) for v in VLG2_st]              
        [0, 0, 0, 0, 0, 0]
\end{verbatim}

\section{\tt \ \ LC\_multipartite\_graph}

\textbf{Syntax:} \verb}LC_multipartite_graph(L, cyclic=False)}\\

This remarkably interesting class of graphs ends up (thanks to
an observation of C. Melles) to be basically the same as a blowup graph,
at least in many interesting cases.

See also \verb}blowup_graph(Gamma, L)}.

\textbf{Description:}
\begin{verbatim}
The function returns what we call a "locally complete multipartite graph",
the complete multipartite graph but with all "long" edges removed:
The vertices are
    V = {0, ..., n-1},
where V = \cup_{j=0}^{k-1} I_j is a disjoint partition,
where n_j = |I_j| (so that n = n_1 + ... + n_{k-1}).
The edges are
    (u,v) for u in I_{j_1} and v in I_{j_2},
where
    j_1 \not= j_2, and j_2 = j_1 \pm 1. 
Here L = [n_1, ... , n_{k-1}]. Each set I_j is
a set of consecutive integers in V:
    I_0 = [0, ..., n_0-1], ..., [n-n_{k-1}, ..., n-1].
Each I_j i an independent set of vertices in the graph.

The option "cyclic = True" constructs the same matrix but where
there are also edges between I_0 and I_{k-1} (ie, the
index set {0, ..., k-1} is identified with ZZ/kZZ).
\end{verbatim}

\textbf{Example:}
\begin{verbatim}
    EXAMPLES:
        sage: Gamma = LC_multipartite_graph([2,3,4])                                       
        sage: Gamma.vertices(sort=True)                                                             
        [0, 1, 2, 3, 4, 5, 6, 7, 8]
        sage: Gamma.edges(labels=None,sort=True)                                                     
        [(0, 2), (0, 3), (0, 4), (1, 2), (1, 3), (1, 4), 
         (2, 5), (2, 6), (2, 7), (2, 8), (3, 5), (3, 6), (3, 7), (3, 8), (4, 5), (4, 6), (4, 7), (4, 8)]
        sage: 
        sage: Gamma2 = LC_multipartite_graph([2,3,2])  
        sage: V2 = Gamma2.vertices(sort=True)   
        sage: Gamma1 = graphs.PathGraph(3) 
        sage: f0 = [0,0,1,1,1,2,2] 
        sage: f = [V2, f0]       
        sage: is_harmonic_graph_morphism(Gamma1, Gamma2, f) 
        True
        sage: [horizontal_multiplicity_of_harmonic_map(Gamma1, Gamma2, f, v) for v in V2] 
        [3, 3, 2, 2, 2, 3, 3]
        sage: [vertical_multiplicity_of_harmonic_map(Gamma1, Gamma2, f, v) for v in V2]  
        [0, 0, 0, 0, 0, 0, 0]
        sage: 
        sage: Gamma2 = LC_multipartite_graph([2,2,3]) 
        sage: V2 = Gamma2.vertices(sort=True)  
        sage: Gamma1 = graphs.PathGraph(3) 
        sage: f0 = [0,0,1,1,2,2,2] 
        sage: f = [V2, f0]        
        sage: is_harmonic_graph_morphism(Gamma1, Gamma2, f) 
        False
        sage: is_graph_morphism(Gamma1, Gamma2, f) 
        True
        sage: [vertical_multiplicity_of_harmonic_map(Gamma1, Gamma2, f, v) for v in V2]
        [0, 0, 0, 0, 0, 0, 0]
        sage: 
        sage: Gamma2 = LC_multipartite_graph([2,2,2], cyclic = True)
        sage: V2 = Gamma2.vertices(sort=True) 
        sage: Gamma1 = graphs.CycleGraph(3)
        sage: f0 = [0,0,1,1,2,2]
        sage: f = [V2, f0] 
        sage: is_harmonic_graph_morphism(Gamma1, Gamma2, f)
        True
        sage: [horizontal_multiplicity_of_harmonic_map(Gamma1, Gamma2, f, v) for v in V2] 
        [2, 2, 2, 2, 2, 2]
        sage: [vertical_multiplicity_of_harmonic_map(Gamma1, Gamma2, f, v) for v in V2] 
        [0, 0, 0, 0, 0, 0]
        sage: 
        sage: Gamma2 = LC_multipartite_graph([2,1,2,1], cyclic = True)  
        sage: V2 = Gamma2.vertices(sort=True)  
        sage: Gamma1 = graphs.CycleGraph(4) 
        sage: f0 = [0,0,1,2,2,3] 
        sage: f = [V2, f0] 
        sage: is_harmonic_graph_morphism(Gamma1, Gamma2, f) 
        True
        sage: [horizontal_multiplicity_of_harmonic_map(Gamma1, Gamma2, f, v) for v in V2] 
        [1, 1, 2, 1, 1, 2]
        sage: [vertical_multiplicity_of_harmonic_map(Gamma1, Gamma2, f, v) for v in V2]
        [0, 0, 0, 0, 0, 0]
\end{verbatim}

\section{\tt \ \ blowup\_graph}

\textbf{Syntax:} \verb}blowup_graph(Gamma, L)}\\

\textbf{Description:}
\begin{verbatim}
This function generalizes
        blowup_balanced_graph(Gamma, r)
to a vector-valued argument r. It also generalizes
        LC_multipartite_graph
to "base" graphs other than the path graph or the cycle graph.
 
INPUT:
  - Gamma is a graph with k vertices and 
  - L = [n_1, ... , n_{k-1}] is a list of integers > 0. 
OUTPUT:
    Gamma2 (defined below).

Let
           n = n_0 + ... + n_{k-1}.
Each set I_j is a set of consecutive integers in V2:
     I_0 = [0, ..., n_0-1], ..., I_{k-1} = [n-n_{k-1}, ..., n-1].
Each I_j is an independent set of vertices in the graph.
Let
         V2 = \cup_{x in V1} I_x.
Connect a vertex v in I_x \subset V2 to each vertex w in I_y by an
edge
                  (v,w) in E2
if and only if (x,y) in E1.

This defines Gamma2 = (V2, E2).
\end{verbatim}

\textbf{Example:}
\begin{verbatim}
    EXAMPLES:
        sage: Gamma1 = graphs.CycleGraph(4)      
        sage: L = [2,3,2,3]                         
        sage: Gamma2a = LC_multipartite_graph(L)    
        sage: Gamma2b = blowup_graph(Gamma1, L)              
        sage: Gamma2a == Gamma2b         
        False
        sage: Gamma2a = LC_multipartite_graph(L, cyclic = True)
        sage: Gamma2a == Gamma2b 
        True
        sage: f0 = [0,0,1,1,1,2,2,3,3,3]   
        sage: V2b = Gamma2b.vertices(sort=True)         
        sage: f = [V2b, f0] 
        sage: is_harmonic_graph_morphism(Gamma1, Gamma2b, f) 
        False
        sage: is_graph_morphism(Gamma1, Gamma2b, f)  
        True
        sage: Gamma1 = graphs.DiamondGraph()      
        sage: L = [3,2,3,2]     
        sage: f0 = [0,0,0,1,1,2,2,2,3,3] 
        sage: Gamma2 = blowup_graph(Gamma1, L)                                   
        sage: V2 = Gamma2.vertices(sort=True) 
        sage: f = [V2, f0]    
        sage: is_harmonic_graph_morphism(Gamma1, Gamma2, f) 
        False
        sage: is_graph_morphism(Gamma1, Gamma2, f)  
        True
        sage: Gamma1 = graphs.DiamondGraph() 
        sage: L = [2,2,2,2]      
        sage: f0 = [0,0,1,1,2,2,3,3]   
        sage: Gamma2 = blowup_graph(Gamma1, L)                                   
        sage: V2 = Gamma2.vertices(sort=True)    
        sage: f = [V2, f0]   
        sage: is_harmonic_graph_morphism(Gamma1, Gamma2, f) 
        True
        sage: L = [3,3,3,3]   
        sage: Gamma2 = blowup_graph(Gamma1, L)                                   
        sage: f0 = [0,0,0,1,1,1,2,2,2,3,3,3] 
        sage: V2 = Gamma2.vertices(sort=True) 
        sage: f = [V2, f0]  
        sage: is_harmonic_graph_morphism(Gamma1, Gamma2, f) 
        True
\end{verbatim}

\section{\tt \ \ blowdown\_morphism}

\textbf{Syntax:} \verb}blowdown_morphism(Gamma, L)}\\

This could also be called the ``projection morphism for the blowup graph.''

\textbf{Description:}
\begin{verbatim}
This function constructs the graph morphism
    phi : blowup_graph(Gamma, L) --> Gamma
While blowup_graph(Gamma, L) replaces the i-th vertex of Gamma by 
"copies" indexed by the i-th set L[i], this blowdown_morphism 
does the opposite: 
    phi sends any of these copies to the original vertex. 

Here Gamma is a graph with k vertices and 
    L = [n_0, ... , n_{k-1}]
is a list of integers > 0. Let
    n = n_0 + ... + n_{k-1}.
The output is a graph morphism in CL notation.

\end{verbatim}

\textbf{Example:}
\begin{verbatim}
   EXAMPLES:
        sage: Gamma = graphs.CycleGraph(4)      
        sage: L = [2,3,2,3]                         
        sage: GammaL = blowup_graph(Gamma, L)              
        sage: blowdown_morphism(Gamma, L)
        [0, 0, 1, 1, 1, 2, 2, 3, 3, 3]
\end{verbatim}

\section{\tt \ \ degree\_of\_harmonic\_morphism}

\textbf{Syntax:} \verb}degree_of_harmonic_morphism(Gamma1, Gamma2, f)}\\

\textbf{Description:}

Returns the degree of the map phi, computed as the total
    horizontal multiplicity of a fiber.

    Assumes both Gamma1 and Gamma2 are in standard form.

\textbf{Example:}
\begin{verbatim}
EXAMPLES:
        sage: Gamma = graphs.CycleGraph(4)
        sage: L = [2,3,2,3]
        sage: GammaL = blowup_graph(Gamma, L)
        sage: VL = GammaL.vertices(sort=True)
        sage: f0 = blowdown_morphism(Gamma, L)
        sage: f = [VL, f0]
        sage: degree_of_harmonic_morphism(Gamma, GammaL, f)
        6
\end{verbatim}

\section{\tt \ \ subdivision\_graph}

\textbf{Syntax:} \verb}subdivision_graph(Gamma1, e)}\\

\textbf{Description:}
\begin{verbatim}
INPUT:
    - Gamma1, a graph in standard form.
OUTPUT:
    A pair (Gamma2, f), where
                  Gamma2
    is the subdivision graph of Gamma1 obtained by subdividing the edge e of Gamma1, 
                    f
    is the CL notation for the graph morphism
            phi: Gamma2 -> Gamma1,
    where the new vertex of Gamma2 has been sent to the "tail"
    (smallest) vertex incident to e. 
    
    This leaves Gamma1 itself unchanged. 
\end{verbatim}

\textbf{Example:}
\begin{verbatim}
EXAMPLES:
        sage: Gamma = graphs.CycleGraph(3)
        sage: e = (0,1)
        sage: Gamma2, f = subdivision_graph(Gamma, e)
        [[0, 1, 2, 3], [0, 1, 2, 0]]
        sage: Gamma2.is_isomorphic(graphs.CycleGraph(4))
        True
        sage: is_harmonic_graph_morphism(Gamma, Gamma2, f, verbose=False)
        False
        sage: is_graph_morphism(Gamma, Gamma2, f, verbose=False)
        True
        sage: Gamma = graphs.PathGraph(3)
        sage: e = (0,1)
        sage: Gamma2, f = subdivision_graph(Gamma, e)
        [[0, 1, 2, 3], [0, 1, 2, 0]]
        sage: is_harmonic_graph_morphism(Gamma, Gamma2, f, verbose=False)
        True
        sage: Gamma2.is_isomorphic(graphs.PathGraph(4))
        True
\end{verbatim}

\section{\tt \ \ lower\_upper\_bound\_sgon}

\textbf{Syntax:} \verb}lower_upper_bound_sgon(Gamma)}\\

\textbf{Description:}
\begin{verbatim}
Returns the upper UB and lower bound LB on the sgonality of a graph, so
              LB <= sgon(Gamma) <= UB.
For more on the bounds implemented here, see pp 2-3 of [1].

REFERENCES
        [1] Gijswijt, Smit, van der Wegen, "Computing graph gonality is hard", 
            ArXiV preprint, 2019.

\end{verbatim}

\textbf{Example:}
\begin{verbatim}
EXAMPLES:
        sage: Gamma = graphs.PathGraph(4)
        sage: lower_upper_bound_sgon(Gamma)
        0.19 1.5
        sage: Gamma = graphs.CycleGraph(5)
        sage: lower_upper_bound_sgon(Gamma)
        0.41 2.0
        sage: Gamma = graphs.DiamondGraph()
        sage: lower_upper_bound_sgon(Gamma)
        0.44 2.5
        sage: Gamma = graphs.CubeGraph(3)
        sage: lower_upper_bound_sgon(Gamma)
        0.44 4.0
        sage: Gamma = flower_graph([4,5])
        sage: lower_upper_bound_sgon(Gamma)
        0.18 2.5


\end{verbatim}

\section{\tt \ \ subdivision\_graph\_morphism}

\textbf{Syntax:} \verb}subdivision_graph_morphism(Gamma1, Gamma2, f, s)}\\

\textbf{Description:}
\begin{verbatim}

INPUT:
       - Gamma1, Gamma2  are two graphs in standard form,
       - s is an edge of Gamma1,
       - f is a surjective morphism from Gamma2 -> Gamma1, in CL notation.

OUTPUT:
       - (Gamma3, Gamma4, g), where
       - Gamma3 is the subdivision graph of Gamma1 obtained by subdividing the edge s of Gamma1, 
       - Gamma4 is the subdivision graph of Gamma2 obtained by subdividing the edges e of Gamma2
                in the fiber of s, 
       - g is the graph morphism from Gamma4 -> Gamma3, corresponding to f.
    
\end{verbatim}

This leaves Gamma1, Gamma2, f unchanged.

\textbf{Example:}
\begin{verbatim}
EXAMPLES:
        sage: Gamma1 = graphs.PathGraph(3)
        sage: Gamma2 = graphs.PathGraph(4)
        sage: f21_0 = [0,1,2,2]
        sage: V2 = Gamma2.vertices(sort=True)
        sage: f21 = [V2, f21_0]
        sage: is_harmonic_graph_morphism(Gamma1, Gamma2, f21, verbose=False)
        True
        sage: E1 = Gamma1.edges(labels=None,sort=True)
        sage: s = E1[0]; s
        (0, 1)
        sage: Gamma3, Gamma4, f43 = subdivision_graph_morphism(Gamma1, Gamma2, f21, s)
        sage: f43
        [[0, 1, 2, 3, 4], [0, 1, 2, 2, 3]]
        sage: is_harmonic_graph_morphism(Gamma3, Gamma4, f43, verbose=False)
        True
\end{verbatim}

\section{\tt \ \ is\_smooth\_vertex}

\textbf{Syntax:} \verb}is_smooth_vertex(Gamma, v0)}\\

\textbf{Description:}
The function checks if vertex v0
 (a) has exactly 2 neighbors, say u and w, and
 (b) those two neighbors are not adjacent in Gamma.
If these conditions hold, v0 is smooth.

Returns True if v0 is smooth and False otherwise.
    Assumes Gamma is a simple graph.

\textbf{Example:}
\begin{verbatim}
EXAMPLES:
        sage: Gamma = graphs.PathGraph(4)
        sage: V = Gamma.vertices(sort=True); V
         [0, 1, 2, 3]
        sage: v0 = V[0]; is_smooth_vertex(Gamma, v0)
         False
        sage: v0 = V[1]; is_smooth_vertex(Gamma, v0)
         True
\end{verbatim}

\section{\tt \ \ smooth\_graph}

\textbf{Syntax:} \verb}smooth_graph(Gamma, v0)}\\

\textbf{Description:}
\begin{verbatim}
INPUT:
    - Gamma a graph in standard form,
    - v0 a smooth vertex of Gamma,

OUTPUT:
    the smoothing of Gamma at v0, in standard form.
    (The labeling of vertices changes by relabeling the vertices
    v0+1, ..., n-1 to v0, ..., n-2.)

If v0 is not smooth, it returns Gamma.
\end{verbatim}

\textbf{Example:}
\begin{verbatim}
EXAMPLES:
        sage: V1 = [0,1,2,3,4]
        sage: E1 = [(0,1),(0,2),(1,3),(2,3),(2,4),(3,4)]
        sage: Gamma1 = Graph([V1,E1])  # 4 is the "tip of the roof", (0,1) is the "base"
        sage: Gamma2 = smooth_graph(Gamma1, 4)
         v0 is not a smooth vertex of Gamma.
        sage: Gamma2 == Gamma1
         True
        sage: Gamma2 = smooth_graph(Gamma1, 0)
        sage: Gamma2.is_isomorphic(graphs.DiamondGraph())
         True
\end{verbatim}

\section{\tt \ \ smooth\_graph\_morphism}

\textbf{Syntax:} \verb}smooth_graph_morphism(Gamma1, Gamma2, f, x0)}\\

\textbf{Description:}
\begin{verbatim}
INPUT:
    -  Gamma1, Gamma2 graphs in standard form,
    -  f is a morphism from Gamma2 -> Gamma1,  in CL notation.
    -  x0 is a vertex of Gamma1, 
OUTPUT:
       - (Gamma3, Gamma4, g), where
       - Gamma3 is the graph obtained by smoothing Gamma1 at x0, 
       - Gamma4 is the graph obtained by smoothing Gamma2 at the 
                vertices in the fiber over x0, 
       - g is the graph morphism from Gamma4 -> Gamma3, corresponding to f.
    
This leaves Gamma1, Gamma2, f unchanged.

This assumes x0 is a smooth vertex of Gamma1 and that all the vertices in the
fiber over x0 are smooth. If not, it returns Gamma1, Gamma2, f and prints a message.

In some examples, this construction can produce a non-simple graph
(with repeated edges). In this case, the graph returned by this function is the
``simplified'' version.

\end{verbatim}

\textbf{Example:}
\begin{verbatim}
EXAMPLES:
        sage: Gamma1 = graphs.CycleGraph(4)
        sage: Gamma2 = graphs.CycleGraph(8)
        sage: f21_0 = [0, 1, 2, 3, 0, 1, 2, 3]
        sage: V2 = Gamma2.vertices(sort=True)
        sage: f21 = [V2, f21_0]
        sage: is_harmonic_graph_morphism(Gamma1, Gamma2, f21, verbose=False)
         True
        sage: Gamma3, Gamma4, f43 = smooth_graph_morphism(Gamma1, Gamma2, f21, 1)
        sage: is_harmonic_graph_morphism(Gamma3, Gamma4, f43, verbose=False)
         True
\end{verbatim}

\section{\tt \ \ delete\_vertex\_graph\_morphism}

\textbf{Syntax:} \verb}delete_vertex_graph_morphism(Gamma1, Gamma2, f, x0)}\\

\textbf{Description:}
\begin{verbatim}
INPUT:
    -  Gamma1, Gamma2 graphs in standard form,
    -  f is a morphism from Gamma2 -> Gamma1,  in CL notation.
    -  x0 is a vertex of Gamma1, 
OUTPUT:
       - (Gamma3, Gamma4, g), where
       - Gamma3 is the graph obtained by deleting x0 from Gamma1,
       - Gamma4 is the graph obtained by removing from Gamma2 all
                vertices in the fiber over x0, 
       - g is the graph morphism from Gamma4 -> Gamma3, corresponding to  
           the restriction of f.
\end{verbatim}
    
    This leaves Gamma1, Gamma2, f unchanged and assumes Gamma1, Gamma2 are in standard form.
    If the result of a deletion leaves a disconnected graph then it returns 
    Gamma1, Gamma2, f and prints a message.

\textbf{Example:}
\begin{verbatim}
EXAMPLES:
        sage: Gamma1 = graphs.CycleGraph(3)
        sage: Gamma2 = graphs.ButterflyGraph()
        sage: V2 = Gamma2.vertices(sort=True)
        sage: f21_0 = [1,2,1,2,0]
        sage: f21 = [V2, f21_0]
        sage: is_harmonic_graph_morphism(Gamma1, Gamma2, f21, verbose=False)
        True
        sage: Gamma3, Gamma4, f43 = delete_vertex_graph_morphism(Gamma1, Gamma2, f21, 1)
        sage: is_harmonic_graph_morphism(Gamma3, Gamma4, f43, verbose=False)
        True
        sage: V3 = Gamma3.vertices(sort=True); V3
        [0, 1]
        sage: V4 = Gamma4.vertices(sort=True); V4
        [0, 1, 2]
        sage: f43
        [[0, 1, 2], [1, 1, 0]]
        sage: Gamma3, Gamma4, f43 = delete_vertex_graph_morphism(Gamma1, Gamma2, f21, 0)
        preimage of x0 under f must not contain a disconnecting vertex of Gamma2.
        sage: Gamma3==Gamma1
        True
        sage: Gamma4==Gamma2
        True
        sage: f43==f21
        True
\end{verbatim}

\section{\tt \ \ table\_of\_local\_horizontal\_multiplicities}

\textbf{Syntax:} \verb}table_of_local_horizontal_multiplicities(Gamma1, Gamma2, f)}\\

\textbf{Description:}

Self-explanatory.

\textbf{Example:}
\begin{verbatim}
EXAMPLES:
        sage: Gamma1 = graphs.PathGraph(2)
        sage: Gamma2 = graphs.PathGraph(4)
        sage: V2 = Gamma2.vertices(sort=True)
        sage: f21_0 = [0,1,1,0]
        sage: f21 = [V2, f21_0]
        sage: is_harmonic_graph_morphism(Gamma1, Gamma2, f21, verbose=True)
        Gamma1 = (V1, E1) =  [0, 1] [(0, 1)]
        Gamma2 = (V2, E2) =  [0, 1, 2, 3] [(0, 1), (1, 2), (2, 3)]
        f = [V2, f0], where f0 =  [0, 1, 1, 0]
        The edge e =  (1, 2)  is vertical under f.
        This [<vertex>, <intersection sizes>]] passes the check: [0, [1]]
        This [<vertex>, <intersection sizes>]] passes the check: [1, [1]]
        This [<vertex>, <intersection sizes>]] passes the check: [2, [1]]
        This [<vertex>, <intersection sizes>]] passes the check: [3, [1]]
        True
        sage: table_of_local_horizontal_multiplicities(Gamma1, Gamma2, f21)
        [[(0, (0, 1), 1)], [(1, (0, 1), 1)], [(2, (0, 1), 1)], [(3, (0, 1), 1)]]
\end{verbatim}

\section{\tt \ \ is\_harmonic\_graphWB\_morphism}

\textbf{Syntax:} \verb}is_harmonic_graphWB_morphism(Gamma1, Gamma2, f, B1, B2, verbose = False)}\\

\textbf{Description:}
\begin{verbatim}
Returns 
      True if f defines a graph-theoretic mapping
           from a graph Gamma2 with boundary B2 to another graph Gamma1
           with boundary B1 that is a harmonic morphism on the interior, and 
      False, otherwise. 

Assumes 
    (a) Gamma1, Gamma2 are in standard form,
    (b) B2 = [s2, t2],
        B1 = [s1, t1], 
and 
    (c) f must send B2 to B1.

Suppose Gamma2 has n vertices. A morphism 
              f: Gamma2 -> Gamma1
is represented by a CL pair of lists [L2, L1], where L2 = V2 is the 
list of all n vertices of Gamma2, and L1 is the list of length n of 
the vertices in Gamma1 that form the corresponding image under
the map f.
\end{verbatim}

\textbf{Example:}
\begin{verbatim}
    EXAMPLES:
        sage: V1 = [0,1,2,3,4]
        sage: E1 = [(0,3),(0,1),(1,2),(2,4)] 
        sage: B1 = [3,4]   
        sage: Gamma1 = Graph([V1,E1]) # P3+s1+t1=P5 
        sage: V2 = [0,1,2,3,4,5]  
        sage: E2 = [(0,4),(0,1),(0,2),(1,3),(2,3),(3,5)] 
        sage: B2 = [4,5]  
        sage: Gamma2 = Graph([V2,E2]) # C4+s2+t2 
        sage: f0_21 = [0,1,1,2,3,4] 
        sage: f21_0 = [0,1,1,2,3,4] 
        sage: f21 = [V2, f21_0] 
        sage: is_harmonic_graphWB_morphism(Gamma1, Gamma2, f_21, B1, B2, verbose = True) 
         Gamma1 = (V1, E1) =  [0, 1, 2, 3, 4] [(0, 1), (0, 3), (1, 2), (2, 4)]
         Boundary of Gamma1 =  [3, 4]
         Gamma2 = (V2, E2) =  [0, 1, 2, 3, 4, 5] [(0, 1), (0, 2), (0, 4), (1, 3), (2, 3), (3, 5)]
         Boundary of Gamma2 =  [4, 5]
         interior of Gamma1: Gamma3 = (V3, E3) =  [0, 1, 2] [(0, 1), (1, 2)]
         interior of Gamma2: Gamma4 = (V4, E4) =  [0, 1, 2, 3] [(0, 1), (0, 2), (1, 3), (2, 3)]
         Gamma1 = (V1, E1) =  [0, 1, 2, 3, 4] [(0, 1), (0, 3), (1, 2), (2, 4)]
         Gamma2 = (V2, E2) =  [0, 1, 2, 3, 4, 5] [(0, 1), (0, 2), (0, 4), (1, 3), (2, 3), (3, 5)]
         f = [V2, f0], where f0 =  [0, 1, 1, 2, 3, 4]
         Gamma1 = (V1, E1) =  [0, 1, 2] [(0, 1), (1, 2)]
         Gamma2 = (V2, E2) =  [0, 1, 2, 3] [(0, 1), (0, 2), (1, 3), (2, 3)]
         f = [V2, f0], where f0 =  [0, 1, 1, 2]
         This [<vertex>, <intersection sizes>]] passes the check: [0, [2]]
         This [<vertex>, <intersection sizes>]] passes the check: [1, [1, 1]]
         This [<vertex>, <intersection sizes>]] passes the check: [2, [1, 1]]
         This [<vertex>, <intersection sizes>]] passes the check: [3, [2]]
         True
        sage: is_harmonic_graphWB_morphism(Gamma1, Gamma2, f_21, B1, B2, verbose = False) 
         True
\end{verbatim}

\section{\tt \ \ partial\_join\_graph}

\textbf{Syntax:} \verb}partial_join_graph(Gamma, S)}\\

\textbf{Description:}
\begin{verbatim}
INPUT:
   - Gamma is a graph with n vertices (assumed to be in standard form),
   -   S   is a subset of vertices of Gamma.
OUTPUT:
    a graph obtained by adding a vertex n and edges (v, n), for v in S, to Gamma.

\end{verbatim}

\textbf{Example:}
\begin{verbatim}
EXAMPLES:
        sage: Gamma = graphs.CycleGraph(4)
        sage: S = [2,3]
        sage: Gamma2 = partial_join_graph(Gamma, S)
        sage: Gamma2.is_isomorphic(graphs.HouseGraph())
         True
        sage: Gamma = graphs.CycleGraph(4)
        sage: S = [0,1,2,3]
        sage: Gamma2 = partial_join_graph(Gamma, S)
        sage: Gamma2.is_isomorphic(graphs.WheelGraph(5))
         True
\end{verbatim}

\section{\tt \ \ leaf\_addition\_graph\_morphism}

\textbf{Syntax:} \verb}leaf_addition_graph_morphism(Gamma1, Gamma2, f, x0)}\\

\textbf{Description:}
\begin{verbatim}
INPUT:
   - Gamma1, Gamma2 graphs in standard form,  
   - x0 is a vertex of Gamma1 that is not the image of an edge under f, 
   - f is a harmonic morphism from Gamma2 -> Gamma1, in CL notation.
    
Assumes 
        * Gamma1, Gamma2 are in standard form, with V1 = [0,...,n1-1],
          V2 = [0,...,n2-1]. 
        * x0 is not the image of an edge.
        * f is harmonic.

OUTPUT:
     a triple (Gamma3, Gamma4, g), where
     - Gamma3 is the leaf-addition Gamma1 obtained by adding a leaf (x0, n1) to Gamma1, 
     - Gamma4 is the graph obtained by adding, for each v in f^(-1)(x0), 
                        mu(v) leafs attached to v. 
                 If f^(-1)(x0) = [v\_1,...,v\_k] (v\_1<v\_2<...<v\_k) then 
                 these new vertices and edges are labeled 
                 V21 = [n2,...,n2+mu(v\_1)-1], (v1,w) for w in V21,
                 V22 = [n2+mu(v\_1),..., n2+mu(v\_1)+mu(v\_2)-1], (v2,w) for w in V22,
                  ..., 
                 V2k = [n2+mu(v\_1)+...+mu(v\_{k-1}),..., n2+mu(v\_1)+...+mu(v\_k)-1], (vk,w) for w in V2k.
     - g is the graph morphism from Gamma4 -> Gamma3, corresponding to this extention of f.
\end{verbatim}
    
    This leaves Gamma1, Gamma2, f unchanged.

\textbf{Example:}
\begin{verbatim}
EXAMPLES:
        sage: Gamma2 = graphs.CycleGraph(4)
        sage: Gamma1 = graphs.PathGraph(3)
        sage: V2 = Gamma2.vertices(sort=True)
        sage: f21_0 = [0,1,2,1]
        sage: f21 = [V2, f21_0]
        sage: is_harmonic_graph_morphism(Gamma1, Gamma2, f21, verbose=False)
         True
        sage: [horizontal_multiplicity_of_harmonic_map(Gamma1, Gamma2, f21, v) for v in V2]
         [2, 1, 2, 1]
        sage: [vertical_multiplicity_of_harmonic_map(Gamma1, Gamma2, f21, v) for v in V2]
         [0, 0, 0, 0]
        sage: Gamma3,Gamma4,g = leaf_addition_graph_morphism(Gamma1, Gamma2, f21, 1)
        sage: is_harmonic_graph_morphism(Gamma3, Gamma4, g, verbose=False)
         True
        sage: [horizontal_multiplicity_of_harmonic_map(Gamma3, Gamma4, g, v) for v in V4]
         [2, 1, 2, 1, 1, 1]
        sage: [vertical_multiplicity_of_harmonic_map(Gamma3, Gamma4, g, v) for v in V4]
         [0, 0, 0, 0, 0, 0]
        sage: degree_of_harmonic_morphism(Gamma3, Gamma4, g)
         2
        sage: degree_of_harmonic_morphism(Gamma1, Gamma2, f21)
         2
        sage: is_harmonic_graph_morphism(Gamma3, Gamma4, g, verbose=False)
         True
        sage: [horizontal_multiplicity_of_harmonic_map(Gamma3, Gamma4, g, v) for v in V4]
         [2, 1, 2, 1, 1, 1]
        sage: [vertical_multiplicity_of_harmonic_map(Gamma3, Gamma4, g, v) for v in V4]
         [0, 0, 0, 0, 0, 0]
        sage: degree_of_harmonic_morphism(Gamma3, Gamma4, g)
         2
\end{verbatim}

\section{\tt \ \ sub\_graph4edge}

\textbf{Syntax:} \verb}sub_graph4edge(Gamma1, e0, Gamma, a, b)}\\

\textbf{Description:}
\begin{verbatim}
INPUT:
     * Gamma1, and Gamma are connected graphs in standard form,
     * e0 = (x0, y0) is an edge of Gamma1, x0<y0.
     * a,b are vertices in Gamma, not necessarily adjacent, a<b.
OUTPUT:
     a pair 
         (Gamma2, f), 
     where
         Gamma2 is the graph obtained by replacing the edge e0=(x0,y0) 
            of Gamma1 by the graph Gamma, 
         where 
           a    is attached at x0 and 
           b    is attached at y0 (so in the new graph a=x0, b=y0).
\end{verbatim}
    
    This leaves Gamma1 itself unchanged. This should agree with the function
    subdivision in the case Gamma = P3 and a,b are the degree 1 vertices.

\textbf{Example:}
\begin{verbatim}
EXAMPLES:
        sage: Gamma1 = graphs.CycleGraph(3)    
        sage: Gamma = graphs.CycleGraph(3)     
        sage: Gamma2 = sub_graph4edge(Gamma1, (0,1), Gamma, 0, 1) 
        sage: Gamma2.is_isomorphic(graphs.DiamondGraph())
         True
        sage: Gamma1 = graphs.CycleGraph(3)    
        sage: Gamma = graphs.CycleGraph(4)     
        sage: Gamma2 = sub_graph4edge(Gamma1, (0,1), Gamma, 0, 1)
        sage: Gamma2.is_isomorphic(graphs.HouseGraph()) 
         True
        sage: Gamma1 = graphs.CycleGraph(3)    
        sage: Gamma = graphs.PathGraph(3)     
        sage: Gamma2 = sub_graph4edge(Gamma1, (1,2), Gamma, 0, 2) 
        sage: Gamma2.is_isomorphic(graphs.CycleGraph(4))  
         True
        sage: Gamma3,f = subdivision_graph(Gamma1, (1,2)) 
        sage: Gamma2.is_isomorphic(Gamma3)     
         True
\end{verbatim}

\section{\tt \ \ has\_harmonic\_quotient}

\textbf{Syntax:} \verb}has_harmonic_quotient(Gamma, verbose=False)}\\

\textbf{Description:}
\begin{verbatim}
This function returns the list (possibly empty) of harmonic graph quotient morphisms, 
                   f: Gamma -> Gamma/G,
where G is a subgroup of Aut(Gamma), where f is in CL format.

Input: Gamma - a graph in standard form,
Output: The harmonic quotient graph morphisms.
\end{verbatim}

\textbf{Example:}
\begin{verbatim}
    EXAMPLES:
        sage: Gamma2 = graphs.LCFGraph(12, [6, -2, 2], 4) # trunc tetrahedron, d3reg12-12? 
        sage: L = has_harmonic_quotient(Gamma2, verbose=True)  
         10  non-trivial harmonic quotient morphism(s) found.
        sage: f0 = L[3][0]; f0     
         [2, 1, 1, 2, 0, 0, 2, 1, 1, 2, 0, 0]
        sage: G0 = L[3][1]; G0.cardinality()            
         4
        sage: Gamma1 = L[3][2]   
        sage: Gamma1.is_isomorphic(graphs.CycleGraph(3))    
         True
        sage: V2 = Gamma2.vertices(sort=True); V2 
         [0, 1, 2, 3, 4, 5, 6, 7, 8, 9, 10, 11]
        sage: f = [V2, f0]   
        sage: is_harmonic_graph_morphism(Gamma1, Gamma2, f, verbose=False)   
         True
\end{verbatim}

\section{\tt \ \ double\_point\_construction}

\textbf{Syntax:} \verb}double_point_construction(Gamma)}\\

\textbf{Description:}

Gamma is a k-regular graph with n vertices, assumed to be
   in standard form. We also assume k<n-1. This adds another vertex, 
   which we connect to k+1 vertices of Gamma by an edge. 
   Call this Gamma2. Make a copy of Gamma2. Connect each 
   vertex of degree k on one copy to its corresponding vertex 
   on the other copy. This new graph is (k+1)-regular, with 
   2n+2 vertices.

\textbf{Example:}
\begin{verbatim}
EXAMPLES:
       sage: Gamma1 = graphs.CycleGraph(4)
       sage: Gamma1.is_regular(2)
        True
       sage: Gamma2 = double_point_construction(Gamma1)
       sage: Gamma2.is_regular(3)
        True
       sage: Gamma1 = graphs.CompleteBipartiteGraph(3,3)
       sage: Gamma1.is_regular(3)
        True
       sage: Gamma2 = double_point_construction(Gamma1)
       sage: Gamma2.is_regular(4)
        True
\end{verbatim}

\section{\tt \ \ graph\_id\_graph6}
\textbf{Syntax:} \verb}graph_id_graph6(Gamma, verbose=False)}\\

\textbf{Description:}
graph6 is a standardized encoding for small graphs used by SageMath and other software.
It turns a graph into a short ASCII string like 'Dbk'.
Each ``small'' graph has its own graph6 ID.

Returns the graph6 id of Gamma = (V,E). 
    If verbose then it also displays the table of all graphs with 
    |V| vertices and |E| edges.


\textbf{Example:}
\begin{verbatim}
EXAMPLES:
        sage: Gamma = graphs.HouseGraph()
        sage: graph_id_graph6(Gamma, verbose=False)
        'Dbk'
        sage: graph_id_graph6(Gamma, verbose=True)

         graphs with 5 vertices and 6 edges:

        Graph6               Aut Grp Size         Degree Sequence     
        ------------------------------------------------------------
        DB{                  2                    [1, 2, 2, 3, 4]     
        DFw                  12                   [2, 2, 2, 3, 3]     
        DJ[                  24                   [0, 3, 3, 3, 3]     
        DJk                  2                    [1, 2, 3, 3, 3]     
        DK{                  8                    [2, 2, 2, 2, 4]     
        Dbk                  2                    [2, 2, 2, 3, 3]     


        'Dbk'
\end{verbatim}

\section{\tt \ \ effective\_resistance\_matrix}

\textbf{Syntax:} \verb}effective_resistance_matrix(Gamma)}\\

Imagine a graph is like a network of wires — each edge is a wire with resistance,
and each vertex is a point where wires meet.
This function answers the question:
``If I connect a battery between two points in the network, what’s the total resistance between them?''
That value is called the effective resistance.
Now imagine doing this for every pair of vertices — this function gives you a matrix where the 
(i,j)-entry tells you the effective resistance between vertex i and vertex j.

\textbf{Description:}

Returns the effective resistance matrix, even when Gamma is
    edge-weighted.

    Gamma is assumed to be in standard form. The default weight of
    each edge is 1 for each edge.

\textbf{Example:}
\begin{verbatim}
EXAMPLES:
        sage: C3 = graphs.CycleGraph(3)
        sage: effective_resistance_matrix(C3) 
         [  0 2/3 2/3]
         [2/3   0 2/3]
         [2/3 2/3   0]
\end{verbatim}

\section{\tt \ \ intersection\_number\_of\_vertex\_subset}

\textbf{Syntax:} \verb}intersection_number_of_vertex_subset(Gamma, C, v)}\\

Think of C as a cell in a partition of the vertices of Gamma.

\textbf{Description:}
\begin{verbatim}
Returns the number of neighbors of v in C.

INPUT:
    - Gamma =(V, E) must be in standard form.
    -   C is represented as a sublist of V.
    -   v is an element of V.
\end{verbatim}

\textbf{Example:}
\begin{verbatim}
    EXAMPLES:
        sage: Gamma = graphs.PathGraph(5)
        sage: C = [1,2]
        sage: intersection_number_of_vertex_subset(Gamma, C, 0)
        1
\end{verbatim}

\section{\tt \ \ is\_equitable\_partition}

\textbf{Syntax:} \verb}is_equitable_partition(Gamma, P, verbose=False)}\\

\textbf{Description:}
\begin{verbatim}
INPUT:
    - Gamma is be in standard form.
    - P     is represented as a list of disjoint lists.

OUTPUT:
    True, if P defines an equitable partition of the vertices
            of Gamma,
    False, otherwise.

If, in addiiton, verbose=True then it also returns the intersection matrix
if P is equitable, and extra info if P is not. 

\end{verbatim}

\textbf{Example:}
\begin{verbatim}
    EXAMPLES:
        sage: Gamma = graphs.PathGraph(5)
        sage: P = [[0,4],[1,3],[2]]
        sage: is_equitable_partition(Gamma, P)
	True
        sage: is_equitable_partition(Gamma, P, verbose=True)
        (
               [0 1 0]
               [1 0 1]
         True, [0 2 0]
         )
        sage: Gamma2 = graphs.HouseGraph()   
        sage: V2 = Gamma2.vertices(sort=True); V2
         [0, 1, 2, 3, 4]
        sage: P = [[0,2],[1,3],[4]] 
        sage: is_equitable_partition(Gamma2, P, verbose=True)   
         Fails to be equitable because deg( [4] [0, 2] ) isn't defined.
         False
\end{verbatim}


\section{\tt \ \ associated\_partition\_to\_morphism}

\textbf{Syntax:} \verb}associated_partition_to_morphism(Gamma1, Gamma2, f)}\\

\textbf{Description:}
\begin{verbatim}
Each surjective graph morphism
             phi: Gamma2 = (V2, E2) -> Gamma1 = (V1, E1)
determines a partition
                        P = P_{phi}
by letting the cels of P be the fibers phi^{-1}(x), x in V1.

Returns the partition of the vertices of Gamma2 associated to f.

Assumes Gamma1, Gamma2 are in standard form and f is surjective.
\end{verbatim}

\textbf{Example:}
\begin{verbatim}
   EXAMPLES:
        sage: Gamma1 = graphs.CycleGraph(3) 
        sage: Gamma2 = graphs.HouseGraph()   
        sage: V2 = Gamma2.vertices(sort=True); V2  
         [0, 1, 2, 3, 4]
        sage: f0 = [0,1,0,1,2]     
        sage: f = [V2, f0]         
        sage: associated_partition_equitable(Gamma1, Gamma2, f) 
         [[0, 2], [1, 3], [4]]
\end{verbatim}

\section{\tt \ \ is\_associated\_partition\_equitable}

\textbf{Syntax:} \verb}is_associated_partition_equitable(Gamma1, Gamma2, f, verbose=False)}\\

\textbf{Description:}
\begin{verbatim}
Returns 
        True, if the fibers of f defines an equitable partition of the vertices
               of Gamma2. If, in addition, verbose=True then it also returns the
               intersection matrix. 
        False, otherwise (and, if verbose=True) then it prints info about the 
               cells where the equitable condition fails.

Gamma1, Gamma2 must be in standard form and f is surjective.
\end{verbatim}

\textbf{Example:}
\begin{verbatim}
    EXAMPLES:
        sage: Gamma1 = graphs.CycleGraph(3) 
        sage: Gamma2 = graphs.HouseGraph()      
        sage: V2 = Gamma2.vertices(sort=True); V2 
         [0, 1, 2, 3, 4]
        sage: f0 = [0,1,0,1,2]      
        sage: f = [V2, f0]          
        sage: is_associated_partition_equitable(Gamma1, Gamma2, f, verbose=False)   
         False
        sage: is_associated_partition_equitable(Gamma1, Gamma2, f, verbose=True)   
         Fails to be equitable because deg( [4] [0, 2] ) isn't defined.
         False
        sage: Gamma1 = graphs.PathGraph(3) 
        sage: Gamma2 = graphs.PathGraph(4) 
        sage: f0 = [0,1,2,2] 
        sage: V2 = Gamma2.vertices(sort=True) 
        sage: f = [V2, f0]                                                                                                    
        sage: is_associated_partition_equitable(Gamma1, Gamma2, f, verbose=True) 
         Fails to be equitable because deg( [1] [2, 3] ) isn't defined.
         False
        sage: is_harmonic_graph_morphism(Gamma1, Gamma2, f, verbose=False) 
         True
\end{verbatim}

\section{\tt \ \ partition\_cell\_edges}

\textbf{Syntax:} \verb}partition_cell_edges(Gamma, C1, C2)}\\

\textbf{Description:}

Returns the collection of edges in Gamma from a vertex in C1 to one in C2.
    Here, C1, C2 are disjoint sublists (without repetition) of V.

\textbf{Example:}
\begin{verbatim}
EXAMPLES:
        sage: V2 = [0,1,2,3,4,5]
        sage: E2 = [(0,1),(0,5),(1,2),(1,4),(2,3),(3,4),(4,5)]
        sage: Gamma2 = Graph([V2,E2])
        sage: C1 = [0,4]
        sage: C2 = [1,3]
        sage: partition_cell_edges(Gamma2, C1, C2)
         [(0, 1), (1, 4), (3, 4)]
\end{verbatim}

\section{\tt \ \ weights\_of\_vertex\_to\_cell}

\textbf{Syntax:} \verb}weights_of_vertex_to_cell(woGamma, C, v)}\\

\textbf{Description:}
\begin{verbatim}

INPUT:
    - woGamma is a triple: [Gamma = (V, E), epsilon, wt], 
              where 
                 --  epsilon is a pair [E,eps0], 
                 --  eps0 is a list of |E| edge orientations,
                 --  wt is a pair [E,w],   
                 --  w is a list of |E| edge weights > 0.
    - C is a sublist of V,
    - v a vertex in V

OUTPUT:
    the list if weights of the edges from vertex v to cell C.
\end{verbatim}

\textbf{Example:}
\begin{verbatim}
    EXAMPLES:
        sage: V2 = [0,1,2,3,4,5]
        sage: E2 = [(0,1),(0,5),(1,2),(1,4),(2,3),(3,4),(4,5)]
        sage: Gamma2 = Graph([V2,E2])
        sage: C1 = [0,4]
        sage: C2 = [1,3]
        sage: w2 = [1, 4, 2, 2, 3, 2, 1]
        sage: w20 = [E2, w2]
        sage: e20 = [E2, [1]*len(E2)]
        sage: woGamma2 = [Gamma2, e20, w20]
        sage: weights_of_vertex_to_cell(woGamma2, C2, 0)
         [1]
        sage: weights_of_vertex_to_cell(woGamma2, C2, 4)
         [2, 2]
\end{verbatim}

\section{\tt \ \ is\_weighted\_equitable\_partition}

\textbf{Syntax:} \verb}is_weighted_equitable_partition(woGamma, P, verbose=False)}\\

\textbf{Description:}
\begin{verbatim}
INPUT:
    - woGamma is a triple: [Gamma = (V, E), epsilon, wt], 
              where 
                 --  epsilon is a pair [E,eps0], 
                 --  eps0 is a list of |E| edge orientations,
                 --  wt is a pair [E,w],   
                 --  w is a list of |E| edge weights > 0.
    - P     is represented as a list of disjoint lists partitioning V.

OUTPUT:  
       True if P defines an equitable partition of the vertices
             of the given weighted oriented graph woGamma.
       False, otherwise.

If, in addition, verbose=True then it returns the intersection matrix, if
P is equitable, and otherwise prints info why not.

    Gamma must be in standard form.

\end{verbatim}

\textbf{Example:}
\begin{verbatim}
    EXAMPLES:
        sage: V2 = [0,1,2,3,4,5]
        sage: E2 = [(0,1),(0,5),(1,2),(1,4),(2,3),(3,4),(4,5)]
        sage: Gamma2 = Graph([V2,E2])
        sage: C1 = [0,4]
        sage: C2 = [1,3]
        sage: w2 = [1, 4, 2, 2, 3, 2, 1]
        sage: w20 = [E2, w2]
        sage: e20 = [E2, [1]*len(E2)]
        sage: woGamma2 = [Gamma2, e20, w20]
	sage: P = [[0,1],[2,3],[4,5]]
        sage: is_weighted_equitable_partition(woGamma2, P, verbose=False)
         False
        sage: is_weighted_equitable_partition(woGamma2, P, verbose=True)
         At vertex  2  the weights to  [0, 1]  are  [2]
         At vertex  3  the weights to  [0, 1]  are  []
         False
\end{verbatim}

\section{\tt \ \ weighted\_cayley\_graph}

\textbf{Syntax:} \verb}weighted_cayley_graph(G, symm_gp_gens, verbose=False)}\\

\textbf{Description:}
\begin{verbatim}
INPUT:
     - G  is a permutation gp (a subgp of Sn, some n),
     - symm_gp_gens,  a symmetric set of generators, 
OUTPUT:
     1) a simple weighted Cayley graph in standard form,
        whose edges e = (u,v) are weighted by the order 
                   wt(e) = ord_gp(g)
	of the generator g such that v = g*u, 
     2) a list of pairs of the vertices of the current
	graph with the corresponding vertex of the original
	graph G.cayley_graph(generators = symm_gp_gens, simple = True)
\end{verbatim}

\textbf{Example:}
\begin{verbatim}
EXAMPLES:
        sage: gp = S3xC3; symm_gp_gens = [(1,2),(2,3),(4,5,6),(4,6,5)]
        sage: weighted_cayley_graph(gp, symm_gp_gens)
         Graph on 18 vertices
\end{verbatim}

\section{\tt \ \ convolve\_function\_with\_filter}

\textbf{Syntax:} \verb}convolve_function_with_filter(g, h, S)}\\

\textbf{Description:}
\begin{verbatim}

Notation:
     *  g is a function on the vertices V of a graph Gamma = (V, E) in standard form,
        represented as a list of length n = |V|,
     * h is a fixed filter function, represented as a polynomial, h(x).
     * S is a shift operator on Gamma represented by a symmetric nxn matrix, 
       S = (s\_{ij}), where s\_{ij} = 0 if (i,j) is not an edge in Gamma.
 
    We assume everything has coefficients in QQ.

    Returns the graph convolution

          h *\_S g = h(S)g,

    another vector of length n.
\end{verbatim}

\textbf{Example:}
\begin{verbatim}
EXAMPLES:
        sage: Gamma = graphs.DiamondGraph()
        sage: V = Gamma.vertices(); n = len(V); n
         4
        sage: g = [2,1,0,-1]
        sage: h(x) = 1+2*x+3*x^2+5*x^7
        sage: S = Gamma.adjacency_matrix()
        sage: convolve_function_with_filter(g, h, S)
         (13, 15, 13, 10)
        sage: g = [1,0,0,0]; g; convolve_function_with_filter(g, h, S)
        [1, 0, 0, 0]
        (11, 8, 8, 10)
        sage: g = [0,1,0,0]; g; convolve_function_with_filter(g, h, S)
        [0, 1, 0, 0]
        (8, 16, 13, 8)
        sage: g = [0,0,1,0]; g; convolve_function_with_filter(g, h, S)
        [0, 0, 1, 0]
        (8, 13, 16, 8)
        sage: g = [0,0,0,1]; g; convolve_function_with_filter(g, h, S)
        [0, 0, 0, 1]
        (10, 8, 8, 11)
\end{verbatim}
\end{document}
